% El patriotismo y la diversidad cultural — Espagnol 1ère A — Cours signé OnBuch+
% Programme officiel MINESEC. Compiler avec : tectonic EA09-patriotismo-diversidad-cultural.tex
\newcommand{\DOCMATIERE}{Espagnol}
\newcommand{\DOCNIVEAU}{1ère A}
\newcommand{\DOCMODULE}{Módulo 3 : Citoyenneté}
\newcommand{\DOCLECON}{EA09}
\newcommand{\DOCTITRE}{El patriotismo y la diversidad cultural}
% OnBuch+ — préambule « cours signés » ENRICHI (compilé avec tectonic / XeLaTeX)
% Système visuel « Pop » d'OnBuch+ : fond crème (jamais blanc), bordures encre
% épaisses, ombres « dures » décalées, titres Archivo Black en capitales.
% Version MATHÉMATIQUES : graphes (pgfplots/tikz), boîtes pédagogiques
% (définition, propriété, méthode, attention, le savais-tu, exercice résolu,
% exercice, TP, bilan), page de garde et signature OnBuch+ ancrée en bas.
%
% Macros attendues dans main.tex AVANT \input{preamble} :
%   \DOCMATIERE  \DOCNIVEAU  \DOCMODULE  \DOCLECON (numéro)  \DOCTITRE
\documentclass[11pt]{article}

\usepackage{fontspec}
\usepackage[spanish,french]{babel} % français = langue principale ; l'espagnol via \esp{..} / espagnol
\usepackage[a4paper,margin=2cm,top=3.4cm,bottom=2.8cm,headheight=1.4cm,footskip=1.3cm]{geometry}
\usepackage{amsmath,amssymb,amsfonts}
\usepackage{mathtools} % \xmapsto, \coloneqq... souvent utilisés par les modèles
\usepackage{cancel} % simplifications barrées (bilans de forces)
% \abs{z} (module, valeur absolue) et \norm{u} (norme), utilisés par les modèles
\providecommand{\abs}{}\let\abs\relax\DeclarePairedDelimiter{\abs}{\lvert}{\rvert}
\providecommand{\norm}{}\let\norm\relax\DeclarePairedDelimiter{\norm}{\lVert}{\rVert}
\usepackage[table]{xcolor} % [table] : fournit aussi \rowcolors (alternance de lignes)
\usepackage{pagecolor}
\usepackage{tikz}
\usetikzlibrary{arrows.meta,positioning,calc,shapes.geometric,shapes.misc,shapes.arrows,decorations.pathmorphing,decorations.pathreplacing,decorations.markings,patterns,shadows,mindmap,trees,fit,backgrounds,matrix,intersections,angles,quotes}
\usepackage{pgfplots}
\usepackage[version=4]{mhchem} % équations nucléaires \ce{^{235}_{92}U}
\usepackage[european,siunitx]{circuitikz} % schémas électriques
\pgfplotsset{compat=1.18}
\usepgfplotslibrary{fillbetween,groupplots} % aires entre courbes (calcul intégral)
\usepackage{enumitem}
\usepackage{fancyhdr}
% [most] charge toutes les bibliothèques tcolorbox (listings, minted, raster,
% documentation...) et gonfle inutilement la mémoire TeX sur un document long
% (~300 boîtes) : on ne charge que ce qui est réellement utilisé.
\usepackage[skins,breakable]{tcolorbox}
\usepackage{graphicx}
\usepackage{colortbl}
\usepackage{booktabs}
\usepackage{tabularx}
\usepackage{diagbox} % case d'en-tête diagonale des tableaux à double entrée
% Colonnes extensibles centrée / gauche / droite (C, L, R), souvent utilisées par les modèles
\newcolumntype{C}{>{\centering\arraybackslash}X}
\newcolumntype{L}{>{\raggedright\arraybackslash}X}
\newcolumntype{R}{>{\raggedleft\arraybackslash}X}
\usepackage{array}
\usepackage{multicol}
\usepackage{siunitx}
% parse-numbers=false : accepte aussi \SI{2,5 \times 10^{-3}}{...}
\sisetup{locale=FR,output-decimal-marker={,},group-minimum-digits=4,parse-numbers=false}
\DeclareSIUnit\ohm{\ensuremath{\symup{\Omega}}} % U+2126 absent de la police
\DeclareSIUnit\division{div} % réglages d'oscilloscope (ms/div, V/div)
\DeclareSIUnit\tr{tr} % tours (vitesse de rotation, ex. \SI{1500}{\tr\per\minute})
\DeclareSIUnit\molar{M} % concentration molaire (ex. \SI{3}{\molar}, \SI{1}{\micro\molar})
\DeclareSIUnit\an{an} % année (le modèle confond parfois avec \year, un compteur TeX) — ex. \SI{5}{\kilo\meter\per\an}
\DeclareSIUnit\FCFA{FCFA} % franc CFA (ex. \SI{500}{\FCFA\per\kilo\gram})
\DeclareSIUnit\degreeBrix{\textdegree Brix} % teneur en sucre d'un sirop/jus (réfractométrie)
\DeclareSIUnit\month{mois} % le modèle utilise parfois \month, un compteur TeX, comme unité de durée
\usepackage{qrcode}
% \degree : commande fréquemment utilisée par les modèles pour un angle ou une
% température en dehors de \SI{}{} (non fournie par les paquets chargés).
\providecommand{\degree}{^{\circ}}
% \qed (fin de démonstration), utilisé par les modèles sans amsthm
\providecommand{\qed}{\hfill$\square$}
% \barre : double barre des valeurs interdites dans un tableau de variations (tkz-tab)
\providecommand{\barre}{\ensuremath{\|}}
% environnement proof (amsthm non chargé) utilisé par les modèles
\ifdefined\proof\else\newenvironment{proof}[1][Démonstration]{\par\noindent\textit{#1.}\ }{\qed\par}\fi
\usepackage{lastpage}
\usepackage{hyperref}
\hypersetup{hidelinks}

% ============================================================
% Palette « Pop » OnBuch+ — mêmes hex que la classe P (plus_ui.dart)
% ============================================================
\definecolor{poporange}{HTML}{F59321}
\definecolor{popdark}{HTML}{E07A0C}
\definecolor{popcream}{HTML}{FFF6E8}
\definecolor{popink}{HTML}{1C1714}
\definecolor{poppurple}{HTML}{7A5AE0}
\definecolor{popgreen}{HTML}{2E9E6A}
\definecolor{poppink}{HTML}{E0457B}
\definecolor{popblue}{HTML}{2F7DE1}
\definecolor{popgold}{HTML}{C98A12}
\definecolor{popmuted}{HTML}{8A7F76}
\definecolor{popline}{HTML}{EADFD2}
\definecolor{popgray}{HTML}{8A7F76}
\colorlet{popgrey}{popgray}
% Alias tolérants (le modèle nomme parfois les couleurs différemment)
\colorlet{popwhite}{white}
\colorlet{popblack}{popink}
\colorlet{popred}{poppink}
\colorlet{popyellow}{popgold}
% Teintes claires (fonds de tableaux/encadrés) — le modèle nomme parfois
% les couleurs avec un suffixe L (ex. popblueL) sans les avoir définies.
\colorlet{poporangeL}{poporange!20}
\colorlet{popdarkL}{popdark!20}
\colorlet{poppurpleL}{poppurple!20}
\colorlet{popgreenL}{popgreen!20}
\colorlet{poppinkL}{poppink!20}
\colorlet{popblueL}{popblue!20}
\colorlet{popgoldL}{popgold!20}
\colorlet{popredL}{popred!20}
\colorlet{popgrayL}{popgray!20}
% Teintes claires pour fonds de tableaux / zones
\colorlet{popblueL}{popblue!10!white}
\colorlet{poporangeL}{poporange!14!white}
\colorlet{popgreenL}{popgreen!12!white}
\colorlet{poppurpleL}{poppurple!12!white}
\colorlet{poppinkL}{poppink!12!white}
\colorlet{popgoldL}{popgold!14!white}

\pagecolor{popcream}

% ============================================================
% Typographie
% ============================================================
\defaultfontfeatures{Ligatures=TeX}
\setmainfont{PlusJakartaSans-Medium.ttf}[Path=../../../fonts/,BoldFont=PlusJakartaSans-Bold.ttf]
% Maths en sans-serif (Fira Math) avec chiffres et lettres droites de Plus
% Jakarta, pour que les formules se fondent dans le texte.
\usepackage[math-style=ISO,bold-style=ISO]{unicode-math}
\setmathfont{FiraMath-Regular.otf}[Path=../../../fonts/]
\setmathfont{PlusJakartaSans-Medium.ttf}[Path=../../../fonts/,range={up/num,up/{Latin,latin},\mathup}]
\usepackage{icomma} % virgule décimale sans espace en mode math
% Caractères mathématiques Unicode absents des polices de texte (Plus Jakarta,
% Archivo Black) : rendus par leur équivalent mathématique, en texte comme en
% formule — sinon ils s'affichent en carré vide (ex. « Arithmétique dans ℤ »).
\usepackage{newunicodechar}
\newunicodechar{ℝ}{\ensuremath{\mathbb{R}}}
\newunicodechar{ℤ}{\ensuremath{\mathbb{Z}}}
\newunicodechar{ℂ}{\ensuremath{\mathbb{C}}}
\newunicodechar{ℕ}{\ensuremath{\mathbb{N}}}
\newunicodechar{ℚ}{\ensuremath{\mathbb{Q}}}
\newunicodechar{𝔻}{\ensuremath{\mathbb{D}}}
\newunicodechar{α}{\ensuremath{\alpha}}
\newunicodechar{β}{\ensuremath{\beta}}
\newunicodechar{γ}{\ensuremath{\gamma}}
\newunicodechar{δ}{\ensuremath{\delta}}
\newunicodechar{ε}{\ensuremath{\varepsilon}}
\newunicodechar{θ}{\ensuremath{\theta}}
\newunicodechar{λ}{\ensuremath{\lambda}}
\newunicodechar{μ}{\ensuremath{\mu}}
\newunicodechar{σ}{\ensuremath{\sigma}}
\newunicodechar{τ}{\ensuremath{\tau}}
\newunicodechar{φ}{\ensuremath{\varphi}}
\newunicodechar{ψ}{\ensuremath{\psi}}
\newunicodechar{ω}{\ensuremath{\omega}}
\newunicodechar{Σ}{\ensuremath{\Sigma}}
% majuscules produites par \MakeUppercase dans les titres (α-aminés -> Α)
\newunicodechar{Α}{\ensuremath{\alpha}}
\newunicodechar{Β}{\ensuremath{\beta}}
\newunicodechar{Γ}{\ensuremath{\Gamma}}
\newunicodechar{Δ}{\ensuremath{\Delta}}
\newunicodechar{Ε}{\ensuremath{\varepsilon}}
\newunicodechar{Θ}{\ensuremath{\Theta}}
\newunicodechar{Λ}{\ensuremath{\Lambda}}
\newunicodechar{Μ}{\ensuremath{\mu}}
\newunicodechar{Τ}{\ensuremath{\tau}}
\newunicodechar{Φ}{\ensuremath{\Phi}}
\newunicodechar{Ψ}{\ensuremath{\Psi}}
\newunicodechar{Ω}{\ensuremath{\Omega}}
\newunicodechar{↦}{\ensuremath{\mapsto}}
\newunicodechar{∈}{\ensuremath{\in}}
\newunicodechar{∉}{\ensuremath{\notin}}
\newunicodechar{∧}{\ensuremath{\wedge}}
\newunicodechar{⇔}{\ensuremath{\Leftrightarrow}}
\newunicodechar{⟺}{\ensuremath{\Longleftrightarrow}}
\newunicodechar{⇒}{\ensuremath{\Rightarrow}}
\newunicodechar{⟹}{\ensuremath{\Longrightarrow}}
\newunicodechar{∘}{\ensuremath{\circ}}
\newunicodechar{∪}{\ensuremath{\cup}}
\newunicodechar{∩}{\ensuremath{\cap}}
\newunicodechar{≡}{\ensuremath{\equiv}}
\newunicodechar{⊂}{\ensuremath{\subset}}
\newunicodechar{⊥}{\ensuremath{\perp}}
\newunicodechar{∃}{\ensuremath{\exists}}
\newunicodechar{∀}{\ensuremath{\forall}}
\newunicodechar{∛}{\ensuremath{\sqrt[3]{\,}}}
\newunicodechar{∜}{\ensuremath{\sqrt[4]{\,}}}
\newunicodechar{⃗}{\ensuremath{{}^{\to}}}
\newunicodechar{ⁿ}{\ifmmode^{n}\else\textsuperscript{\ensuremath{n}}\fi}
\newunicodechar{ˣ}{\ifmmode^{x}\else\textsuperscript{\ensuremath{x}}\fi}
\newunicodechar{ᵇ}{\ifmmode^{b}\else\textsuperscript{\ensuremath{b}}\fi}
\newunicodechar{ʳ}{\ifmmode^{r}\else\textsuperscript{\ensuremath{r}}\fi}
\newunicodechar{ᵘ}{\ifmmode^{u}\else\textsuperscript{\ensuremath{u}}\fi}
\newunicodechar{ʸ}{\ifmmode^{y}\else\textsuperscript{\ensuremath{y}}\fi}
\newunicodechar{ᵃ}{\ifmmode^{a}\else\textsuperscript{\ensuremath{a}}\fi}
\newunicodechar{ᵉ}{\ifmmode^{e}\else\textsuperscript{\ensuremath{e}}\fi}
\newunicodechar{ᵏ}{\ifmmode^{k}\else\textsuperscript{\ensuremath{k}}\fi}
\newunicodechar{ᵖ}{\ifmmode^{p}\else\textsuperscript{\ensuremath{p}}\fi}
\newunicodechar{ᴺ}{\ifmmode^{N}\else\textsuperscript{\ensuremath{N}}\fi}
\newunicodechar{ᶜ}{\ifmmode^{c}\else\textsuperscript{\ensuremath{c}}\fi}
\newunicodechar{ʰ}{\ifmmode^{h}\else\textsuperscript{\ensuremath{h}}\fi}
\newunicodechar{ᵠ}{\ifmmode^{q}\else\textsuperscript{\ensuremath{q}}\fi}
\newunicodechar{ₙ}{\ifmmode_{n}\else\textsubscript{\ensuremath{n}}\fi}
\newunicodechar{ₐ}{\ifmmode_{a}\else\textsubscript{\ensuremath{a}}\fi}
\newunicodechar{ᵢ}{\ifmmode_{i}\else\textsubscript{\ensuremath{i}}\fi}
\newunicodechar{ᵣ}{\ifmmode_{r}\else\textsubscript{\ensuremath{r}}\fi}
\newunicodechar{ₖ}{\ifmmode_{k}\else\textsubscript{\ensuremath{k}}\fi}
\newunicodechar{ᵦ}{\ifmmode_{\beta}\else\textsubscript{\ensuremath{\beta}}\fi}
\newunicodechar{ₓ}{\ifmmode_{x}\else\textsubscript{\ensuremath{x}}\fi}
\newfontfamily\popdisplay{ArchivoBlack-Regular.ttf}[Path=../../../fonts/]
\newfontfamily\poplogo{SpaceGrotesk-Bold.ttf}[Path=../../../fonts/]
\newfontfamily\popsemi{PlusJakartaSans-SemiBold.ttf}[Path=../../../fonts/]
\newfontfamily\popxbold{PlusJakartaSans-ExtraBold.ttf}[Path=../../../fonts/]
\newfontfamily\popxboldls{PlusJakartaSans-ExtraBold.ttf}[Path=../../../fonts/,LetterSpace=3.0]

\setlist[enumerate]{leftmargin=1.7em,itemsep=5pt,topsep=4pt}
\setlist[itemize]{leftmargin=1.7em,itemsep=4pt,topsep=4pt,label={\color{poporange}\textbullet}}
\setlength{\parindent}{0pt}
\setlength{\parskip}{5pt}
\renewcommand{\arraystretch}{1.25}

% pgfplots : style Pop par défaut
\pgfplotsset{
  every axis/.append style={axis line style={popink,line width=1pt},tick style={popink},
    grid style={popline,line width=0.5pt},label style={font=\small},tick label style={font=\footnotesize}},
  popcurve/.style={poporange,line width=1.8pt,no markers,smooth},
}
% Même style disponible dans un \draw TikZ simple (hors pgfplots)
\tikzset{popcurve/.style={poporange,line width=1.8pt}}
\tikzset{popgrid/.style={popline,very thin}}
\colorlet{popcurve}{poporange} % le nom du style est parfois utilisé comme couleur

% ============================================================
% Pill / squiggle
% ============================================================
\newcommand{\poppill}[3]{%
  \tikz[baseline=-0.55ex]\node[draw=popink,line width=1.2pt,rounded corners=2.6pt,%
    fill=#2,inner xsep=6.5pt,inner ysep=3pt]{%
    \popxboldls\fontsize{7}{7}\selectfont\color{#3}\MakeUppercase{#1}};%
}
\newcommand{\popsquiggle}[1]{%
  \tikz[baseline=-0.4ex]{
    \draw[#1,line width=2.6pt,line cap=round]
      plot [smooth] coordinates {(0,0.09) (0.34,0.01) (0.68,0.21) (1.02,0.01) (1.36,0.10)};
  }%
}
% Mot-clé surligné (marqueur orange sous le texte)
\newcommand{\cle}[1]{{\popxbold\color{popdark}#1}}

% ============================================================
% En-tête / pied
% ============================================================
\pagestyle{fancy}
\fancyhf{}
\renewcommand{\headrulewidth}{0pt}
\renewcommand{\footrulewidth}{0pt}
\lhead{\raisebox{-0.15ex}{\poplogo\fontsize{13}{13}\selectfont\color{popink}OnBuch\color{poporange}+}}
\rhead{\poppill{\DOCMATIERE\ \textbullet\ \DOCNIVEAU\ \textbullet\ Leçon \DOCLECON}{white}{popink}}
\lfoot{\popxbold\fontsize{7.6}{7.6}\selectfont\color{popmuted}\MakeUppercase{Cours signé OnBuch+}\ \textbullet\ {\popsemi\DOCTITRE}}
\rfoot{\tikz[baseline=-0.6ex]\node[rounded corners=8.5pt,fill=popink,minimum height=17pt,minimum width=17pt,inner xsep=5pt,inner sep=0]{\popxbold\fontsize{8}{8}\selectfont\color{popcream}\thepage\,/\,\pageref*{LastPage}};}

% ============================================================
% Titres
% ============================================================
\newcommand{\chaptitre}[2]{%
  \begin{tcolorbox}[enhanced,colback=poporange,colframe=popink,boxrule=2.6pt,arc=11pt,
      drop shadow=popink,left=15pt,right=15pt,top=12pt,bottom=11pt,before skip=3pt,after skip=18pt]
    \poppill{#2}{popink}{popcream}\par\vspace{9pt}
    {\raggedright\hyphenpenalty=10000\exhyphenpenalty=10000\popdisplay\fontsize{22}{24}\selectfont\color{popcream}\MakeUppercase{#1}\par}
  \end{tcolorbox}
}
% Section numérotée : pastille orange avec le numéro + titre Archivo
\newcounter{popsec}
\newcommand{\coursec}[1]{%
  \stepcounter{popsec}\setcounter{popsub}{0}%
  \par\vspace{16pt}\noindent
  \begin{minipage}{\linewidth}
  \tikz[baseline=-0.7ex]\node[draw=popink,line width=1.6pt,fill=poporange,rounded corners=4pt,minimum size=22pt,inner sep=0]{\popdisplay\fontsize{12}{12}\selectfont\color{popcream}\thepopsec};%
  \hspace{8pt}{\popdisplay\fontsize{15}{17}\selectfont\color{popink}\MakeUppercase{#1}}\par
  \vspace{4pt}\noindent{\color{poporange}\rule{36pt}{3.6pt}}
  \end{minipage}\par\nopagebreak\vspace{8pt}
}
\newcounter{popsub}
\newcommand{\courssub}[1]{%
  \stepcounter{popsub}%
  \par\vspace{9pt}\noindent{\popxbold\fontsize{11.5}{13}\selectfont\color{popink}{\color{poporange}\thepopsec.\thepopsub}\hspace{6pt}#1}\par\nopagebreak\vspace{4pt}
}

% ============================================================
% Boîtes pédagogiques — même recette : fond blanc, bordure épaisse colorée,
% ombre dure encre, badge plein. Titre optionnel [#1].
% ============================================================
\tcbset{popbase/.style={breakable,enhanced,colback=white,boxrule=2.3pt,arc=9pt,
  shadow={3pt}{-3pt}{0pt}{popink},left=14pt,right=14pt,top=11pt,bottom=11pt,before skip=11pt,after skip=15pt,
  fontupper=\popsemi\fontsize{10.6}{14.5}\selectfont\color{popink},
  fontlower=\popsemi\fontsize{10.6}{14.5}\selectfont\color{popink}}}
% Coche dessinée (le glyphe ✓ n'existe pas dans les polices du document)
\newcommand{\popcheck}[1]{\tikz[baseline=-0.5ex]\draw[#1,line width=1.6pt,line cap=round,line join=round](-0.13,0.0)--(-0.04,-0.09)--(0.13,0.1);}
\providecommand{\coche}{\popcheck{popgreen}}
\providecommand{\resultat}[1]{\par\smallskip\noindent\fcolorbox{popgreen}{popgreenL}{\parbox{\dimexpr\linewidth-2\fboxsep-2\fboxrule\relax}{\textbf{Résultat.} #1}}\par\smallskip}
\newcommand{\popboxhead}[3]{\poppill{#1}{#2}{white}\ifx\relax#3\relax\else\hspace{6pt}{\popxbold\fontsize{10.6}{12}\selectfont\color{popink}#3}\fi\par\vspace{7pt}}

\newtcolorbox{aretenir}[1][]{popbase,colframe=popgreen,before upper={\popboxhead{À retenir}{popgreen}{#1}}}
% Texte en espagnol (document de lecture, dialogue, poème, extrait) : cadre
% d'étude. Un titre optionnel (source, auteur) s'affiche dans l'en-tête.
\newtcolorbox{textebox}[1][]{popbase,colframe=popblue,before upper={\popboxhead{Texto}{popblue}{#1}\begin{otherlanguage}{spanish}},after upper={\end{otherlanguage}}}
% Marques ✓ / ✗ (forme correcte / fautive) souvent utilisées par les modèles
\providecommand{\cross}{\textcolor{poppink}{\ensuremath{\times}}}
\providecommand{\xmark}{\cross}\providecommand{\crossmark}{\cross}\providecommand{\cmark}{\ensuremath{\checkmark}}
% Ligne de réponse / justification à compléter (inventée par les modèles)
\providecommand{\justification}[1]{\par\noindent\textit{Justification~:} \dotfill\par}
% \textipa (package tipa non chargé) : la police ne couvre pas l'API -> simple passage
\providecommand{\textipa}[1]{#1}
% Soulignement (ulem non chargé) : \uline, \ul
\providecommand{\uline}[1]{\underline{#1}}\providecommand{\ul}[1]{\underline{#1}}
% Ordinaux français écrits en macro par les modèles : 1\ière, 3\ième
\newcommand{\ière}{\textsuperscript{re}}\newcommand{\ième}{\textsuperscript{e}}
% Mot ou phrase en espagnol à l'intérieur d'un paragraphe français
\newcommand{\esp}[1]{\ifmmode\text{\textit{#1}}\else\begingroup\selectlanguage{spanish}\itshape #1\endgroup\fi}
\newtcolorbox{exemplebox}[1][]{popbase,colframe=poporange,before upper={\popboxhead{Exemple}{poporange}{#1}}}
\newtcolorbox{definition}[1][]{popbase,colframe=popblue,before upper={\popboxhead{Définition}{popblue}{#1}}}
\newtcolorbox{propriete}[1][]{popbase,colframe=poppurple,before upper={\popboxhead{Propriété}{poppurple}{#1}}}
\newtcolorbox{methode}[1][]{popbase,colframe=popgold,before upper={\popboxhead{Méthode}{popgold}{#1}}}
\newtcolorbox{attention}[1][]{popbase,colframe=poppink,before upper={\popboxhead{Attention}{poppink}{#1}}}
\newtcolorbox{experience}[1][]{popbase,colframe=popblue,colback=popblueL,before upper={\popboxhead{Expérience}{popblue}{#1}}}
% « Le savais-tu ? » : carte violette pleine (contexte, culture, Cameroun)
\newtcolorbox{savaistu}[1][]{popbase,colback=poppurple,colframe=popink,
  fontupper=\popsemi\fontsize{10.4}{14.2}\selectfont\color{white},
  before upper={\poppill{Le savais-tu\,?}{white}{poppurple}\ifx\relax#1\relax\else\hspace{6pt}{\popxbold\color{white}#1}\fi\par\vspace{7pt}}}
% Exercice résolu : énoncé en haut, \tcblower puis la solution sur fond crème
\newcounter{popexo}\newcounter{popexr}
\newtcolorbox{exoresolu}[1][]{popbase,colframe=poporange,colbacklower=popcream,
  segmentation style={popink,line width=1.2pt,dashed},
  before upper={\stepcounter{popexr}\popboxhead{Exercice résolu \thepopexr}{poporange}{#1}},
  before lower={\poppill{Solution}{popink}{popcream}\par\vspace{6pt}}}
% Exercice (non résolu en place) — la correction est dans la partie Corrigés
\newtcolorbox{exercice}[1][]{popbase,colframe=popink,
  before upper={\stepcounter{popexo}\popboxhead{Exercice \thepopexo}{popink}{#1}}}
% Corrigé
\newtcolorbox{corrige}[1][]{popbase,colframe=popgreen,colback=popgreenL,
  before upper={\popboxhead{Corrigé}{popgreen}{#1}}}
% Objectifs / prérequis / bilan
\newtcolorbox{objectifs}{popbase,colframe=popink,colback=poporangeL,
  before upper={\popboxhead{Objectifs de la leçon}{popink}{}}}
\newtcolorbox{prerequis}{popbase,colframe=popink,colback=popblueL,
  before upper={\popboxhead{Prérequis}{popblue}{}}}
\newtcolorbox{bilan}{popbase,colframe=popink,colback=white,boxrule=2.8pt,
  before upper={\popboxhead{Fiche bilan}{poporange}{L'essentiel à connaître pour l'examen}}}
% Figure encadrée avec légende
\newtcolorbox{popfigure}[1][]{enhanced,breakable=false,colback=white,colframe=popline,boxrule=1.4pt,arc=8pt,
  left=10pt,right=10pt,top=10pt,bottom=8pt,before skip=10pt,after skip=12pt,halign=center,
  lower separated=false,halign lower=center,
  fontlower=\popsemi\fontsize{9}{11.5}\selectfont\color{popmuted},
  #1}
\newcommand{\legende}[1]{\par\vspace{6pt}{\popsemi\fontsize{9}{11.5}\selectfont\color{popmuted}#1}}

% ============================================================
% Signature OnBuch+
% ============================================================
\newcommand{\poplogoinline}[1]{{\poplogo\fontsize{#1}{#1}\selectfont\color{popink}OnBuch\color{poporange}+}}
\newcommand{\signatureonbuch}{%
  \begin{tcolorbox}[enhanced,colback=popink,colframe=popink,boxrule=2.2pt,arc=10pt,
      drop shadow=poporange,left=14pt,right=14pt,top=10pt,bottom=10pt,before skip=0pt,after skip=0pt]
    \tikz[baseline=-0.7ex]\node[circle,fill=popgreen,draw=popcream,line width=1.3pt,minimum size=22pt,inner sep=0]{\popcheck{white}};%
    \hspace{9pt}\begin{minipage}[c]{0.62\linewidth}
      {\popxbold\fontsize{12}{13}\selectfont\color{popcream}Signé\ \textbullet\ {\poplogo OnBuch\color{poporange}+}}\par\vspace{2pt}
      {\raggedright\popsemi\fontsize{8}{10.5}\selectfont\color{popline}\MakeUppercase{Équipe pédagogique OnBuch+}\\\MakeUppercase{Conforme au programme officiel MINESEC}\par}
    \end{minipage}\hfill
    {\poplogo\fontsize{22}{22}\selectfont\color{popcream}OnBuch\color{poporange}+}
  \end{tcolorbox}%
}

% ============================================================
% Page de garde
% #1 titre · #2 matière · #3 niveau · #4 module · #5 n° leçon · #6 durée ·
% #7 description / objectifs (texte ou itemize)
% ============================================================
\newcommand{\pagegarde}[7]{%
  \thispagestyle{empty}
  \newgeometry{a4paper,margin=1.7cm,top=1.6cm,bottom=1.4cm}%
  \begin{tikzpicture}[remember picture,overlay]
    \fill[poporange!18] ([xshift=-3.2cm,yshift=-3.6cm]current page.north east) circle (4.6cm);
    \fill[poppurple!14] ([xshift=2.4cm,yshift=7.6cm]current page.south west) circle (3.8cm);
  \end{tikzpicture}%
  {\poplogo\fontsize{30}{30}\selectfont\color{popink}OnBuch\color{poporange}+}\hfill
  \raisebox{0.6ex}{\poppill{Cours signé \textbullet\ Premium}{popink}{popcream}}\par
  \vspace{1pt}\popsquiggle{poppurple}\par
  \vspace{8pt}
  \begin{tcolorbox}[enhanced,colback=poporange,colframe=popink,boxrule=2.8pt,arc=13pt,
      drop shadow=popink,left=18pt,right=18pt,top=12pt,bottom=13pt,after skip=10pt]
    \poppill{#2 \textbullet\ #3}{popink}{popcream}\hspace{5pt}\poppill{#4}{white}{popink}\par\vspace{8pt}
    {\popdisplay\fontsize{14}{14}\selectfont\color{popink}LEÇON #5}\par\vspace{4pt}
    {\raggedright\hyphenpenalty=10000\exhyphenpenalty=10000\popdisplay\fontsize{25}{27}\selectfont\color{popcream}\MakeUppercase{#1}\par}
    \vspace{8pt}
    {\popxbold\fontsize{9.5}{9.5}\selectfont\color{popink}\MakeUppercase{Durée conseillée : #6}}
  \end{tcolorbox}
  \begin{tcolorbox}[enhanced,colback=white,colframe=popink,boxrule=2.2pt,arc=10pt,
      drop shadow=popink,left=16pt,right=16pt,top=10pt,bottom=10pt,after skip=8pt]
    \poppill{Dans cette leçon}{poppurple}{white}\par\vspace{5pt}
    \popsemi\fontsize{9.3}{12.6}\selectfont\color{popink}#7
  \end{tcolorbox}
  \noindent\begin{minipage}[t]{0.487\linewidth}
    \begin{tcolorbox}[enhanced,colback=popgreen,colframe=popink,boxrule=2.2pt,arc=10pt,
        drop shadow=popink,left=11pt,right=11pt,top=10pt,bottom=10pt,height=3.1cm,valign=center]
      \tcbox[colback=white,colframe=popink,boxrule=1.3pt,arc=5pt,boxsep=0pt,nobeforeafter,
        left=3pt,right=3pt,top=3pt,bottom=3pt,valign=center]{\qrcode[height=1.9cm]{https://play.google.com/store/apps/details?id=com.luvvix.onbuch&hl=fr}}%
      \hspace{8pt}\begin{minipage}[c]{0.5\linewidth}
        {\popdisplay\fontsize{11}{12.5}\selectfont\color{white}\MakeUppercase{Télécharge OnBuch}}\par\vspace{4pt}
        {\raggedright\popsemi\fontsize{8.4}{11}\selectfont\color{white}Gratuit \textbullet\ Google Play\\Tuteur IA, annales, concours, résultats\par}
      \end{minipage}
    \end{tcolorbox}
  \end{minipage}\hfill
  \begin{minipage}[t]{0.487\linewidth}
    \begin{tcolorbox}[enhanced,colback=poppurple,colframe=popink,boxrule=2.2pt,arc=10pt,
        drop shadow=popink,left=11pt,right=11pt,top=10pt,bottom=10pt,height=3.1cm,valign=center]
      \tikz[baseline=-0.7ex]\node[circle,fill=white,draw=popink,line width=1.3pt,minimum size=1.6cm,inner sep=0]{\popdisplay\fontsize{24}{24}\selectfont\color{poppurple}+};%
      \hspace{8pt}\begin{minipage}[c]{0.6\linewidth}
        {\popdisplay\fontsize{11}{12.5}\selectfont\color{white}\MakeUppercase{Passe à OnBuch+}}\par\vspace{4pt}
        {\raggedright\popsemi\fontsize{8.4}{11}\selectfont\color{white}Tous les cours signés, Tuteur IA illimité, fiches PDF — onglet OnBuch+ de l'app\par}
      \end{minipage}
    \end{tcolorbox}
  \end{minipage}
  \vspace{8pt}
  \popcontenu
  \vfill
  \signatureonbuch
  \restoregeometry
  \newpage
}

% Rangée « Ce cours contient » (page de garde)
\newcommand{\poptile}[3]{% #1 couleur #2 gros texte #3 légende
  \begin{tcolorbox}[enhanced,colback=#1,colframe=popink,boxrule=2pt,arc=9pt,shadow={2.5pt}{-2.5pt}{0pt}{popink},
      left=6pt,right=6pt,top=9pt,bottom=9pt,halign=center,valign=center,height=1.75cm,width=\linewidth,nobeforeafter]
    {\popdisplay\fontsize{12.5}{13}\selectfont\color{white}\MakeUppercase{#2}}\par\vspace{3pt}
    {\centering\popsemi\fontsize{7.8}{9.5}\selectfont\color{white}#3\par}
  \end{tcolorbox}}
\newcommand{\popcontenu}{%
  \par\noindent{\popxboldls\fontsize{7.5}{7.5}\selectfont\color{popmuted}CE COURS SIGNÉ CONTIENT}\par\vspace{7pt}
  \noindent\begin{minipage}[t]{0.235\linewidth}\poptile{popblue}{Cours}{complet, illustré, pas à pas}\end{minipage}\hfill
  \begin{minipage}[t]{0.235\linewidth}\poptile{poporange}{Méthodes}{et exercices résolus}\end{minipage}\hfill
  \begin{minipage}[t]{0.235\linewidth}\poptile{poppink}{Exercices}{type examen + corrigés}\end{minipage}\hfill
  \begin{minipage}[t]{0.235\linewidth}\poptile{popgreen}{Fiche bilan}{l'essentiel à retenir}\end{minipage}\par}

% Sommaire
\newtcolorbox{sommaire}{enhanced,colback=white,colframe=popink,boxrule=2.2pt,arc=10pt,
  drop shadow=popink,left=18pt,right=18pt,top=13pt,bottom=13pt,before skip=6pt,after skip=20pt,
  before upper={\poppill{Sommaire}{poppurple}{white}\par\vspace{9pt}\popsemi\fontsize{10.6}{14.5}\selectfont\color{popink}}}

% Signature de fin de cours — poussée tout en bas de la dernière page
\newcommand{\findecours}{%
  \par\vfill\nopagebreak
  \begin{center}\popsquiggle{poporange}\end{center}\vspace{4pt}
  \signatureonbuch
}
\AtBeginDocument{\def\llbracket{\mathopen{[\mkern-3.5mu[}}\def\rrbracket{\mathclose{]\mkern-3.5mu]}}} % intervalles d'entiers (glyphe ⟦ absent de la police)
% Glyphes absents de Fira Math : redéfinis à partir de glyphes disponibles
\AtBeginDocument{%
  \def\setminus{\mathbin{\backslash}}\let\smallsetminus\setminus
  \def\vdots{\mathord{\vbox{\baselineskip3.2pt\lineskiplimit0pt\kern4pt\hbox{.}\hbox{.}\hbox{.}}}}%
  \def\ddots{\mathinner{\mkern1mu\raise6.5pt\hbox{.}\mkern2mu\raise3.6pt\hbox{.}\mkern2mu\raise0.7pt\hbox{.}\mkern1mu}}%
  \def\triangle{\mathord{\scalebox{0.9}{$\symup{\Delta}$}}}%
  \def\nabla{\mathord{\raisebox{\depth}{\scalebox{1}[-1]{$\symup{\Delta}$}}}}%
  \def\checkmark{\popcheck{popdark}}%
  \def\star{\mathbin{\ast}}\def\bigstar{\mathord{\ast}}%
  \def\diamond{\mathbin{\scalebox{0.6}{$\Diamond$}}}%
  \def\bigcup{\mathop{\mathchoice{\raisebox{-0.3em}{\scalebox{1.7}{$\cup$}}}{\scalebox{1.25}{$\cup$}}{\cup}{\cup}}\displaylimits}%
  \def\bigcap{\mathop{\mathchoice{\raisebox{-0.3em}{\scalebox{1.7}{$\cap$}}}{\scalebox{1.25}{$\cap$}}{\cap}{\cap}}\displaylimits}%
  \def\lesseqqgtr{\mathrel{\lessgtr}}\def\bigoplus{\mathop{\oplus}\displaylimits}%
}
\providecommand{\ssi}{si et seulement si} % abréviation fréquente
\AtBeginDocument{\providecommand{\wideparen}{\overparen}} % arc de cercle

%% FIGFIX-BEGIN v5 — correctifs globaux des figures (docs/onbuch-generation/tools/figfix)
\usepackage{adjustbox}\usepackage{etoolbox}\tcbuselibrary{raster}
% toute figure TikZ/pgfplots plus large que la ligne est réduite à la largeur disponible
\BeforeBeginEnvironment{tikzpicture}{\begin{adjustbox}{max width=\linewidth}}
\AfterEndEnvironment{tikzpicture}{\end{adjustbox}}
% légendes pgfplots lisibles par-dessus les courbes
\pgfplotsset{every axis/.append style={legend style={fill=white,fill opacity=0.92,text opacity=1,draw=black!15,inner sep=3pt}}}
% étiquettes d'arêtes (syntaxe edge["..."]) avec fond blanc
\tikzset{every edge quotes/.append style={fill=white,inner sep=1.2pt}}
%% FIGFIX-END
\begin{document}

\pagegarde{El patriotismo y la diversidad cultural}{Espagnol}{1ère A}{Módulo 3 : Citoyenneté}{EA09}{2 h}{Cette leçon propose la lecture et le commentaire d'un texte sur le patriotisme et la diversité culturelle au Cameroun et dans le monde hispanophone. Elle permet d'acquérir le lexique des symboles nationaux et des traditions, et de mobiliser ser/estar et le présent pour présenter une fête ou promouvoir la diversité culturelle.
\vspace{4pt}
\begin{multicols}{2}\small
\begin{itemize}
\item À la fin de cette leçon, je sais lire et comprendre un texte sur le patriotisme et la diversité culturelle
\item À la fin de cette leçon, je sais identifier et utiliser le lexique des symboles nationaux (patria, bandera, himno) et des traditions (fiesta, lengua, patrimonio, mestizaje)
\item À la fin de cette leçon, je sais dégager l'idée principale et les idées secondaires d'un texte argumentatif simple
\item À la fin de cette leçon, je sais rédiger un commentaire structuré (introduction, développement, conclusion) sur un texte culturel
\item À la fin de cette leçon, je sais utiliser ser et estar pour décrire une tradition et situer une fête
\item À la fin de cette leçon, je sais employer le présent de l'indicatif pour parler des coutumes et des habitudes culturelles
\end{itemize}\end{multicols}}

\begin{sommaire}
\begin{multicols}{2}
\begin{enumerate}
\item Texte support : « Mi patria, mis culturas »
\item Compréhension du texte
\item Lexique thématique : patria, símbolos y tradiciones
\item Grammaire outillée : ser/estar et le présent des traditions
\item Méthode du commentaire de texte
\item Commentaire modèle et production guidée
\item Activité d'intégration
\item Méthodes et exercices résolus
\item Exercices
\item Corrigés des exercices
\item Fiche bilan
\end{enumerate}
\end{multicols}
\end{sommaire}

\begin{objectifs}
\begin{itemize}
\item À la fin de cette leçon, je sais lire et comprendre un texte sur le patriotisme et la diversité culturelle
\item À la fin de cette leçon, je sais identifier et utiliser le lexique des symboles nationaux (patria, bandera, himno) et des traditions (fiesta, lengua, patrimonio, mestizaje)
\item À la fin de cette leçon, je sais dégager l'idée principale et les idées secondaires d'un texte argumentatif simple
\item À la fin de cette leçon, je sais rédiger un commentaire structuré (introduction, développement, conclusion) sur un texte culturel
\item À la fin de cette leçon, je sais utiliser ser et estar pour décrire une tradition et situer une fête
\item À la fin de cette leçon, je sais employer le présent de l'indicatif pour parler des coutumes et des habitudes culturelles
\item À la fin de cette leçon, je sais présenter oralement une tradition camerounaise ou hispanophone
\item À la fin de cette leçon, je sais rédiger un court texte pour promouvoir la diversité culturelle en espagnol
\end{itemize}
\end{objectifs}

\begin{prerequis}
\begin{itemize}
\item Le présent de l'indicatif des verbes réguliers et irréguliers courants (Seconde)
\item Première approche de ser et estar (Seconde)
\item Le lexique de la famille, de la ville et des loisirs (Seconde)
\item Les adjectifs possessifs et les nationalités (Seconde)
\item La méthode de compréhension de texte : repérage du thème et des personnages (Seconde)
\end{itemize}
\end{prerequis}

\newpage

\coursec{Texte support : « Mi patria, mis culturas »}

\courssub{Situation d'accroche et problématique}

C'est le 20 mai à Yaoundé. Sur le boulevard du 20 mai, les élèves des lycées défilent au pas cadencé, les fanfares résonnent, et des milliers de voix chantent l'hymne national. Partout flottent les couleurs vert-rouge-jaune du drapeau camerounais, fièrement portées sur les pagnes, les casquettes et les visages peints. Quelques semaines plus tard, Carlos, un correspondant mexicain, raconte dans une lettre comment, le 16 septembre, son pays célèbre son indépendance avec le célèbre \esp{Grito de Dolores}. Et une amie de Malabo décrit les danses traditionnelles et les langues multiples de la Guinée équatoriale.

Trois pays, trois continents de cultures, trois façons d'aimer sa patrie. Comment exprimer en espagnol cet amour de la patrie et cette richesse des cultures ? Qu'est-ce qui unit un Camerounais, un Espagnol et un Mexicain malgré leurs différences ? C'est tout l'enjeu du \cle{patriotisme culturel} et de la \cle{diversité culturelle}, que nous allons découvrir à travers le texte de cette leçon.

\courssub{Comment lire un texte : la démarche}

Avant de découvrir le texte, rappelons la méthode de lecture efficace que tu utiliseras toute l'année, en classe comme à l'examen.

\begin{methode}[Lire et comprendre un texte en espagnol]
\begin{enumerate}
\item \textbf{Lire le titre} et deviner le thème : le titre \esp{Mi patria, mis culturas} annonce un texte personnel (\esp{mi}, \esp{mis}) sur la patrie et les cultures.
\item \textbf{Première lecture globale}, silencieuse, sans dictionnaire : chercher le sens général, pas chaque mot.
\item \textbf{Souligner les mots inconnus} et vérifier leur sens au glossaire ou grâce au contexte.
\item \textbf{Repérer les idées principales} paragraphe par paragraphe : qui parle ? de quoi ? où ? quand ?
\item \textbf{Deuxième lecture attentive} : relever les mots-clés du champ lexical (patrie, fête, langue) et les connecteurs (\esp{pero}, \esp{porque}, \esp{por eso}).
\end{enumerate}
\end{methode}

\begin{experience}[Lecture en classe]
\begin{enumerate}
\item \textbf{Lecture silencieuse} (3 minutes) : chaque élève lit le texte seul et souligne les mots inconnus.
\item \textbf{Lecture à haute voix} par l'enseignant : écoutez la prononciation, l'intonation et les pauses. Attention : \esp{patria} se prononce « pa-tria », \esp{bandera} « ban-dé-ra », \esp{Guinea} « gui-né-a », \esp{orgullo} « or-gou-llo » (le \esp{ll} se prononce comme un « y » doux, comme dans « fille »).
\item \textbf{Lecture par des élèves volontaires}, phrase par phrase, avec correction de la prononciation.
\end{enumerate}
\end{experience}

\courssub{Le texte}

\begin{textebox}[Mi patria, mis culturas]
Me llamo André y vivo en Yaoundé, la capital de Camerún. El 20 de mayo es un día muy especial para mi país: es la fiesta nacional. Ese día, los alumnos de todos los colegios desfilan con orgullo delante de la bandera verde, roja y amarilla. Cantamos el himno nacional en francés y en inglés, porque Camerún es un país bilingüe. Mi padre dice que el verde representa los bosques del sur, el rojo la unidad y el amarillo el sol y las sabanas del norte.

Tengo un amigo mexicano, Carlos, que estudia conmigo. Él me cuenta que el 16 de septiembre los mexicanos celebran su independencia con el famoso «Grito de Dolores». Esa noche, la gente grita «¡Viva México!» en las plazas, come platos típicos y baila hasta muy tarde. Carlos está muy orgulloso de su patria y de sus tradiciones.

También tengo una amiga de Malabo, en Guinea Ecuatorial. Ella me explica que en su país la gente habla español, fang y bubi: ¡qué riqueza! El español es la lengua oficial, pero las lenguas africanas viven en las canciones, en los cuentos y en las fiestas de los pueblos.

Cuando comparo estos tres países, comprendo una cosa importante: nuestras tradiciones, nuestras lenguas y nuestras fiestas son diferentes, pero el amor a la patria es el mismo. La bandera, el himno y las fiestas son símbolos que unen a los ciudadanos. Y la diversidad cultural no es un problema: es un tesoro, un patrimonio que debemos proteger y transmitir a nuestros hijos.

Por eso digo con orgullo: mi patria es Camerún, pero mi corazón habla todas las lenguas del mundo.
\end{textebox}

\textbf{Traduction du texte.}

Je m'appelle André et je vis à Yaoundé, la capitale du Cameroun. Le 20 mai est un jour très spécial pour mon pays : c'est la fête nationale. Ce jour-là, les élèves de tous les collèges défilent avec fierté devant le drapeau vert, rouge et jaune. Nous chantons l'hymne national en français et en anglais, parce que le Cameroun est un pays bilingue. Mon père dit que le vert représente les forêts du sud, le rouge l'unité et le jaune le soleil et les savanes du nord.

J'ai un ami mexicain, Carlos, qui étudie avec moi. Il me raconte que le 16 septembre, les Mexicains célèbrent leur indépendance avec le célèbre « Cri de Dolores ». Cette nuit-là, les gens crient « Vive le Mexique ! » sur les places, mangent des plats typiques et dansent jusqu'à très tard. Carlos est très fier de sa patrie et de ses traditions.

J'ai aussi une amie de Malabo, en Guinée équatoriale. Elle m'explique que dans son pays, les gens parlent espagnol, fang et bubi : quelle richesse ! L'espagnol est la langue officielle, mais les langues africaines vivent dans les chansons, dans les contes et dans les fêtes des villages.

Quand je compare ces trois pays, je comprends une chose importante : nos traditions, nos langues et nos fêtes sont différentes, mais l'amour de la patrie est le même. Le drapeau, l'hymne et les fêtes sont des symboles qui unissent les citoyens. Et la diversité culturelle n'est pas un problème : c'est un trésor, un patrimoine que nous devons protéger et transmettre à nos enfants.

C'est pourquoi je dis avec fierté : ma patrie est le Cameroun, mais mon cœur parle toutes les langues du monde.

\courssub{Glossaire et lexique du texte}

\begin{center}
\begin{tabularx}{\linewidth}{|X|X|X|}
\hline
\rowcolor{popblueL} \textbf{Español} & \textbf{Français} & \textbf{Exemple du texte} \\
\hline
la patria & la patrie & \esp{el amor a la patria} \\
\hline
la bandera & le drapeau & \esp{la bandera verde, roja y amarilla} \\
\hline
el himno (nacional) & l'hymne (national) & \esp{cantamos el himno nacional} \\
\hline
desfilar & défiler & \esp{los alumnos desfilan} \\
\hline
el orgullo / orgulloso(a) & la fierté / fier(ère) & \esp{con orgullo} ; \esp{está orgulloso} \\
\hline
la fiesta nacional & la fête nationale & \esp{es la fiesta nacional} \\
\hline
la lengua & la langue & \esp{la lengua oficial} \\
\hline
la tradición & la tradition & \esp{nuestras tradiciones} \\
\hline
el patrimonio & le patrimoine & \esp{un patrimonio que debemos proteger} \\
\hline
la diversidad & la diversité & \esp{la diversidad cultural} \\
\hline
el mestizaje & le métissage & \esp{el mestizaje de las culturas} \\
\hline
el tesoro & le trésor & \esp{es un tesoro} \\
\hline
unir & unir & \esp{símbolos que unen} \\
\hline
proteger & protéger & \esp{debemos proteger} \\
\hline
gritar / el grito & crier / le cri & \esp{el Grito de Dolores} \\
\hline
bilingüe & bilingue & \esp{un país bilingüe} \\
\hline
la sierra / la sabana & la chaîne de montagnes / la savane & \esp{las sabanas del norte} \\
\hline
transmitir & transmettre & \esp{transmitir a nuestros hijos} \\
\hline
\end{tabularx}
\end{center}

\begin{attention}[Faux amis et pièges de traduction]
\begin{itemize}
\item \esp{la patria} = la patrie (et non « le père » : le père se dit \esp{el padre}).
\item \esp{la lengua} = la langue (au sens de langue parlée ET d'organe) ; ne pas confondre avec \esp{el lenguaje} (le langage, la faculté de langage).
\item \esp{gritar} = crier, pousser un cri ; ne pas traduire \esp{el grito} par « le gril » (insecte) !
\item \esp{está orgulloso} : on utilise \esp{estar} car c'est un état, un sentiment passager (nous y reviendrons dans la section grammaire).
\item Attention au genre : \esp{el himno} est masculin, \esp{la bandera} est féminin, \esp{el orgullo} masculin, \esp{la diversidad} féminin.
\item \esp{la sierra} = chaîne de montagnes ; \esp{la sabana} = savane. Ne pas confondre \esp{sabana} (savane) et \esp{sábana} (drap, avec accent sur le premier a).
\end{itemize}
\end{attention}

\courssub{Repérage des éléments de situation}

Après la lecture, répondons aux quatre questions fondamentales qui permettent de situer n'importe quel texte.

\begin{center}
\begin{tabularx}{\linewidth}{|X|X|}
\hline
\rowcolor{popgreenL} \textbf{Question} & \textbf{Réponse} \\
\hline
¿Quién habla? (Qui parle?) & André, un jeune Camerounais de Yaoundé, qui témoigne à la première personne (\esp{me llamo}, \esp{vivo}, \esp{cantamos}, \esp{digo}). \\
\hline
¿De qué habla? (De quoi parle-t-il?) & De la fête nationale camerounaise du 20 mai, de la fête de l'Indépendance mexicaine (16 sept.) et des langues de la Guinée équatoriale. \\
\hline
¿Dónde? (Où?) & À Yaoundé (Camerún), avec des évocations du Mexique et de Malabo (Guinea Ecuatorial). \\
\hline
¿Cuándo? (Quand?) & Aujourd'hui (présent de narration) ; les fêtes évoquées ont lieu le \esp{20 de mayo} (Camerún) et le \esp{16 de septiembre} (México). \\
\hline
\end{tabularx}
\end{center}

\courssub{Le genre du texte et son thème}

\begin{definition}[Récit-témoignage à visée argumentative]
Un \cle{récit-témoignage} est un texte où le narrateur raconte une expérience vécue à la première personne (\esp{yo}). Il est dit « à visée argumentative » lorsque, derrière le récit, l'auteur veut défendre une idée, convaincre le lecteur. Ici, André raconte son expérience des fêtes nationales \textbf{pour défendre une thèse} : la diversité culturelle est une richesse et non un problème.
\end{definition}

Les indices du genre dans le texte :
\begin{itemize}
\item Première personne : \esp{me llamo André}, \esp{vivo en Yaoundé}, \esp{mi corazón}.
\item Faits vécus et concrets : le défilé du 20 mai, le récit de Carlos, les explications de l'amie de Malabo.
\item Prise de position argumentative : \esp{comprendo una cosa importante}, \esp{la diversidad cultural no es un problema: es un tesoro}, \esp{por eso digo con orgullo}.
\end{itemize}

\begin{aretenir}[Le thème du texte]
Le thème principal est \textbf{l'amour de la patrie à travers la richesse des cultures} : chaque pays a ses symboles (drapeau, hymne), ses fêtes et ses langues, mais le sentiment patriotique est universel. L'idée directrice : \esp{La diversidad cultural no es un problema: es un tesoro que debemos proteger.} « La diversité culturelle n'est pas un problème : c'est un trésor que nous devons protéger. »
\end{aretenir}

\begin{definition}[Patriotismo]
Le \cle{patriotismo} est le sentiment d'amour et de fierté envers la patrie, ses symboles et sa culture. En espagnol : \esp{sentimiento de amor y de orgullo hacia la patria, sus símbolos y su cultura}. Exemple : \esp{El patriotismo se vive intensamente durante la fiesta nacional.}
\end{definition}

\begin{definition}[Diversidad cultural]
La \cle{diversidad cultural} est la coexistence de langues, de traditions et de coutumes différentes à l'intérieur d'un pays ou entre les pays. En espagnol : \esp{coexistencia de lenguas, tradiciones y costumbres diferentes dentro de un país o entre países}. Exemple : \esp{Camerún es un ejemplo de diversidad cultural por sus más de 200 lenguas.}
\end{definition}

\courssub{Carte mentale du lexique}

\begin{popfigure}
\begin{tikzpicture}[mindmap, grow cyclic, every node/.style={concept, execute at begin node=\hskip0pt}, concept color=popblue!60, text=white,
  level 1/.append style={level distance=3.2cm, sibling angle=90},
  level 2/.append style={level distance=2.2cm, sibling angle=45}]
\node[concept, font=\large\bfseries] {Patriotismo y diversidad}
  child[concept color=poporange!80] { node {Símbolos}
    child { node {bandera} }
    child { node {himno} }
    child { node {escudo} }
  }
  child[concept color=popgreen!80] { node {Tradiciones}
    child { node {fiestas} }
    child { node {lenguas} }
    child { node {gastronomía} }
  }
  child[concept color=poppurple!80] { node {Países}
    child { node[align=center] {Camerún} }
    child { node[align=center] {México} }
    child { node[align=center] {Guinea\\Ecuatorial} }
    child { node[align=center] {España} }
  }
  child[concept color=poppink!80] { node {Valores}
    child { node {orgullo} }
    child { node {respeto} }
    child { node {unidad} }
    child { node {mestizaje} }
  };
\end{tikzpicture}
\legende{Carte mentale du champ lexical : le patriotisme et la diversité culturelle.}
\end{popfigure}

\courssub{Questions de compréhension}

\begin{exercice}[Compréhension globale du texte]
Réponds en espagnol, avec des phrases complètes :
\begin{enumerate}
\item ¿Quién es el narrador y dónde vive?
\item ¿Qué hacen los alumnos el 20 de mayo en Yaoundé?
\item ¿Por qué los alumnos cantan el himno en dos lenguas?
\item ¿Qué celebran los mexicanos el 16 de septiembre?
\item ¿Qué lenguas se hablan en Guinea Ecuatorial?
\item Según el narrador, ¿la diversidad cultural es un problema?
\end{enumerate}
\end{exercice}

\begin{exoresolu}[Question de compréhension résolue]
Énoncé : ¿Por qué los alumnos cantan el himno nacional en francés y en inglés?
\tcblower
Solution détaillée : il faut d'abord repérer dans le texte la phrase qui contient la réponse : \esp{Cantamos el himno nacional en francés y en inglés, porque Camerún es un país bilingüe.} Le connecteur \esp{porque} (« parce que ») introduit la cause. On répond donc avec une phrase complète qui reprend la question :

\esp{Los alumnos cantan el himno nacional en francés y en inglés porque Camerún es un país bilingüe.}

« Les élèves chantent l'hymne national en français et en anglais parce que le Cameroun est un pays bilingue. »

Méthode : ne jamais répondre par un simple « oui » ou « non » ; reprendre les mots de la question et ajouter l'information du texte.
\end{exoresolu}

\begin{attention}[Erreurs fréquentes des francophones à l'écrit]
\begin{itemize}
\item Ne dis pas \esp{el 20 de mayo es la fiesta de Camerún es muy especial} : évite les phrases à deux verbes conjugués sans lien. Dis : \esp{El 20 de mayo es un día muy especial para mi país.}
\item Les dates : en espagnol, on dit \esp{el 20 de mayo} avec l'article \esp{el} et la préposition \esp{de}, jamais \esp{en 20 mayo} ni \esp{el 20 mayo}.
\item Les nationalités et les langues ne prennent pas de majuscule en espagnol : \esp{mexicano}, \esp{español}, \esp{francés}, \esp{camerunés}. Seuls les noms de pays prennent la majuscule : \esp{México}, \esp{Camerún}, \esp{España}.
\item N'oublie jamais les points d'exclamation et d'interrogation ouvrants : \esp{¡qué riqueza!}, \esp{¿dónde vives?}, \esp{¡Viva México!}.
\item \esp{Guinea Ecuatorial} s'écrit sans accent sur \esp{Guinea} et avec accent sur \esp{Ecuatorial} (sur le \esp{a}).
\end{itemize}
\end{attention}

\begin{savaistu}[La Guinée équatoriale, pont entre le Cameroun et le monde hispanophone]
La Guinée équatoriale est le seul pays d'Afrique où l'espagnol est langue officielle (depuis 1844). Ce petit pays, voisin du Cameroun (frontière terrestre de 189 km), partage aussi des langues africaines comme le fang, parlé de part et d'autre de la frontière. C'est donc un pont naturel entre le Cameroun et le monde hispanophone : en apprenant l'espagnol, tu peux communiquer avec tes voisins de Malabo et de Bata, mais aussi avec plus de vingt pays d'Espagne et d'Amérique latine. Le mot \esp{mestizaje} (métissage culturel) est né dans ce monde hispanophone pour décrire le mélange des cultures européennes, amérindiennes et africaines.
\end{savaistu}

\begin{aretenir}[Ce qu'il faut retenir de cette section]
\begin{itemize}
\item Le texte \esp{Mi patria, mis culturas} est un récit-témoignage à visée argumentative : André raconte son expérience pour défendre l'idée que la diversité culturelle est un trésor.
\item Les mots-clés du thème : \esp{patria}, \esp{bandera}, \esp{himno}, \esp{tradición}, \esp{fiesta}, \esp{lengua}, \esp{patrimonio}, \esp{diversidad}, \esp{mestizaje}, \esp{orgullo}.
\item Pour situer un texte, toujours répondre à quatre questions : ¿quién habla? ¿de qué? ¿dónde? ¿cuándo?
\item Le Cameroun (20 mai), le Mexique (16 septembre) et la Guinée équatoriale montrent trois visages du même amour de la patrie.
\end{itemize}
\end{aretenir}

\courssub{Grammaire outillée : \textit{ser} / \textit{estar} et le présent de l'indicatif}

\begin{propriete}[Le grand dilemme : \textit{ser} ou \textit{estar} ?]
En français, le verbe « être » couvre tout. En espagnol, le choix entre \cle{ser} et \cle{estar} change le sens de la phrase.
\begin{itemize}
\item \textbf{SER} : identité, origine, nationalité, caractéristique essentielle, définition, heure, possession, matière. \rightarrow \emph{C'est « ce que je suis » de façon permanente.}
\item \textbf{ESTAR} : localisation (où?), état physique ou moral passager, résultat d'action, position. \rightarrow \emph{C'est « comment je suis » ou « où je suis » à un moment donné.}
\end{itemize}
\end{propriete}

\begin{center}
\begin{tabular}{|c|c|c|}
\hline
\rowcolor{popblueL} \textbf{Contexte} & \textbf{Verbe} & \textbf{Exemples du texte / du thème} \\
\hline
Identité / Définition & \esp{ser} & \esp{Mi patria \textbf{es} Camerún.} / \esp{El 20 de mayo \textbf{es} la fiesta nacional.} \\
\hline
Origine / Nationalité & \esp{ser} & \esp{Carlos \textbf{es} mexicano.} / \esp{Yo \textbf{soy} camerunés.} \\
\hline
Caractéristique essentielle & \esp{ser} & \esp{La diversidad cultural \textbf{es} un tesoro.} / \esp{El himno \textbf{es} bilingüe.} \\
\hline
Heure / Date & \esp{ser} & \esp{\textbf{Son} las diez.} / \esp{Hoy \textbf{es} 20 de mayo.} \\
\hline
Localisation (Où?) & \esp{estar} & \esp{Yaoundé \textbf{está} en Camerún.} / \esp{Malabo \textbf{está} en Guinea Ecuatorial.} \\
\hline
État / Sentiment passager & \esp{estar} & \esp{Carlos \textbf{está} orgulloso.} / \esp{Estamos contentos.} \\
\hline
Action en cours (gérondif) & \esp{estar} + gérondif & \esp{Los alumnos \textbf{están desfilando}.} \\
\hline
\end{tabular}
\end{center}

\begin{attention}[Pièges classiques \textit{ser} / \textit{estar}]
\begin{itemize}
\item \esp{Ser aburrido} = être une personne ennuyeuse (caractère). \quad \esp{Estar aburrido} = s'ennuyer (état momentané).
\item \esp{Ser listo} = être intelligent. \quad \esp{Estar listo} = être prêt.
\item \esp{Ser verde} = être de couleur verte (le drapeau). \quad \esp{Estar verde} = ne pas être mûr (fruit) / être novice.
\item \textbf{Localisation des événements} : Pour situer un événement (fête, concert, match), on utilise \esp{ser} : \esp{La fiesta \textbf{es} en el estadio.} (Mais : \esp{El estadio \textbf{está} en el centro.})
\item \textbf{Mort} : \esp{Estar muerto} (état résultatif), jamais \esp{ser muerto}.
\end{itemize}
\end{attention}

\begin{propriete}[Conjugaison au présent de l'indicatif (verbes du thème)]
Le présent sert à décrire des habitudes, des vérités générales, des traditions ancrées.
\end{propriete}

\begin{center}
\begin{tabular}{|l|l|l|l|}
\hline
\rowcolor{popblueL} \textbf{Personne} & \textbf{ser (irrégulier)} & \textbf{estar (irrégulier)} & \textbf{vivir / cantar / hablar (réguliers)} \\
\hline
yo & soy & estoy & vivo / canto / hablo \\
\hline
tú & eres & estás & vives / cantas / hablas \\
\hline
él / ella / usted & es & está & vive / canta / habla \\
\hline
nosotros / as & somos & estamos & vivimos / cantamos / hablamos \\
\hline
vosotros / as & sois & estáis & vivís / cantáis / habláis \\
\hline
ellos / ellas / ustedes & son & están & viven / cantan / hablan \\
\hline
\end{tabular}
\end{center}

\begin{exemplebox}[Exemples authentiques au présent]
\begin{itemize}
\item \esp{Los cameruneses \textbf{cantamos} el himno en dos lenguas.} (Habitude, fait culturel récurrent)
\item \esp{En México, la gente \textbf{grita} «¡Viva México!» el 16 de septiembre.} (Tradition annuelle)
\item \esp{El español \textbf{es} la lengua oficial de Guinea Ecuatorial.} (Fait permanent, identité)
\item \esp{Las lenguas africanas \textbf{viven} en las canciones.} (Métaphore : présence continue)
\item \esp{Mi amigo Carlos \textbf{está} orgulloso de su patria.} (État actuel, sentiment)
\end{itemize}
\end{exemplebox}

\begin{exercice}[Entraînement : \textit{ser} ou \textit{estar} au présent ?]
Complète par la forme correcte du verbe entre parenthèses.
\begin{enumerate}
\item Camerún \_\_\_\_\_\_ (ser/estar) un país bilingüe.
\item Yaoundé \_\_\_\_\_\_ (ser/estar) la capital de Camerún.
\item Los alumnos \_\_\_\_\_\_ (ser/estar) orgullosos el 20 de mayo.
\item El himno nacional \_\_\_\_\_\_ (ser/estar) un símbolo de la patria.
\item Malabo \_\_\_\_\_\_ (ser/estar) en la isla de Bioko.
\item ¿Qué día \_\_\_\_\_\_ (ser/estar) hoy? — Hoy \_\_\_\_\_\_ (ser/estar) 20 de mayo.
\item Las tradiciones \_\_\_\_\_\_ (ser/estar) un patrimonio que \_\_\_\_\_\_ (deber) proteger.
\end{enumerate}
\end{exercice}

\begin{corrige}[Exercice Grammaire : \textit{ser} / \textit{estar}]
\begin{enumerate}
\item \textbf{es} (caractéristique essentielle / définition)
\item \textbf{está} (localisation géographique d'une ville)
\item \textbf{están} (état passager, sentiment à un moment donné)
\item \textbf{es} (identité, définition d'un symbole)
\item \textbf{está} (localisation)
\item \textbf{es} / \textbf{es} (date / heure → toujours \textit{ser})
\item \textbf{son} (caractéristique) / \textbf{debemos} (verbe \textit{deber} au présent, 1ère personne du pluriel : obligation morale)
\end{enumerate}
\end{corrige}

\courssub{Méthode du commentaire de texte (Épreuve écrite)}

\begin{methode}[Commenter un texte en espagnol (Bac / Évaluation)]
\begin{enumerate}
\item \textbf{Introduction} (3-4 lignes) :
  \begin{itemize}
  \item \textbf{Accroche} : une phrase sur le thème général (patriotisme, diversité).
  \item \textbf{Présentation} : titre (entre guillemets), auteur (si connu, ici « texte anonyme / didactique »), nature du texte (récit-témoignage argumentatif), source (manuel / cours).
  \item \textbf{Problématique} : la question à laquelle le texte répond (ex. : \esp{¿Cómo se expresa el patriotismo en contextos culturales diversos?}).
  \item \textbf{Annonce du plan} : \esp{En primer lugar...; En segundo lugar...; Finalmente...}
  \end{itemize}
\item \textbf{Développement} (2 ou 3 parties, 1 paragraphe chacune) :
  \begin{itemize}
  \item Une idée directrice par paragraphe.
  \item Argumentation : \esp{En el texto se dice que...}, \esp{El autor afirma que...}, \esp{Por ejemplo...}
  \item Citation courte intégrée : \esp{«la diversidad cultural no es un problema»}.
  \item Analyse : expliquer la portée de la citation, le vocabulaire utilisé (\esp{tesoro, patrimonio, unir}).
  \item Connecteurs logiques : \esp{Por un lado... por otro lado...}, \esp{Además}, \esp{Sin embargo}, \esp{Por tanto}.
  \end{itemize}
\item \textbf{Conclusion} (2-3 lignes) :
  \begin{itemize}
  \item \textbf{Bilan} : réponse synthétique à la problématique.
  \item \textbf{Ouverture} : lien avec l'actualité, expérience personnelle, autre pays hispanophone (ex. : l'Espagne et ses \esp{comunidades autónomas}, la \esp{Hispanidad}).
  \end{itemize}
\end{enumerate}
\end{methode}

\courssub{Commentaire modèle : « Mi patria, mis culturas »}

\begin{textebox}[Commentaire modèle (niveau A2+/B1)]
\textbf{Introducción}
El patriotismo no es solo ondear una bandera; es un sentimiento que se expresa a través de la cultura. El texto «Mi patria, mis culturas» es un relato testimonial argumentativo escrito con fines didácticos. Presenta la experiencia de André, un joven camerunés, para reflexionar sobre la relación entre amor a la patria y diversidad cultural. La pregunta central es: ¿puede la diversidad cultural fortalecer la unidad nacional?

\textbf{Desarrollo}
\textbf{1. Los símbolos patrios como base de la identidad colectiva.}
En el primer párrafo, André describe la fiesta nacional del 20 de mayo en Yaoundé. Los símbolos unen a los ciudadanos: «los alumnos desfilan con orgullo delante de la bandera» y «cantamos el himno nacional en francés y en inglés». El uso del presente de indicativo (\textit{desfilan, cantamos, es}) muestra que son ritos vivos y actuales. Los colores de la bandera tienen un significado geográfico y moral: el verde, el rojo, el amarillo. Así, el patriotismo se enseña y se vive en la escuela.

\textbf{2. La diversidad en el mundo hispano: México y Guinea Ecuatorial.}
El narrador abre su horizonte comparando su país con México y Guinea Ecuatorial. En México, el «Grito de Dolores» (\textit{gritan «¡Viva México!»}) es una tradición histórica que genera orgullo (\textit{está orgulloso}). En Guinea Ecuatorial, la riqueza es lingüística: «hablan español, fang y bubi». Aquí el español es lengua oficial, pero las lenguas africanas «viven en las canciones». La diversidad no divide, enriquece.

\textbf{3. La tesis: la diversidad es un tesoro, no un problema.}
En el párrafo final, André argumenta: «nuestras tradiciones... son diferentes, pero el amor a la patria es el mismo». Los símbolos (\textit{banderas, himnos, fiestas}) «unen a los ciudadanos». La frase clave resume la tesis: «la diversidad cultural no es un problema: es un tesoro, un patrimonio que debemos proteger». El verbo \textit{deber} + infinitivo expresa un deber moral colectivo.

\textbf{Conclusión}
En definitiva, el texto demuestra que el patriotismo cultural se construye sobre el respeto a la diferencia. Camerún, México y Guinea Ecuatorial comparten el mismo amor a la patria expresado de formas distintas. Como dice André: «mi corazón habla todas las lenguas del mundo». Esta lección vale para todo el mundo hispano: la unidad no exige uniformidad.
\end{textebox}

\courssub{Production guidée : Présenter une fête ou une tradition}

\begin{experience}[Tâche finale : « Mi fiesta, mi orgullo » — Expression orale / écrite]
\textbf{Consigne :} Choisis une fête traditionnelle du Cameroun (ex. : Ngondo, Medumba, Nyem-Nyem, fête des moissons, Tabaski, Noël, 20 mai) ou d'un pays hispanophone que tu connais. Rédige un paragraphe de 80-100 mots en espagnol pour la présenter à des correspondants mexicains ou équato-guinéens. Utilise le vocabulaire de la leçon et la grammaire travaillée (\textit{ser/estar}, présent).

\textbf{Structure suggérée :}
\begin{enumerate}
\item \textbf{Nom et date} : \esp{Mi fiesta favorita es... Se celebra el...}
\item \textbf{Lieu et participants} : \esp{Tiene lugar en... Participan...}
\item \textbf{Activités (présent habitude)} : \esp{La gente baila / come / canta / viste trajes típicos...}
\item \textbf{Sens / Symboles (ser)} : \esp{Es una tradición muy importante porque... Representa...}
\item \textbf{Sentiment personnel (estar)} : \esp{Yo estoy muy orgulloso/a de esta fiesta.}
\end{enumerate}

\textbf{Exemple de production élève (modèle) :}
\esp{Mi fiesta favorita es el Ngondo, la gran fiesta del pueblo Sawa en Douala. Se celebra cada año en diciembre a orillas del río Wouri. Ese día, la gente viste telas tradicionales coloridas y los jefes tradicionales llegan en piraguas decoradas. La gente baila el essewe y canta en duala. Es una tradición muy importante porque une a las familias y honra a los antepasados. El río es sagrado: es el símbolo de la vida. Yo estoy muy orgulloso de esta fiesta porque muestra la riqueza de mi cultura. ¡Viva el Ngondo!}

\textbf{Évaluation (grille simplifiée) :}
\begin{itemize}
\item Contenu / Respect de la structure : 4 pts
\item Lexique du thème (patria, tradición, orgullo, símbolo, patrimonio...) : 4 pts
\item Grammaire : \textit{ser/estar} corrects, présent indicatif régulier/irrégulier : 4 pts
\item Orthographe (accents, ñ, ¿ ¡, majuscules pays / minuscules nationalités) : 3 pts
\item Cohérence / Connecteurs (porque, además, por eso) : 3 pts
\item \textbf{Total : 18/20} (convertir sur 20 selon barème établissement)
\end{itemize}
\end{experience}

\begin{aretenir}[Synthèse de la leçon EA09]
\begin{itemize}
\item \textbf{Lexique} : \esp{patria, bandera, himno, tradición, fiesta, lengua, patrimonio, diversidad, mestizaje, orgullo, grito, bilingüe}.
\item \textbf{Grammaire} : Distinction \textbf{ser} (identité, origine, essence, heure) / \textbf{estar} (localisation, état, gérondif). Présent de l'indicatif pour les habitudes et traditions.
\item \textbf{Culture} : 20 mai (Camerún), 16 septembre (México), Guinée équatoriale (seul pays hispanophone d'Afrique, langues fang/bubi).
\item \textbf{Compétences} : Lire un texte argumentatif, situer un texte (Qui/Quoi/Où/Quand), commenter (Intro/Développement/Conclusion), présenter une fête à l'oral/écrit.
\end{itemize}
\end{aretenir}

\coursec{Compréhension du texte}

Dans la section précédente, nous avons lu et écouté le texte \esp{« Mi patria, mis culturas »}, dans lequel un élève camerounais raconte comment sa classe célèbre la fête nationale et découvre les traditions des pays hispanophones. Avant d'étudier le vocabulaire et la grammaire, il faut maintenant vérifier que nous avons bien \cle{compris} le texte : c'est l'étape de la \cle{comprensión lectora}. Nous allons répondre à des questions globales, puis à des questions de détail, puis interpréter une phrase clé, et enfin repérer les champs lexicaux.

\begin{textebox}[Rappel du texte : « Mi patria, mis culturas » (extraits)]
Me llamo André y soy alumno de una escuela de Yaundé. El veinte de mayo celebramos la fiesta nacional de Camerún. Ese día, cantamos el himno nacional, saludamos la bandera y desfilamos delante de los profesores. Estoy muy orgulloso de mi país.

Mi profesora de español nos dice que la diversidad cultural es un tesoro. En Camerún hablamos muchas lenguas: el francés, el inglés y más de doscientas lenguas nacionales. Cada región tiene sus fiestas y sus tradiciones.

En los países hispanohablantes también hay mucha diversidad. El dieciséis de septiembre, los mexicanos celebran su independencia con música y fuegos artificiales. En España, cada comunidad tiene sus fiestas: los sanfermines en Pamplona, la feria de Sevilla... En Guinea Ecuatorial, un país africano como el mío, la gente habla español, francés, portugués y lenguas africanas como el fang y el bubi.

Mi patria es Camerún, pero mi corazón habla muchas lenguas. La diversidad no nos divide: nos une.
\end{textebox}

\textbf{Glossaire :} \esp{alumno} = élève ; \esp{himno} = hymne ; \esp{bandera} = drapeau ; \esp{desfilar} = défiler ; \esp{orgulloso} = fier ; \esp{tesoro} = trésor ; \esp{fuegos artificiales} = feux d'artifice ; \esp{corazón} = cœur ; \esp{unir} = unir.

\courssub{Questions de compréhension globale}

Les questions globales portent sur le texte entier : le thème, les personnages, les lieux. On y répond sans entrer dans les détails.

\begin{methode}[Comment répondre à une question de compréhension]
\begin{enumerate}
\item Lis la question et souligne le mot interrogatif : \esp{¿quién?} (qui), \esp{¿qué?} (que), \esp{¿dónde?} (où), \esp{¿cuándo?} (quand), \esp{¿por qué?} (pourquoi).
\item Repère dans le texte la phrase qui contient la réponse.
\item Rédige une \textbf{phrase complète en espagnol}, en reprenant les mots de la question : \esp{¿Dónde estudia André?} → \esp{André estudia en una escuela de Yaundé.}
\item Cite un mot ou un groupe de mots du texte pour justifier ta réponse.
\end{enumerate}
\end{methode}

\begin{exemplebox}[Questions globales et réponses modèles]
\textbf{1. ¿De qué habla el texto?} (De quoi parle le texte ?)

\esp{El texto habla del patriotismo y de la diversidad cultural en Camerún y en los países hispanohablantes.} « Le texte parle du patriotisme et de la diversité culturelle au Cameroun et dans les pays hispanophones. »

\textbf{2. ¿Quién es el narrador?} (Qui est le narrateur ?)

\esp{El narrador es André, un alumno de una escuela de Yaundé.} « Le narrateur est André, un élève d'une école de Yaoundé. »

\textbf{3. ¿Qué países se mencionan en el texto?} (Quels pays sont mentionnés ?)

\esp{El texto menciona Camerún, México, España y Guinea Ecuatorial.} « Le texte mentionne le Cameroun, le Mexique, l'Espagne et la Guinée équatoriale. »
\end{exemplebox}

\courssub{Questions de détail}

Les questions de détail exigent de retrouver une \cle{information précise} : une date, une action, une liste. Il faut relire le passage concerné.

\begin{exemplebox}[Questions de détail et réponses modèles]
\textbf{1. ¿Qué hacen los alumnos el veinte de mayo?} (Que font les élèves le 20 mai ?)

\esp{El veinte de mayo, los alumnos cantan el himno nacional, saludan la bandera y desfilan delante de los profesores.} « Le 20 mai, les élèves chantent l'hymne national, saluent le drapeau et défilent devant les professeurs. »

\textbf{2. ¿Qué se celebra el dieciséis de septiembre en México?} (Que célèbre-t-on le 16 septembre au Mexique ?)

\esp{El dieciséis de septiembre, los mexicanos celebran su independencia con música y fuegos artificiales.} « Le 16 septembre, les Mexicains célèbrent leur indépendance avec de la musique et des feux d'artifice. »

\textbf{3. ¿Qué lenguas se hablan en Guinea Ecuatorial?} (Quelles langues parle-t-on en Guinée équatoriale ?)

\esp{En Guinea Ecuatorial, la gente habla español, francés, portugués y lenguas africanas como el fang y el bubi.} « En Guinée équatoriale, les gens parlent espagnol, français, portugais et des langues africaines comme le fang et le bubi. »

\textbf{4. ¿Qué fiestas hay en España?}

\esp{En España hay fiestas como los sanfermines en Pamplona y la feria de Sevilla.} « En Espagne, il y a des fêtes comme les sanfermines à Pampelune et la foire de Séville. »
\end{exemplebox}

\begin{exemplebox}[Deux phrases clés traduites]
\esp{Cantamos el himno nacional.} = « Nous chantons l'hymne national. »

\esp{La diversidad cultural es un tesoro.} = « La diversité culturelle est un trésor. »
\end{exemplebox}

\courssub{Question d'interprétation}

La question d'interprétation demande d'aller \textbf{au-delà du texte} : il faut expliquer une image, une opinion, une intention de l'auteur. La réponse n'est pas écrite mot pour mot dans le texte ; il faut la construire.

\begin{exoresolu}[Pourquoi la diversité est-elle « un tesoro » ?]
\textbf{Question :} \esp{¿Por qué dice el narrador que la diversidad cultural es « un tesoro »?} (Pourquoi le narrateur dit-il que la diversité culturelle est « un trésor » ?)
\tcblower
\textbf{Solution détaillée :}

1. Un \esp{tesoro} est quelque chose de très précieux, comme l'or ou les bijoux. Le narrateur utilise donc une \textbf{comparaison} : la diversité culturelle est précieuse comme un trésor.

2. Le texte donne des indices : \esp{« cada región tiene sus fiestas y sus tradiciones »} et \esp{« la diversidad no nos divide: nos une »}.

3. Réponse complète en espagnol : \esp{El narrador dice que la diversidad cultural es un tesoro porque las lenguas, las fiestas y las tradiciones de cada pueblo son una riqueza preciosa. Además, la diversidad no divide a la gente: la une.} « Le narrateur dit que la diversité culturelle est un trésor parce que les langues, les fêtes et les traditions de chaque peuple sont une richesse précieuse. De plus, la diversité ne divise pas les gens : elle les unit. »
\end{exoresolu}

\begin{attention}[Ne traduis pas mot à mot]
Dans tes réponses, \textbf{ne traduis pas le français mot à mot} et ne recopie pas de longs passages du texte sans les comprendre. Reformule avec des mots simples que tu maîtrises. Par exemple, pour dire « la diversité est une richesse », écris simplement \esp{la diversidad es una riqueza}, et non une phrase compliquée calquée du français. Attention aussi aux faux amis : \esp{asistir a una fiesta} signifie « assister à une fête », mais \esp{una fiesta} ne se traduit jamais par « une fieste » ; et \esp{la gente} (les gens) est singulier en espagnol : \esp{la gente habla}, pas \esp{la gente hablan}.
\end{attention}

\courssub{Vrai ou faux : justifier par le texte}

Pour chaque affirmation, écris \esp{Verdadero} (vrai) ou \esp{Falso} (faux), puis \textbf{cite la phrase du texte} qui le prouve.

\begin{exercice}[¿Verdadero o falso?]
\begin{enumerate}
\item \esp{André es alumno de una escuela de Duala.}
\item \esp{El veinte de mayo, los alumnos cantan el himno nacional.}
\item \esp{En Camerún se hablan más de doscientas lenguas nacionales.}
\item \esp{Los mexicanos celebran su independencia el veinte de mayo.}
\item \esp{En Guinea Ecuatorial se habla español.}
\item \esp{Para el narrador, la diversidad divide a la gente.}
\end{enumerate}
\end{exercice}

\begin{corrige}[Exercice « ¿Verdadero o falso? »]
\begin{enumerate}
\item \textbf{Falso.} Le texte dit : \esp{« soy alumno de una escuela de Yaundé »} (Yaoundé, pas Douala).
\item \textbf{Verdadero.} \esp{« Ese día, cantamos el himno nacional »}.
\item \textbf{Verdadero.} \esp{« hablamos [...] más de doscientas lenguas nacionales »}.
\item \textbf{Falso.} \esp{« El dieciséis de septiembre, los mexicanos celebran su independencia »} (le 16 septembre, pas le 20 mai).
\item \textbf{Verdadero.} \esp{« la gente habla español, francés, portugués y lenguas africanas »}.
\item \textbf{Falso.} \esp{« La diversidad no nos divide: nos une »} (elle unit, elle ne divise pas).
\end{enumerate}
\end{corrige}

\courssub{Les champs lexicaux du texte}

Un \cle{champ lexical} est l'ensemble des mots d'un texte qui se rapportent à une même idée. Repérer les champs lexicaux aide à comprendre la structure du texte. Ici, trois champs dominent.

\begin{center}
\begin{tabularx}{\linewidth}{|X|X|X|}
\hline
\rowcolor{popblueL} \textbf{Champ lexical} & \textbf{Mots du texte} & \textbf{Traduction} \\
\hline
Symboles nationaux & \esp{patria, bandera, himno nacional, fiesta nacional} & patrie, drapeau, hymne national, fête nationale \\
\hline
Fêtes et traditions & \esp{fiesta, tradición, desfilar, fuegos artificiales, feria, sanfermines} & fête, tradition, défiler, feux d'artifice, foire, sanfermines \\
\hline
Langues & \esp{lengua, francés, inglés, español, portugués, fang, bubi} & langue, français, anglais, espagnol, portugais, fang, bubi \\
\hline
\end{tabularx}
\end{center}

On remarque que le champ des \textbf{langues} est le plus riche du texte : c'est un indice que le thème principal est la \cle{diversité}, plus que le seul patriotisme.

\courssub{Tableau comparatif à compléter : Camerún / Mundo hispanohablante}

Le texte met en parallèle le Cameroun et le monde hispanophone. Complétons ce tableau avec les fêtes, les langues et les symboles cités.

\begin{popfigure}
\begin{tikzpicture}[
  box/.style={draw=popblue, thick, rounded corners=3pt, fill=popblueL, align=left, text width=6.2cm, inner sep=6pt, font=\small},
  head/.style={draw=popdark, thick, rounded corners=3pt, fill=popblue, text=white, align=center, text width=6.2cm, inner sep=5pt, font=\small\bfseries}
]
\node[head] at (0,0) {Camerún};
\node[head] at (7.4,0) {Mundo hispanohablante};
\node[box, anchor=north, align=center] at (0,-0.5){\textbf{Fiesta:} el veinte de mayo, fiesta nacional\\[2pt]
\textbf{Symboles nationaux:} la bandera, el himno nacional\\[2pt]
\textbf{Lenguas:} francés, inglés y más de doscientas lenguas nacionales};
\node[box, anchor=north, align=center] at (7.4,-0.5){\textbf{Fiestas:} el dieciséis de septiembre (México), los sanfermines y la feria de Sevilla (España)\\[2pt]
\textbf{Éléments festifs:} música, fuegos artificiales\\[2pt]
\textbf{Lenguas:} español; en Guinea Ecuatorial también francés, portugués, fang y bubi};
\end{tikzpicture}
\legende{Tableau comparatif des fêtes, symboles et langues cités dans le texte}
\end{popfigure}

\begin{exercice}[À toi de compléter]
Recopie le tableau ci-dessus et ajoute une ligne personnelle : une fête de ton village ou de ta région (\esp{en mi región celebramos...}), une langue que tu parles (\esp{yo hablo...}) et un symbole que tu aimes (\esp{me gusta...}).
\end{exercice}

\courssub{Reformuler l'idée principale}

La dernière étape de la compréhension consiste à résumer le texte \textbf{en une seule phrase en espagnol}, avec tes propres mots. C'est un exercice essentiel pour le commentaire de texte.

\begin{methode}[Reformuler l'idée principale en une phrase]
\begin{enumerate}
\item Repère le sujet du texte : \esp{el patriotismo y la diversidad cultural}.
\item Repère l'opinion du narrateur : \esp{la diversidad es un tesoro, nos une}.
\item Assemble les deux dans une phrase simple avec \esp{porque} ou \esp{y} : sujet + verbe + opinion.
\end{enumerate}
\end{methode}

\begin{exemplebox}[Reformulations possibles]
\esp{André está orgulloso de Camerún y piensa que la diversidad cultural es un tesoro que une a la gente.} « André est fier du Cameroun et pense que la diversité culturelle est un trésor qui unit les gens. »

\esp{El texto muestra que el patriotismo y la diversidad de las lenguas y las fiestas son una riqueza para todos los países.} « Le texte montre que le patriotisme et la diversité des langues et des fêtes sont une richesse pour tous les pays. »
\end{exemplebox}

\begin{exoresolu}[Construis ta propre phrase de synthèse]
\textbf{Consigne :} Résume le texte en une phrase en espagnol, sans recopier aucune phrase entière du texte.
\tcblower
\textbf{Solution détaillée :}

1. Je choisis mon sujet : \esp{el narrador} ou \esp{el texto}.

2. Je choisis un verbe d'opinion : \esp{decir, pensar, mostrar, expliquer}.

3. Je formule : \esp{El texto explica que cada país tiene sus fiestas y sus lenguas, y que esta diversidad es una riqueza que une a las personas.} « Le texte explique que chaque pays a ses fêtes et ses langues, et que cette diversité est une richesse qui unit les personnes. »

Cette phrase est correcte, personnelle, et ne recopie pas le texte : c'est une bonne reformulation.
\end{exoresolu}

\begin{savaistu}[Le « Grito de Dolores »]
Le \esp{« Grito de Dolores »} (le « Cri de Dolores ») est l'appel lancé en 1810 par le prêtre Miguel Hidalgo dans le village de Dolores : c'est le cri qui a déclenché la guerre d'indépendance du Mexique. Aujourd'hui encore, chaque 15 septembre au soir, le président du Mexique répète ce cri depuis le balcon du palais national, à Mexico, devant une immense foule : \esp{« ¡Viva México! »}. C'est pourquoi la fête de l'indépendance mexicaine, célébrée le 16 septembre, commence en réalité la veille au soir, avec des cloches, des drapeaux et des feux d'artifice.
\end{savaistu}

\begin{aretenir}[L'essentiel de la compréhension]
\begin{itemize}
\item Réponds toujours par une \textbf{phrase complète en espagnol}, en reprenant les mots de la question.
\item Pour le vrai/faux, \textbf{cite la phrase du texte} qui justifie ta réponse.
\item Repère les \textbf{champs lexicaux} (symboles, fêtes, langues) pour comprendre le thème.
\item Termine par une \textbf{reformulation personnelle} de l'idée principale, en une phrase simple.
\item Évite la traduction mot à mot et les calques du français.
\end{itemize}
\end{aretenir}

Maintenant que le texte est bien compris, nous pouvons passer à l'étude de son vocabulaire : les mots de la patrie, des symboles et des traditions.

\coursec{Lexique thématique : patria, símbolos y tradiciones}

Dans la section précédente, nous avons lu et compris le texte \esp{Mi patria, mis culturas}. Nous y avons rencontré des mots comme \esp{bandera}, \esp{himno}, \esp{tradición} ou \esp{diversidad}. Il est temps maintenant de nous arrêter sur ce vocabulaire : le connaître par cœur, savoir l'employer dans une phrase et le classer par familles de sens. Un bon lexique, c'est la moitié d'un bon commentaire de texte et d'une bonne expression écrite au Bac.

\courssub{Observation : le vocabulaire dans son contexte}

Avant d'apprendre des listes de mots, observons-les d'abord dans un court texte. Lisez ce texte original et soulignez mentalement tous les mots qui parlent de la patrie, des fêtes et de la diversité.

\begin{textebox}[« Un país de muchas culturas » — texto de vocabulario]
Camerún es mi patria y estoy orgulloso de ella. El 20 de mayo, fiesta nacional, los ciudadanos desfilan delante de la bandera verde, roja y amarilla, y cantan el himno nacional con emoción. En las escuelas de Yaoundé y de Douala, los alumnos llevan ropa tradicional y presentan danzas de las diez regiones del país.

Camerún es un país bilingüe: se hablan el francés y el inglés, pero también más de doscientas cincuenta lenguas locales como el ewondo, el duala o el fulfuldé. Esta diversidad es una riqueza extraordinaria. Cada región tiene sus costumbres, su comida típica — el ndolé, el poulet DG, el eru — y sus fiestas tradicionales.

En los países hispanohablantes ocurre lo mismo. México celebra su independencia el 16 de septiembre con un gran desfile, y España tiene fiestas famosas como la Tomatina o la Semana Santa. Guinea Ecuatorial, país hispanohablante de África, mezcla las tradiciones africanas con la lengua española: es un ejemplo de mestizaje cultural.

El patrimonio de un pueblo es su tesoro. Debemos proteger nuestras lenguas, nuestras danzas y nuestras tradiciones, porque la unidad no significa que todos seamos iguales: significa vivir juntos en paz, con respeto. La diversidad es una riqueza, nunca un problema.
\end{textebox}

\textbf{Glossaire :} \esp{desfilar} = défiler ; \esp{llevar ropa} = porter des vêtements ; \esp{ocurrir} = se passer, se produire ; \esp{mezclar} = mélanger ; \esp{el tesoro} = le trésor ; \esp{significar} = signifier ; \esp{juntos} = ensemble.

\textbf{Questions rapides :} 1. ¿Cuándo es la fiesta nacional de Camerún? 2. ¿Qué lenguas se hablan en Camerún? 3. ¿Qué fiesta celebra México el 16 de septiembre? 4. Según el texto, ¿qué debemos proteger?

\courssub{Les quatre champs lexicaux}

\begin{definition}[Qu'est-ce qu'un champ lexical ?]
Un \cle{champ lexical} est un ensemble de mots qui se rapportent à une même idée ou à un même thème. Ici, le thème « patriotisme et diversité culturelle » se divise en quatre champs : les symboles de la patrie, les traditions et les fêtes, la diversité, et les sentiments et valeurs. Classer le vocabulaire par champs aide à le mémoriser et à le mobiliser dans une rédaction.
\end{definition}

\textbf{Champ lexical 1 — la patrie et les symboles}

\begin{itemize}
\item \esp{la patria} : la patrie (mot féminin, très utilisé dans les discours officiels : \esp{amar la patria}, aimer la patrie).
\item \esp{la bandera} : le drapeau. \esp{La bandera de Camerún es verde, roja y amarilla.} « Le drapeau du Cameroun est vert, rouge et jaune. »
\item \esp{el himno} : l'hymne. \esp{Cantamos el himno nacional.} « Nous chantons l'hymne national. »
\item \esp{el escudo} : le blason, les armoiries.
\item \esp{el desfile} : le défilé. \esp{Hay un desfile militar el 20 de mayo.} « Il y a un défilé militaire le 20 mai. »
\item \esp{la independencia} : l'indépendance.
\item \esp{el ciudadano / la ciudadana} : le citoyen / la citoyenne.
\end{itemize}

\textbf{Champ lexical 2 — les traditions et les fêtes}

\begin{itemize}
\item \esp{la tradición} : la tradition ; \esp{la costumbre} : la coutume (plus proche de l'habitude quotidienne).
\item \esp{la fiesta} : la fête ; \esp{la fiesta nacional} : la fête nationale.
\item \esp{la danza} : la danse ; \esp{la comida típica} : le plat typique.
\item Verbes : \esp{celebrar} (célébrer), \esp{desfilar} (défiler), \esp{cantar} (chanter), \esp{bailar} (danser).
\item \esp{El 20 de mayo es la fiesta nacional de Camerún.} « Le 20 mai est la fête nationale du Cameroun. »
\item \esp{En mi región celebramos la fiesta del ñame.} « Dans ma région, nous célébrons la fête de l'igname. »
\end{itemize}

\textbf{Champ lexical 3 — la diversité}

\begin{itemize}
\item \esp{la diversidad} : la diversité ; \esp{la riqueza} : la richesse.
\item \esp{la lengua / el idioma} : la langue ; \esp{la lengua materna} : la langue maternelle.
\item \esp{el patrimonio} : le patrimoine ; \esp{el mestizaje} : le métissage.
\item Adjectifs : \esp{bilingüe} (bilingue), \esp{multilingüe} (multilingue).
\item \esp{Camerún es un país bilingüe.} « Le Cameroun est un pays bilingue. »
\item \esp{Guinea Ecuatorial es un ejemplo de mestizaje.} « La Guinée équatoriale est un exemple de métissage. »
\end{itemize}

\textbf{Champ lexical 4 — les sentiments et les valeurs}

\begin{itemize}
\item \esp{el orgullo} : la fierté ; \esp{estar orgulloso/a de} : être fier/fière de. \esp{Estoy orgullosa de mi cultura.} « Je suis fière de ma culture. »
\item \esp{el respeto} : le respect ; \esp{la unidad} : l'unité ; \esp{la paz} : la paix.
\item Verbes : \esp{amar} (aimer), \esp{proteger} (protéger), \esp{valorar} (valoriser, apprécier à sa juste valeur).
\item \esp{Debemos proteger nuestro patrimonio.} « Nous devons protéger notre patrimoine. »
\end{itemize}

\begin{attention}[Genre et faux amis]
\begin{itemize}
\item \esp{el himno} et \esp{el patrimonio} sont \textbf{masculins} (terminaison en \emph{-o}) : on dit \esp{el himno nacional}, \esp{el patrimonio cultural}. Ne dites jamais \esp{la himno} ni \esp{la patrimonio}.
\item \esp{la lengua} et \esp{el idioma} sont synonymes, mais dans l'usage courant au Bac on privilégie \esp{la lengua española} et non \esp{el idioma español}. Attention aussi : \esp{la lengua} signifie aussi « la langue » (organe) — le contexte décide.
\item \esp{estar orgulloso de} se construit avec la préposition \textbf{de} : \esp{Estoy orgulloso de mi país.} Ne calquez pas le français « fier à ».
\item \esp{la paz} est féminin : \esp{la paz}, pas \esp{el paz}.
\item \esp{bilingüe} et \esp{multilingüe} sont invariables en genre (terminaison en \emph{-e}) : \esp{un país bilingüe}, \esp{una persona bilingüe}.
\end{itemize}
\end{attention}

\courssub{Tableau récapitulatif espagnol-français}

\begin{center}
\begin{tabularx}{\linewidth}{|X|X|X|}\hline
\rowcolor{popblueL} Español & Français & Remarque \\ \hline
la patria & la patrie & féminin ; \esp{amar la patria} \\ \hline
la bandera & le drapeau & féminin \\ \hline
el himno & l'hymne & \textbf{masculin} : \esp{el himno nacional} \\ \hline
el desfile & le défilé & verbe : \esp{desfilar} \\ \hline
la fiesta & la fête & \esp{la fiesta nacional} \\ \hline
la lengua / el idioma & la langue & dire \esp{la lengua española} \\ \hline
el patrimonio & le patrimoine & \textbf{masculin} \\ \hline
la diversidad & la diversité & \esp{la diversidad cultural} \\ \hline
el mestizaje & le métissage & culturel ou humain \\ \hline
estar orgulloso/a de & être fier de & accord : \esp{orgulloso} / \esp{orgullosa} \\ \hline
\end{tabularx}
\end{center}

\courssub{Carte mentale du lexique}

\begin{popfigure}
\begin{tikzpicture}[mindmap, grow cyclic, every node/.style={concept, text=white, align=center, font=\small}, concept color=popblue, level 1/.append style={level distance=3.4cm, sibling angle=90}, level 2/.append style={level distance=2.1cm, sibling angle=45}]
\node[concept color=popdark, font=\large\bfseries] {Mi patria}
  child[concept color=popblue] { node {Símbolos}
    child { node[concept color=popblue!70] {la bandera} }
    child { node[concept color=popblue!70] {el himno} }
    child { node[concept color=popblue!70] {el escudo} }
  }
  child[concept color=poporange] { node {Fiestas}
    child { node[concept color=poporange!70] {la tradición} }
    child { node[concept color=poporange!70] {la danza} }
    child { node[concept color=poporange!70] {celebrar} }
  }
  child[concept color=popgreen] { node {Diversidad}
    child { node[concept color=popgreen!70] {la lengua} }
    child { node[concept color=popgreen!70] {el mestizaje} }
    child { node[concept color=popgreen!70] {bilingüe} }
  }
  child[concept color=poppurple] { node {Valores}
    child { node[concept color=poppurple!70] {el orgullo} }
    child { node[concept color=poppurple!70] {el respeto} }
    child { node[concept color=poppurple!70] {la paz} }
  };
\end{tikzpicture}
\legende{Carte mentale du lexique : quatre champs autour du thème « Mi patria ».}
\end{popfigure}

\courssub{Expressions utiles pour parler de son pays}

Ces phrases toutes faites vous serviront dans l'expression écrite et orale, notamment pour présenter une tradition ou une fête :

\begin{exemplebox}[Frases para hablar de mi país]
\begin{itemize}
\item \esp{Mi país es Camerún, un país bilingüe y multilingüe.} « Mon pays est le Cameroun, un pays bilingue et multilingue. »
\item \esp{La fiesta nacional se celebra el 20 de mayo.} « La fête nationale se célèbre le 20 mai. » (notez la voix passive réfléchie \esp{se celebra}, très fréquente pour les fêtes)
\item \esp{Estoy orgulloso/a de mi cultura y de mis tradiciones.} « Je suis fier/fière de ma culture et de mes traditions. »
\item \esp{La diversidad es una riqueza, no un problema.} « La diversité est une richesse, pas un problème. »
\item \esp{En mi región, la gente baila danzas tradicionales y come comida típica.} « Dans ma région, les gens dansent des danses traditionnelles et mangent des plats typiques. »
\end{itemize}
\end{exemplebox}

\courssub{Exercice résolu et activités de mémorisation}

\begin{exoresolu}[Clasificar y completar]
\textbf{Énoncé.} a) Classez les mots suivants dans le bon champ lexical (símbolos / fiestas / diversidad / valores) : \esp{el himno, la danza, el mestizaje, la paz, la bandera, el patrimonio, el orgullo, celebrar}. b) Complétez : \esp{Camerún es un país \ldots (bilingüe) ; estoy \ldots (orgulloso) de mi \ldots (patria).}
\tcblower
\textbf{Solution.} a) Símbolos : \esp{el himno, la bandera}. Fiestas : \esp{la danza, celebrar}. Diversidad : \esp{el mestizaje, el patrimonio}. Valores : \esp{la paz, el orgullo}. b) \esp{Camerún es un país bilingüe ; estoy orgulloso (ou orgullosa si c'est une fille qui parle) de mi patria.} Remarque : l'adjectif \esp{bilingüe} est invariable en genre, mais \esp{orgulloso} s'accorde avec le sujet.
\end{exoresolu}

\begin{experience}[Activités de mémorisation en classe]
\begin{enumerate}
\item \textbf{Association mot-image :} le professeur montre une image (drapeau camerounais, défilé du 20 mai, danseurs traditionnels, carte des langues) et les élèves crient le mot espagnol correspondant : \esp{¡la bandera! ¡el desfile! ¡la danza!}
\item \textbf{Classement par champ :} en équipes, distribuez vingt cartes de vocabulaire et classez-les le plus vite possible dans les quatre branches de la carte mentale.
\item \textbf{Mini-dictée :} le professeur dicte huit mots au choix (\esp{la patria, el himno, la tradición, la diversidad, el mestizaje, el orgullo, la paz, bilingüe}) ; chaque élève les écrit avec l'article correct, puis rédige une phrase avec deux d'entre eux.
\end{enumerate}
\end{experience}

\begin{exercice}[À vous de jouer]
1. Traduisez en espagnol : « Le drapeau du Cameroun est vert, rouge et jaune. » / « Je suis fier de mes traditions. » / « La diversité est une richesse. » 2. Rédigez quatre phrases pour présenter une fête de votre village ou de votre région en utilisant : \esp{se celebra, la gente, danzas, comida típica}.
\end{exercice}

\begin{corrige}[Exercice « À vous de jouer »]
1. \esp{La bandera de Camerún es verde, roja y amarilla.} / \esp{Estoy orgulloso (orgullosa) de mis tradiciones.} / \esp{La diversidad es una riqueza.} 2. Modèle : \esp{En mi pueblo se celebra la fiesta del ñame en diciembre. La gente baila danzas tradicionales y canta en la lengua local. Preparamos comida típica como el ndolé. Estoy orgulloso de esta tradición porque une a las familias.}
\end{corrige}

\begin{savaistu}[Un pays, 250 langues]
Le Cameroun compte plus de 250 langues locales : c'est l'un des pays les plus diversifiés linguistiquement au monde. En espagnol, on peut dire : \esp{Camerún es uno de los países con más diversidad lingüística del mundo.} C'est un argument parfait pour un essai ou un commentaire sur \esp{la diversidad} : la diversité linguistique n'est pas un obstacle, c'est une richesse nationale. Et n'oubliez pas la Guinée équatoriale, voisine du Cameroun : c'est le seul pays d'Afrique où l'espagnol est langue officielle, un beau symbole de \esp{mestizaje} entre l'Afrique et le monde hispanique.
\end{savaistu}

\begin{aretenir}[L'essentiel du lexique]
Retenez les quatre champs : \textbf{símbolos} (\esp{patria, bandera, himno, escudo}), \textbf{fiestas} (\esp{tradición, fiesta, danza, celebrar}), \textbf{diversidad} (\esp{lengua, patrimonio, mestizaje, bilingüe}) et \textbf{valores} (\esp{orgullo, respeto, unidad, paz}). Surveillez le genre de \esp{el himno} et \esp{el patrimonio}, dites \esp{la lengua española}, et employez \esp{estar orgulloso/a de} avec la préposition \esp{de}. Avec ces trente mots, vous pouvez déjà présenter votre pays et défendre la diversité culturelle en espagnol.
\end{aretenir}

\coursec{Grammaire outillée : ser/estar et le présent des traditions}

Dans la section précédente, nous avons découvert le lexique de la patrie et des symboles : \esp{patria}, \esp{bandera}, \esp{himno}, \esp{tradición}, \esp{patrimonio}, \esp{diversidad}. Mais un vocabulaire ne vit que dans des phrases. Or, dès que l'on veut décrire son pays, situer une ville ou dire sa fierté, l'espagnol nous impose un choix que le français ignore : celui entre \esp{ser} et \esp{estar}, deux verbes qui correspondent tous les deux à notre verbe « être ». C'est l'une des distinctions les plus importantes de la grammaire espagnole — et l'une des plus rentables à maîtriser pour l'examen. Observons, réglons, entraînons-nous.

\courssub{Observation : deux verbes pour un seul « être »}

Reprenons deux phrases inspirées de notre texte support \emph{« Mi patria, mis culturas »} :

\esp{Camerún es un país bilingüe.} « Le Cameroun est un pays bilingue. »

\esp{La bandera está delante del colegio.} « Le drapeau est (se trouve) devant le collège. »

Dans les deux cas, le français dit « est ». L'espagnol, lui, choisit \esp{es} (verbe \esp{ser}) dans la première phrase et \esp{está} (verbe \esp{estar}) dans la seconde. Pourquoi ? Parce que dans la première, on dit \textbf{ce qu'est} le Cameroun (son identité, sa nature permanente), tandis que dans la seconde, on dit \textbf{où se trouve} le drapeau (sa localisation, qui peut changer demain).

\begin{propriete}[La règle fondamentale : ser = esencia, estar = estado y lugar]
\begin{itemize}
\item \cle{Ser} exprime l'\textbf{essence} : l'identité (qui on est), l'origine (d'où l'on vient), la nationalité, la profession, la caractéristique permanente, la matière, l'heure, ainsi que la date et le lieu des événements.
\item \cle{Estar} exprime l'\textbf{état et le lieu} : la localisation (où l'on se trouve) et l'état temporaire (comment on se sent à un moment donné).
\end{itemize}
Astuce de mémorisation : \esp{ser} répond à la question \esp{¿qué es?} (qu'est-ce que c'est ?) ; \esp{estar} répond aux questions \esp{¿dónde está?} (où est-ce ?) et \esp{¿cómo está?} (comment ça va ?).
\end{propriete}

\begin{popfigure}
\begin{tikzpicture}[grow=right, level distance=3.4cm, edge from parent/.style={draw=popmuted, thick}, every node/.style={align=center}]
\node[fill=popblue, text=white, rounded corners, font=\bfseries] {SER / ESTAR}
child { node[fill=poporange, text=white, rounded corners, font=\bfseries, align=center]{ESTAR\\estado y lugar}
  child { node[fill=poporangeL, rounded corners, font=\small, align=center]{sentimiento\\del momento\\\esp{estoy orgulloso}} }
  child { node[fill=poporangeL, rounded corners, font=\small, align=center]{estado\\temporal\\\esp{está cansado}} }
  child { node[fill=poporangeL, rounded corners, font=\small, align=center]{lugar\\\esp{está en África}} }
}
child { node[fill=popgreen, text=white, rounded corners, font=\bfseries, align=center]{SER\\esencia}
  child { node[fill=popgreenL, rounded corners, font=\small, align=center]{fecha de\\eventos\\\esp{la fiesta es en mayo}} }
  child { node[fill=popgreenL, rounded corners, font=\small, align=center]{característica\\\esp{es grande}} }
  child { node[fill=popgreenL, rounded corners, font=\small, align=center]{origen\\\esp{es de Bafoussam}} }
  child { node[fill=popgreenL, rounded corners, font=\small, align=center]{identidad\\\esp{es un país}} }
};
\end{tikzpicture}
\legende{Carte mentale : les emplois de \esp{ser} (essence) et \esp{estar} (état et lieu).}
\end{popfigure}

\courssub{Conjugaison : ser et estar au présent de l'indicatif}

Ces deux verbes sont irréguliers ; leurs formes doivent être sues par cœur.

\begin{center}
\begin{tabular}{|l|l|l|l|}
\hline
\rowcolor{popblueL} \textbf{Personne} & \textbf{ser} & \textbf{estar} & \textbf{Français} \\
\hline
yo & soy & estoy & je suis \\
\hline
tú & eres & estás & tu es \\
\hline
él / ella / usted & es & está & il / elle est \\
\hline
nosotros/as & somos & estamos & nous sommes \\
\hline
vosotros/as & sois & estáis & vous êtes \\
\hline
ellos / ellas / ustedes & son & están & ils / elles sont \\
\hline
\end{tabular}
\end{center}

\begin{attention}[Accents et confusions]
Notez les accents écrits de \esp{estar} : \esp{estás}, \esp{está}, \esp{están} (accent sur la dernière syllabe) et \esp{estáis} (accent sur le \emph{a}). Ne confondez pas \esp{está} (il est, verbe \esp{estar}) et \esp{esta} (cette, adjectif démonstratif) : \esp{Esta bandera está limpia.} « Ce drapeau est propre. » De même, \esp{están} (ils sont) ne s'écrit pas comme \esp{estan} sans accent.
\end{attention}

\courssub{Ser : l'identité, l'origine, la caractéristique permanente}

On emploie \cle{ser} suivi d'un \textbf{nom} ou d'un \textbf{adjectif} de nationalité ou de caractéristique durable. Pour l'origine, la construction est \esp{ser de} + lieu :

\begin{exemplebox}[Ser = esencia]
\esp{México es un país grande.} « Le Mexique est un grand pays. » (identité)

\esp{Los cameruneses son acogedores.} « Les Camerounais sont accueillants. » (caractéristique permanente)

\esp{Mi profesora es de Douala.} « Ma professeure est de Douala. » (origine : \esp{ser de} + lieu)

\esp{El himno nacional es un símbolo de la patria.} « L'hymne national est un symbole de la patrie. » (définition)

\esp{Soy camerunés y estoy orgulloso de serlo.} « Je suis camerounais et je suis fier de l'être. » (nationalité ; \esp{serlo} = « l'être », le pronom \esp{lo} reprend l'adjectif)
\end{exemplebox}

\courssub{Estar : la localisation et l'état temporaire}

On emploie \cle{estar} suivi d'un \textbf{complément de lieu} (souvent introduit par \esp{en}, \esp{delante de}, \esp{cerca de}, \esp{lejos de}) ou d'un \textbf{adjectif d'état} :

\begin{exemplebox}[Estar = lugar y estado]
\esp{Guinea Ecuatorial está en África Central.} « La Guinée équatoriale se trouve en Afrique centrale. » (localisation)

\esp{España está en Europa.} « L'Espagne se trouve en Europe. » (localisation)

\esp{Estoy orgulloso de mi patria.} « Je suis fier de ma patrie. » (sentiment du moment)

\esp{Los alumnos están contentos antes de la fiesta nacional.} « Les élèves sont contents avant la fête nationale. » (état temporaire)

\esp{¿Dónde está la bandera?} « Où est le drapeau ? »
\end{exemplebox}

\begin{exemplebox}[Le même pays, les deux verbes]
\esp{España es un país europeo.} « L'Espagne est un pays européen. » (\esp{ser} : identité, nature)

\esp{España está en Europa.} « L'Espagne se trouve en Europe. » (\esp{estar} : localisation)

\esp{Estamos orgullosos de nuestras tradiciones.} « Nous sommes fiers de nos traditions. » (\esp{estar} : sentiment)
\end{exemplebox}

\begin{attention}[Le piège n° 1 des francophones]
Le français n'a qu'un seul verbe, « être » : la tentation est grande de traduire systématiquement par \esp{ser}. C'est faux ! Pour \textbf{situer une ville, un pays, un bâtiment}, l'espagnol exige \esp{estar} :

\esp{Yaoundé está en el centro de Camerún.} « Yaoundé est au centre du Cameroun. » (et non \esp{Yaoundé es en el centro})

\esp{El mercado está cerca del colegio.} « Le marché est près du collège. »
\end{attention}

\begin{attention}[Deux pièges subtils]
\textbf{1.} \esp{ser orgulloso} désigne un trait de caractère (quelqu'un d'orgueilleux, de fier par nature) ; \esp{estar orgulloso de} exprime le sentiment éprouvé à un moment précis : \esp{Estoy orgulloso de mi país.} « Je suis fier de mon pays » (en ce moment, à juste titre).

\textbf{2.} Pour situer un \textbf{événement} dans le temps ou dans l'espace, l'espagnol utilise \esp{ser} et non \esp{estar} : \esp{La fiesta es en mayo.} « La fête a lieu en mai. » ; \esp{El desfile es en la avenida principal.} « Le défilé a lieu sur l'avenue principale. » On ne dit pas \esp{la fiesta está en mayo}. Logique : un événement n'a pas de position fixe comme un bâtiment ; sa date et son lieu font partie de sa définition.
\end{attention}

\begin{center}
\begin{tabularx}{\linewidth}{|X|X|X|}
\hline
\rowcolor{popblueL} \textbf{Français} & \textbf{Español} & \textbf{Remarque} \\
\hline
Le Cameroun est un pays bilingue. & Camerún es un país bilingüe. & ser : identité \\
\hline
Le Cameroun se trouve en Afrique centrale. & Camerún está en África Central. & estar : localisation \\
\hline
Nous sommes fiers de notre pays. & Estamos orgullosos de nuestro país. & estar : sentiment \\
\hline
Les Camerounais sont hospitaliers. & Los cameruneses son hospitalarios. & ser : caractère permanent \\
\hline
La fête est (a lieu) le 20 mai. & La fiesta es el 20 de mayo. & ser : date d'un événement \\
\hline
\end{tabularx}
\end{center}

\courssub{Le présent de l'indicatif : le temps des traditions}

Pour parler des \textbf{habitudes}, des \textbf{traditions} et des \textbf{vérités culturelles générales}, l'espagnol utilise le présent de l'indicatif, exactement comme le français. C'est le temps naturel du commentaire sur les fêtes et les coutumes. Les marqueurs fréquents de ce temps sont : \esp{cada año} (chaque année), \esp{siempre} (toujours), \esp{normalmente} (normalement), \esp{todos los años} (tous les ans).

\begin{exemplebox}[Le présent des traditions]
\esp{Cada año, los alumnos desfilan el 20 de mayo.} « Chaque année, les élèves défilent le 20 mai. »

\esp{En México, la gente come tacos y baila.} « Au Mexique, les gens mangent des tacos et dansent. »

\esp{En mi región, la gente baila el bikutsi en las fiestas.} « Dans ma région, les gens dansent le bikutsi pendant les fêtes. »

\esp{Los españoles celebran muchas fiestas tradicionales.} « Les Espagnols célèbrent de nombreuses fêtes traditionnelles. »

\esp{En Guinea Ecuatorial, se habla español.} « En Guinée équatoriale, on parle espagnol. » (tournure impersonnelle avec \esp{se})
\end{exemplebox}

Rappel de conjugaison au présent (verbes réguliers) : \esp{celebrar} → celebro, celebras, celebra, celebramos, celebráis, celebran ; \esp{comer} → como, comes, come, comemos, coméis, comen ; \esp{vivir} → vivo, vives, vive, vivimos, vivís, viven. Attention aux verbes à changement de voyelle déjà rencontrés : \esp{jugar} → \esp{juega}, \esp{poder} → \esp{puede}. Notez aussi que \esp{la gente} (les gens) est un nom \textbf{singulier} en espagnol : \esp{la gente baila}, et non \esp{bailan}.

\courssub{Texte d'application : le présent et ser/estar en contexte}

\begin{textebox}[« Una fiesta de diversidad »]
En mi colegio de Yaoundé, cada año celebramos la Fiesta de la Diversidad Cultural. Es un día muy especial: los alumnos llevan ropa tradicional de todas las regiones de Camerún. Camerún es un país con más de doscientas lenguas, y estamos muy orgullosos de esta riqueza.

Durante la fiesta, la bandera camerunesa está delante de la entrada principal. Es verde, roja y amarilla, con una estrella en el centro. Los profesores explican el significado de los colores: el verde es el símbolo de los bosques del sur, el rojo representa la unidad y el amarillo es el sol del norte.

En el patio, la gente canta y baila. Los alumnos del oeste bailan danzas bamiléké; los del litoral presentan el makossa. Mi amiga Lucía, que es de Guinea Ecuatorial, lleva un vestido precioso y habla de las tradiciones de su país. Guinea Ecuatorial está en África Central y es el único país africano donde el español es lengua oficial.

Al final del día, todos comemos platos típicos: ndolé, eru, poulet DG. Estamos cansados pero muy contentos. Esta fiesta es importante porque nos enseña que la diversidad no es un problema: es un tesoro. Como dice nuestro director: «Somos diferentes, pero somos una sola patria».
\end{textebox}

\textbf{Glossaire :} \esp{llevar} = porter (un vêtement) ; \esp{la riqueza} = la richesse ; \esp{la entrada} = l'entrée ; \esp{el bosque} = la forêt ; \esp{el patio} = la cour ; \esp{precioso} = magnifique ; \esp{el plato típico} = le plat typique ; \esp{el tesoro} = le trésor ; \esp{enseñar} = enseigner, montrer ; \esp{ndolé}, \esp{eru}, \esp{poulet DG} = plats typiques du Cameroun.

\textbf{Questions de grammaire :}
\begin{enumerate}
\item Relevez trois emplois de \esp{ser} et trois emplois de \esp{estar} dans le texte et justifiez chacun.
\item Pourquoi dit-on \esp{la bandera está delante de la entrada} et non \esp{es delante} ?
\item Pourquoi dit-on \esp{esta fiesta es importante} et non \esp{está importante} ?
\end{enumerate}

\courssub{Exercice résolu et entraînement}

\begin{exoresolu}[Ser ou estar ?]
Complétez avec la forme correcte de \esp{ser} ou \esp{estar} au présent :
\begin{enumerate}
\item Camerún \_\_\_ un país de África Central.
\item La bandera \_\_\_ delante del colegio.
\item Nosotros \_\_\_ orgullosos de nuestras tradiciones.
\item La fiesta nacional \_\_\_ el 20 de mayo.
\item Mis abuelos \_\_\_ de Bafoussam.
\item Hoy los alumnos \_\_\_ muy contentos.
\end{enumerate}
\tcblower
\textbf{Solution détaillée :}
\begin{enumerate}
\item \esp{es} — identité, définition du pays (ser).
\item \esp{está} — localisation du drapeau (estar + lieu).
\item \esp{estamos} — sentiment du moment (estar + état).
\item \esp{es} — date d'un événement : les événements se situent avec ser.
\item \esp{son} — origine : \esp{ser de} + lieu.
\item \esp{están} — état temporaire (« aujourd'hui » indique le moment).
\end{enumerate}
\end{exoresolu}

\begin{exercice}[À vous de jouer]
Traduisez en espagnol :
\begin{enumerate}
\item Le Mexique est un grand pays d'Amérique.
\item La Guinée équatoriale se trouve en Afrique centrale.
\item Je suis fier de ma culture.
\item Chaque année, nous célébrons la fête nationale le 20 mai.
\item Les Camerounais sont accueillants et généreux.
\end{enumerate}
\end{exercice}

\begin{corrige}[Exercice : traduction]
1. \esp{México es un país grande de América.} 2. \esp{Guinea Ecuatorial está en África Central.} 3. \esp{Estoy orgulloso de mi cultura.} 4. \esp{Cada año, celebramos la fiesta nacional el 20 de mayo.} 5. \esp{Los cameruneses son acogedores y generosos.}
\end{corrige}

\begin{aretenir}[L'essentiel de la section]
\begin{itemize}
\item \esp{Ser} = essence : identité, origine (\esp{ser de}), nationalité, profession, caractéristique permanente, date et lieu d'un événement.
\item \esp{Estar} = état et lieu : localisation, sentiment, état temporaire.
\item Le présent de l'indicatif exprime les traditions, les habitudes et les vérités culturelles.
\item Une seule question à se poser : « est-ce que je dis ce que c'est (\esp{ser}) ou où/comment c'est (\esp{estar}) ? »
\end{itemize}
\end{aretenir}

\begin{savaistu}[Ser/estar au baccalauréat]
Au Bac, la distinction \esp{ser}/\esp{estar} est un point de grammaire presque systématiquement évalué, en version comme en expression écrite. Les correcteurs repèrent immédiatement la faute \esp{Camerún es en África} au lieu de \esp{Camerún está en África}. Maîtriser cette distinction, c'est sécuriser des points faciles dans toutes les épreuves de la série A. Et petit bonus culturel : la Guinée équatoriale, voisine du Cameroun, est le seul pays d'Afrique où l'espagnol est langue officielle — une raison de plus de bien conjuguer \esp{ser} et \esp{estar} !
\end{savaistu}

\coursec{Méthode du commentaire de texto}

Dans la section précédente, nous avons appris à utiliser \esp{ser} et \esp{estar} pour décrire une tradition, et le présent de l'indicatif pour raconter les fêtes de notre pays. Nous savons donc maintenant \emph{parler} de la diversité culturelle. Mais à l'examen comme dans la vie scolaire, il ne suffit pas de parler : il faut aussi savoir \emph{lire} un texte et en rendre compte de manière organisée. C'est exactement le rôle du \cle{comentario de texto}, l'exercice roi de la classe de Première. Cette section vous donne une méthode simple, en trois étapes, pour réussir à tous les coups.

\begin{definition}[Lexique clé de la leçon : patria, símbolos y tradiciones]
Voici les mots essentiels du module, à connaître pour lire, commenter et produire :
\begin{itemize}
\item \esp{la patria} (la patrie, le pays natal) — \esp{la bandera} (le drapeau) — \esp{el himno} (l'hymne).
\item \esp{la tradición} (la tradition) — \esp{la fiesta} (la fête) — \esp{el patrimonio} (le patrimoine).
\item \esp{la lengua} (la langue) — \esp{la diversidad} (la diversité) — \esp{el mestizaje} (le métissage).
\end{itemize}
\emph{Attention :} \esp{patria} est féminin ; \esp{himno}, \esp{patrimonio}, \esp{mestizaje} sont masculins.
\end{definition}

\courssub{Qu'est-ce qu'un commentaire de texto ?}

\begin{definition}[El comentario de texto]
Le \cle{commentaire de texto} (\esp{comentario de texto}) est une \textbf{lecture organisée} d'un texte. Il fait trois choses, dans cet ordre : il \textbf{présente} le texte (de quoi s'agit-il ?), il \textbf{analyse} ses idées (que dit le texte, et comment ?), et il \textbf{apprécie} (quel est mon avis personnel ?). En espagnol : \esp{presentar, analizar y valorar un texto}.
\end{definition}

Observez d'abord la différence entre trois exercices que les élèves confondent souvent :

\begin{center}
\begin{tabularx}{\linewidth}{|X|X|X|}
\hline
\rowcolor{popblueL} \textbf{Exercice} & \textbf{Ce qu'on fait} & \textbf{Ce qu'on ne fait pas} \\
\hline
\esp{Resumen} (résumé) & Reprendre les idées du texte avec d'autres mots, plus court & Donner son avis \\
\hline
\esp{Traducción} (traduction) & Passer le texte d'une langue à l'autre & Analyser les idées \\
\hline
\esp{Comentario} (commentaire) & Présenter, analyser avec des citations, donner son point de vue & Recopier le texte \\
\hline
\end{tabularx}
\end{center}

\begin{attention}[Un commentaire n'est ni un résumé ni une traduction]
L'erreur la plus fréquente des élèves francophones est de \textbf{raconter le texte} phrase par phrase, ou de le traduire. Un commentaire doit \textbf{analyser les idées} (les expliquer, les illustrer par une citation) et se terminer par un \textbf{point de vue personnel}. Si votre copie ne contient aucune phrase commençant par \esp{En mi opinión...} ou \esp{Para mí...}, ce n'est pas un commentaire !
\end{attention}

\courssub{Les trois étapes du commentaire}

\begin{propriete}[La structure en trois parties]
Tout commentaire de texto suit un plan fixe, comme un voyage en trois étapes :
\begin{enumerate}
\item \textbf{Introducción} : présenter le texte (titre, auteur ou narrateur, thème) et annoncer l'\textbf{idée principale} (\esp{la idea principal}).
\item \textbf{Desarrollo} : dégager \textbf{deux ou trois idées secondaires} (\esp{ideas secundarias}), chacune appuyée par une \textbf{citation du texte} entre guillemets, traduite ou expliquée.
\item \textbf{Conclusión} : faire le \textbf{bilan} (\esp{balance}) et ouvrir : comparaison avec son propre pays, avis personnel.
\end{enumerate}
\end{propriete}

\begin{popfigure}
\begin{tikzpicture}[
node distance=0.9cm,
etape/.style={rectangle, rounded corners, draw=popblue, fill=popblueL, text width=9.5cm, align=center, font=\small, inner sep=6pt},
fleche/.style={-{Stealth[length=3mm]}, very thick, poporange}]
\node[etape, align=center] (intro){\textbf{1. INTRODUCCIÓN}\\ Presentación del texto (título, autor, tema) + \emph{idea principal}};
\node[etape, below=of intro, draw=poppurple, fill=poppurpleL, align=center] (des){\textbf{2. DESARROLLO}\\ Idea secundaria 1 + cita del texto \quad / \quad Idea secundaria 2 + cita del texto};
\node[etape, below=of des, draw=popgreen, fill=popgreenL, align=center] (conc){\textbf{3. CONCLUSIÓN}\\ Balance + opinión personal (\emph{En mi opinión...})};
\draw[fleche] (intro) -- (des);
\draw[fleche] (des) -- (conc);
\end{tikzpicture}
\legende{Organigramme du commentaire de texto : trois étapes, toujours dans le même ordre.}
\end{popfigure}

\textbf{Étape 1 — L'introduction.}\par
L'introduction répond à trois questions : \emph{Quel texte ? De qui ? De quoi ?} On présente le titre, l'auteur ou le narrateur, le thème, puis on annonce l'idée principale en une phrase.

\begin{exemplebox}[Modèles d'introduction]
\esp{El texto trata del patriotismo y de la diversidad cultural.} = Le texte traite du patriotisme et de la diversité culturelle.\par
\esp{El narrador presenta las fiestas y las tradiciones de su país.} = Le narrateur présente les fêtes et les traditions de son pays.\par
\esp{La idea principal del texto es que la diversidad cultural es una riqueza.} = L'idée principale du texte est que la diversité culturelle est une richesse.
\end{exemplebox}

\textbf{Étape 2 — Le développement.}\par
C'est le cœur du commentaire. On relit le texte et on dégage deux ou trois idées secondaires. Pour chaque idée : une phrase d'annonce, une \textbf{citation exacte} du texte entre guillemets (\esp{« ... »}), puis une explication ou une traduction de la citation. Une idée sans citation n'est pas recevable : la citation est votre preuve.

\begin{exemplebox}[Modèle d'idée développée]
\esp{Primero, el texto muestra el amor por la patria. El narrador dice que «mi patria es mi familia y mi tierra». Esta frase significa que la patria no es solo un país, sino también las personas que queremos.} = D'abord, le texte montre l'amour de la patrie. Le narrateur dit que « ma patrie, c'est ma famille et ma terre ». Cette phrase signifie que la patrie n'est pas seulement un pays, mais aussi les personnes que nous aimons.
\end{exemplebox}

\textbf{Étape 3 — La conclusion.}\par
La conclusion fait le bilan en une ou deux phrases, puis \textbf{ouvre} : on compare avec le Cameroun, on donne son avis. C'est ici que le correcteur voit votre personnalité.

\begin{exemplebox}[Modèle de conclusion]
\esp{En conclusión, el texto defiende la diversidad cultural. En mi opinión, la diversidad cultural es una riqueza para Camerún, porque tenemos muchas lenguas y tradiciones. Para mí, ser patriota es respetar todas las culturas de mi país.} = En conclusion, le texte défend la diversité culturelle. À mon avis, la diversité culturelle est une richesse pour le Cameroun, car nous avons beaucoup de langues et de traditions. Pour moi, être patriote, c'est respecter toutes les cultures de mon pays.
\end{exemplebox}

\courssub{Les connecteurs du commentaire}

Un bon commentaire est un texte \emph{lié} : les phrases ne sont pas jetées les unes après les autres, elles sont enchaînées par des \cle{conectores}.

\begin{center}
\begin{tabularx}{\linewidth}{|X|X|X|}
\hline
\rowcolor{popblueL} \textbf{Conector} & \textbf{Français} & \textbf{Emploi dans le commentaire} \\
\hline
\esp{primero} & d'abord, premièrement & annoncer la première idée \\
\hline
\esp{además} & de plus, en outre & ajouter une deuxième idée \\
\hline
\esp{sin embargo} & cependant, pourtant & nuancer, opposer deux idées \\
\hline
\esp{por eso} & c'est pourquoi & tirer une conséquence \\
\hline
\esp{en conclusión} & en conclusion & ouvrir la conclusion \\
\hline
\esp{en mi opinión} & à mon avis & donner son point de vue personnel \\
\hline
\end{tabularx}
\end{center}

\begin{exemplebox}[Connecteurs en contexte]
\esp{Además, el narrador habla del mestizaje.} = De plus, le narrateur parle du métissage.\par
\esp{Sin embargo, la diversidad no es un problema.} = Cependant, la diversité n'est pas un problème.\par
\esp{Por eso, debemos respetar las tradiciones de los demás.} = C'est pourquoi nous devons respecter les traditions des autres.
\end{exemplebox}

\begin{attention}[Pièges des francophones]
\begin{itemize}
\item \esp{Sin embargo} ne signifie jamais « sans embargo » : c'est « cependant ». Ne le confondez pas avec \esp{sin} (sans).
\item \esp{En mi opinión} : ne dites pas \esp{en mi opinion} sans accent, ni \esp{según mí} calqué du français « selon moi » ; la forme correcte est \esp{en mi opinión} ou \esp{para mí}.
\item Les citations se mettent entre guillemets et se recopient \textbf{mot à mot}, avec les accents du texte.
\end{itemize}
\end{attention}

\courssub{Un commentaire modèle à étudier}

Lisez ce commentaire complet du texte « Mi patria, mis culturas » étudié dans les sections précédentes. Repérez les trois parties et les connecteurs.

\begin{textebox}[Comentario modelo : « Mi patria, mis culturas »]
El texto « Mi patria, mis culturas » trata del patriotismo y de la diversidad cultural. El narrador presenta las tradiciones de su país. La idea principal es que la diversidad cultural es una riqueza y no un problema.\par
Primero, el texto muestra el amor por la patria. El narrador dice que « mi patria es mi familia, mi tierra y mi historia ». Esta frase significa que la patria es más que un territorio: es también la gente y la memoria.\par
Además, el narrador habla de las fiestas y de las lenguas. Según el texto, « cada región tiene su fiesta, su música y su lengua ». Esta cita muestra que la diversidad está en todas partes: en las fiestas, en la música y en las palabras.\par
Sin embargo, el texto explica que esta diversidad no divide al país. El narrador afirma que « somos diferentes, pero estamos unidos ». Por eso, la diversidad cultural es una fuerza.\par
En conclusión, el texto defiende el patriotismo y el respeto de todas las culturas. En mi opinión, la diversidad cultural es muy importante para Camerún, porque nuestro país tiene muchas lenguas y muchas tradiciones. Para mí, ser patriota es amar todas las culturas de mi país.
\end{textebox}

\textbf{Glossaire :} \esp{la tierra} = la terre ; \esp{la gente} = les gens ; \esp{la memoria} = la mémoire ; \esp{unidos} = unis ; \esp{la fuerza} = la force ; \esp{el respeto} = le respect ; \esp{afirmar} = affirmer.

\begin{methode}[Pour réussir votre commentaire en cinq gestes]
\begin{enumerate}
\item \textbf{Lire deux fois le texte} : une fois pour comprendre, une fois pour analyser.
\item \textbf{Souligner les idées} : une couleur pour l'idée principale, une autre pour les idées secondaires.
\item \textbf{Rédiger un plan en trois parties} : introducción, desarrollo (2 ou 3 idées), conclusión.
\item \textbf{Citer le texte entre guillemets} : chaque idée est prouvée par une citation exacte, expliquée ou traduite.
\item \textbf{Donner son avis en conclusion} : \esp{En mi opinión...}, avec une comparaison avec le Cameroun.
\end{enumerate}
\end{methode}

\begin{textebox}[Frases hechas : les phrases toutes faites du commentaire]
\esp{El texto trata de...} (Le texte traite de...)\\
\esp{El narrador dice que...} (Le narrateur dit que...)\\
\esp{Según el texto...} (Selon le texte...)\\
\esp{La idea principal es que...} (L'idée principale est que...)\\
\esp{Esta cita muestra que...} (Cette citation montre que...)\\
\esp{En mi opinión, la diversidad cultural es...} (À mon avis, la diversité culturelle est...)
\end{textebox}

\begin{exoresolu}[Repérage des trois parties]
\textbf{Énoncé.} Dans le commentaire modèle ci-dessus : a) recopiez la phrase qui annonce l'idée principale ; b) recopiez les deux citations du développement ; c) recopiez la phrase qui donne l'avis personnel ; d) soulignez trois connecteurs.
\tcblower
\textbf{Solution.} a) Idée principale : \esp{« La idea principal es que la diversidad cultural es una riqueza y no un problema. »} b) Citations : \esp{« mi patria es mi familia, mi tierra y mi historia »} et \esp{« cada región tiene su fiesta, su música y su lengua »} (on peut ajouter \esp{« somos diferentes, pero estamos unidos »}). c) Avis personnel : \esp{« En mi opinión, la diversidad cultural es muy importante para Camerún... »}. d) Connecteurs : \esp{primero}, \esp{además}, \esp{sin embargo}, \esp{por eso}, \esp{en conclusión} (trois suffisent).
\end{exoresolu}

\begin{exercice}[À vous de commenter]
Choisissez un court texte sur une fête de votre région (ou reprenez le texte « Mi patria, mis culturas »). Rédigez un commentaire de 8 à 10 phrases en suivant le plan : introduction avec \esp{El texto trata de...}, deux idées avec citations, conclusion avec \esp{En mi opinión...}. Utilisez au moins quatre connecteurs de la liste.
\end{exercice}

\begin{savaistu}[Le patriotisme au Cameroun et dans l'espace hispanophone]
Au Cameroun, le patriotisme s'exprime à travers le bilinguisme (français/anglais), les \esp{fiestas} locales (Ngondo, Nyem-Nyem, Medumba), la diversité des \esp{lenguas} nationales (plus de 250) et le \esp{patrimonio} commun (monts, fleuves, forêt). En Espagne, chaque \esp{comunidad autónoma} a son \esp{himno}, son \esp{bandera} et ses \esp{tradiciones} (Castells en Catalogne, Fallas à Valence). Au Mexique, le \esp{mestizaje} indigène-espagnol fonde l'identité nationale (\esp{Día de la Independencia}, \esp{Día de Muertos}). En Guinée équatoriale, seul pays africain hispanophone, le patriotisme unit les cultures Fang, Bubi, Ndowe et Annobonaise sous le drapeau national. \textbf{Point commun :} partout, \esp{la diversidad cultural es una riqueza, no un problema}.
\end{savaistu}

\begin{aretenir}[L'essentiel]
Le commentaire de texto est une lecture organisée en trois parties : \textbf{introducción} (présentation + idée principale), \textbf{desarrollo} (2 ou 3 idées, chacune avec une citation expliquée), \textbf{conclusión} (bilan + avis personnel). On enchaîne avec les connecteurs \esp{primero, además, sin embargo, por eso, en conclusión, en mi opinión}. Un commentaire n'est ni un résumé ni une traduction : il analyse et il prend position.
\end{aretenir}

\coursec{Commentaire modèle et production guidée}

Dans la section précédente, nous avons construit pas à pas la \cle{méthode du commentaire de texte} : lire deux fois, dégager le thème, rédiger une introduction, développer deux idées avec des citations, conclure avec un avis personnel. Il est temps maintenant de voir cette méthode \emph{en action} : nous allons lire un \cle{commentaire modèle} rédigé sur le texte « Mi patria, mis culturas », l'analyser ensemble pour repérer ses rouages (connecteurs, citations, opinion), puis vous produirez vous-mêmes un commentaire guidé sur une fête camerounaise, avant de préparer une présentation orale d'une tradition de votre région. C'est l'aboutissement de toute la leçon.

\courssub{Lecture du commentaire modèle}

Lisons d'abord un commentaire court, rédigé selon le plan imposé : introduction, deux idées développées avec citations, conclusion avec ouverture sur le Cameroun.

\begin{textebox}[Commentaire modèle sur « Mi patria, mis culturas »]
El texto « Mi patria, mis culturas » trata del amor a la patria en Camerún y en los países hispanohablantes. El narrador presenta su país y sus símbolos con mucho orgullo.

Primero, el narrador describe la fiesta del 20 de mayo en Camerún: « los alumnos desfilan con orgullo ». Esta descripción muestra que la juventud camerunesa ama su patria y respeta sus símbolos, como la bandera y el himno nacional.

Además, el texto compara esta fiesta con el « Grito de Dolores » en México, la fiesta de la independencia mexicana. Esta comparación demuestra que el patriotismo existe en todos los países, en África y en América.

En conclusión, este texto es un homenaje a la patria y a la diversidad cultural. En mi opinión, la diversidad cultural es una riqueza para todos los pueblos: en Camerún, nuestras lenguas, nuestras danzas y nuestras fiestas son un tesoro que debemos proteger.
\end{textebox}

\textbf{Traduction du commentaire modèle}

« Le texte "Ma patrie, mes cultures" traite de l'amour de la patrie au Cameroun et dans les pays hispanophones. Le narrateur présente son pays et ses symboles avec beaucoup de fierté.

D'abord, le narrateur décrit la fête du 20 mai au Cameroun : "les élèves défilent avec fierté". Cette description montre que la jeunesse camerounaise aime sa patrie et respecte ses symboles, comme le drapeau et l'hymne national.

De plus, le texte compare cette fête au "Cri de Dolores" au Mexique, la fête de l'indépendance mexicaine. Cette comparaison démontre que le patriotisme existe dans tous les pays, en Afrique et en Amérique.

En conclusion, ce texte est un hommage à la patrie et à la diversité culturelle. À mon avis, la diversité culturelle est une richesse pour tous les peuples : au Cameroun, nos langues, nos danses et nos fêtes sont un trésor que nous devons protéger. »

\courssub{Analyse collective du modèle}

Observons maintenant \emph{comment} ce commentaire est construit. En classe, soulignez dans le texte les éléments suivants avec trois couleurs différentes : les \cle{connecteurs}, les \cle{citations} et l'\cle{avis personnel}.

\begin{center}
\begin{tabularx}{\linewidth}{|X|X|X|}
\hline
\rowcolor{popblueL} Élément & Exemples du modèle & Fonction \\
\hline
Connecteurs logiques & \esp{primero}, \esp{además}, \esp{en conclusión} & Organiser le commentaire en étapes claires \\
\hline
Verbes de présentation & \esp{trata de}, \esp{describe}, \esp{muestra}, \esp{compara}, \esp{demuestra} & Présenter le texte et ses procédés \\
\hline
Citations entre guillemets & « los alumnos desfilan con orgullo » ; « Grito de Dolores » & Prouver chaque idée par le texte \\
\hline
Avis personnel & \esp{en mi opinión}, \esp{debemos proteger} & Conclure avec une prise de position \\
\hline
Ouverture & \esp{en Camerún, nuestras lenguas... son un tesoro} & Relier le texte à la réalité de l'élève \\
\hline
\end{tabularx}
\end{center}

\begin{aretenir}[Les trois ingrédients d'un bon commentaire]
\begin{itemize}
\item \cle{Connecteurs} : \esp{primero} (d'abord), \esp{además} (de plus), \esp{por eso} (c'est pourquoi), \esp{sin embargo} (cependant), \esp{en conclusión} (en conclusion).
\item \cle{Citation + commentaire} : chaque idée s'appuie sur une courte citation entre guillemets, suivie d'une phrase qui l'explique : \esp{Esta descripción muestra que...} (Cette description montre que...).
\item \cle{Avis personnel} en conclusion : \esp{en mi opinión} (à mon avis), \esp{pienso que} (je pense que), \esp{me parece que} (il me semble que).
\end{itemize}
\end{aretenir}

\begin{attention}[Pièges des francophones dans le commentaire]
\begin{itemize}
\item Ne dites pas \esp{el texto habla de} à répétition : variez avec \esp{trata de}, \esp{presenta}, \esp{describe}, \esp{evoca}.
\item Attention au faux ami : \esp{actualmente} signifie « actuellement », pas « en fait ». Pour « en fait », dites \esp{en realidad}.
\item Les citations espagnoles s'ouvrent avec « et se ferment avec » ; n'oubliez pas les signes doubles ¿ et ¡ si vous citez une question ou une exclamation.
\item N'écrivez jamais \esp{según a mí} : la forme correcte est \esp{según yo} ou mieux \esp{en mi opinión}.
\end{itemize}
\end{attention}

\courssub{Production guidée : un commentaire sur une fête camerounaise}

À vous de jouer. Lisez le texte suivant, puis rédigez un commentaire de 10 à 15 lignes en suivant le \cle{plan imposé}.

\begin{textebox}[El Ngondo, fiesta del pueblo sawa]
El Ngondo es la fiesta tradicional del pueblo sawa; se celebra cada año en Douala, a orillas del río Wouri. Durante esta fiesta, los sawa rinden homenaje a los espíritus del agua, los miengu, y piden protección y prosperidad para su pueblo.

La fiesta comienza con ceremonias en la orilla del río: un hombre bucea en las aguas del Wouri con un calabazo y sale del agua con el calabazo lleno, sin mojar el contenido. Este acto misterioso es el momento más importante del Ngondo.

Después, hay danzas tradicionales, cantos en lengua sawa y carreras de piraguas sobre el río. Los jóvenes y los ancianos participan juntos: la fiesta une a todas las generaciones. Los tambores suenan día y noche y la ciudad entera vive al ritmo de la tradición.

Para los sawa, el Ngondo no es solamente una fiesta: es también una escuela de la memoria. Los ancianos transmiten a los jóvenes las historias, las canciones y los valores de sus antepasados. Así, la cultura sawa permanece viva, orgullosa y fuerte, en el corazón de la ciudad moderna de Douala.
\end{textebox}

\textbf{Glossaire}

\begin{itemize}
\item \esp{a orillas de} : au bord de ; \esp{el río} : le fleuve, la rivière
\item \esp{rendir homenaje a} : rendre hommage à ; \esp{los espíritus} : les esprits
\item \esp{bucear} : plonger ; \esp{el calabazo} : la calebasse ; \esp{mojar} : mouiller
\item \esp{la carrera de piraguas} : la course de pirogues ; \esp{los ancianos} : les anciens, les vieillards
\item \esp{los tambores} : les tambours ; \esp{sonar} : sonner, retentir
\item \esp{transmitir} : transmettre ; \esp{los antepasados} : les ancêtres
\item \esp{permanecer} : rester, demeurer ; \esp{orgullosa} : fière
\end{itemize}

\begin{methode}[Plan imposé pour votre commentaire (10 à 15 lignes)]
\begin{enumerate}
\item \textbf{Introduction (2-3 lignes)} : présentez le texte avec \esp{El texto trata de...} et annoncez le thème : la fête du Ngondo et la tradition sawa.
\item \textbf{Première idée (3-4 lignes)} : le Ngondo est une fête religieuse et culturelle. Appuyez-vous sur une citation, par exemple « rinden homenaje a los espíritus del agua », puis commentez-la.
\item \textbf{Deuxième idée (3-4 lignes)} : le Ngondo unit les générations et transmet la mémoire. Citez « los ancianos transmiten a los jóvenes las historias » et commentez avec \esp{Esta cita demuestra que...}.
\item \textbf{Conclusion personnelle (2-3 lignes)} : donnez votre avis avec \esp{En mi opinión...} et ouvrez sur une fête de votre région (fête de la jeunesse, fête traditionnelle de votre village).
\end{enumerate}
\end{methode}

\begin{exoresolu}[Premier paragraphe de développement : modèle de rédaction]
Énoncé : rédigez la première idée du commentaire sur « El Ngondo, fiesta del pueblo sawa » en une phrase d'idée, une citation et une phrase de commentaire.
\tcblower
Solution détaillée :
\esp{Primero, el Ngondo es una fiesta religiosa y cultural del pueblo sawa. El narrador explica que los sawa « rinden homenaje a los espíritus del agua ». Esta cita muestra que la fiesta tiene una dimensión espiritual muy importante: los sawa piden protección a los miengu, los espíritus del río Wouri.}

Traduction : « D'abord, le Ngondo est une fête religieuse et culturelle du peuple sawa. Le narrateur explique que les Sawa "rendent hommage aux esprits de l'eau". Cette citation montre que la fête a une dimension spirituelle très importante : les Sawa demandent la protection des miengu, les esprits du fleuve Wouri. »

Remarquez la structure en trois temps : connecteur + idée (\esp{Primero, el Ngondo es...}), citation entre guillemets, commentaire introduit par \esp{Esta cita muestra que...}. Reproduisez exactement ce schéma pour votre deuxième idée.
\end{exoresolu}

\begin{exemplebox}[Phrases utiles pour citer et commenter]
\begin{itemize}
\item \esp{El Ngondo es la fiesta tradicional del pueblo sawa; se celebra cada año en Douala, a orillas del río Wouri.} « Le Ngondo est la fête traditionnelle du peuple sawa ; elle se célèbre chaque année à Douala, au bord du fleuve Wouri. »
\item \esp{El texto dice que « la fiesta une a todas las generaciones ».} « Le texte dit que "la fête unit toutes les générations". »
\item \esp{Esta frase demuestra que la tradición es un puente entre los ancianos y los jóvenes.} « Cette phrase démontre que la tradition est un pont entre les anciens et les jeunes. »
\item \esp{En mi opinión, el Ngondo es un ejemplo de la riqueza cultural de Camerún.} « À mon avis, le Ngondo est un exemple de la richesse culturelle du Cameroun. »
\end{itemize}
\end{exemplebox}

\courssub{Préparation orale : présenter une tradition de sa région}

Après l'écrit, l'oral. Vous allez présenter \cle{en une minute} une tradition ou une fête de votre région. C'est le moment de réutiliser la grammaire de la leçon : \esp{ser} et \esp{estar} pour décrire et situer, le \cle{présent} pour raconter les activités habituelles.

\begin{methode}[Présenter une tradition à l'oral en quatre étapes]
\begin{enumerate}
\item \textbf{Nom de la fête} avec \esp{ser} : \esp{El Ngondo es la fiesta tradicional del pueblo sawa.} (identité, définition).
\item \textbf{Date et lieu} avec \esp{ser} ou \esp{estar} : \esp{Se celebra en diciembre. Douala está a orillas del río Wouri.} (rappel : \esp{estar} pour la localisation d'une ville ; \esp{celebrarse en} + lieu pour un événement).
\item \textbf{Activités} au présent : \esp{Los jóvenes bailan, los tambores suenan y la gente canta canciones tradicionales.}
\item \textbf{Sentiment personnel} avec \esp{estar orgulloso de} : \esp{Estoy muy orgulloso de esta tradición porque une a mi pueblo.}
\end{enumerate}
\end{methode}

\begin{exemplebox}[Modèle de présentation orale (une minute)]
\esp{La fiesta de la juventud es una fiesta nacional de Camerún. Se celebra el 11 de febrero en todo el país. En mi ciudad, los alumnos desfilan por la mañana y por la tarde hay partidos de fútbol y concursos de danza. Los jóvenes están muy contentos porque es su fiesta. Estoy orgulloso de esta fiesta porque celebra la juventud y el futuro de mi país.}

Traduction : « La fête de la jeunesse est une fête nationale du Cameroun. Elle se célèbre le 11 février dans tout le pays. Dans ma ville, les élèves défilent le matin et l'après-midi il y a des matchs de football et des concours de danse. Les jeunes sont très contents parce que c'est leur fête. Je suis fier de cette fête parce qu'elle célèbre la jeunesse et l'avenir de mon pays. »
\end{exemplebox}

\begin{center}
\begin{tabularx}{\linewidth}{|X|X|X|}
\hline
\rowcolor{popgreenL} Étape & Verbe / temps & Exemple \\
\hline
Nom et définition & \esp{ser} (identité) & \esp{El Ngondo es una fiesta tradicional.} \\
\hline
Localisation & \esp{estar} (lieu) & \esp{Douala está en el litoral.} \\
\hline
Événement & \esp{celebrarse en} & \esp{Se celebra en diciembre.} \\
\hline
Activités habituelles & présent de l'indicatif & \esp{La gente baila y canta.} \\
\hline
Sentiment & \esp{estar} + adjectif & \esp{Estoy orgulloso de mi cultura.} \\
\hline
\end{tabularx}
\end{center}

\begin{attention}[À l'oral : la simplicité est une force]
Évitez les phrases trop longues avec plusieurs subordonnées : à l'oral, elles provoquent des erreurs d'accord et des oublis de verbe. Privilégiez des \cle{phrases simples au présent} avec un lexique précis : \esp{La fiesta es en diciembre. Los jóvenes bailan. Estoy orgulloso.} Trois phrases correctes valent mieux qu'une longue phrase fautive. Pensez aussi à prononcer clairement les voyelles espagnoles, qui ne changent jamais de son.
\end{attention}

\begin{experience}[Tourniquet des traditions]
Par groupes de quatre, chaque élève présente en une minute une tradition de sa région (danse, fête, plat, cérémonie) en suivant les quatre étapes de la méthode. Les autres élèves écoutent et notent : un mot de lexique appris, une phrase correcte avec \esp{ser} ou \esp{estar}, et une question à poser. Puis chaque groupe choisit son meilleur orateur pour une présentation devant la classe entière.
\end{experience}

\courssub{Les étapes de la production et la grille d'évaluation}

Toute production, écrite ou orale, suit le même chemin en cinq étapes. Mémorisez cette frise : c'est votre feuille de route pour toutes les évaluations.

\begin{popfigure}
\begin{tikzpicture}[node distance=0.4cm]
\node[draw=popblue, fill=popblueL, rounded corners, minimum width=2.2cm, minimum height=1cm, align=center] (a) {\textbf{1. leer}\\ lire};
\node[draw=popgreen, fill=popgreenL, rounded corners, minimum width=2.2cm, minimum height=1cm, align=center, right=of a] (b) {\textbf{2. planificar}\\ planifier};
\node[draw=poporange, fill=poporangeL, rounded corners, minimum width=2.2cm, minimum height=1cm, align=center, right=of b] (c) {\textbf{3. redactar}\\ rédiger};
\node[draw=poppurple, fill=poppurpleL, rounded corners, minimum width=2.2cm, minimum height=1cm, align=center, right=of c] (d) {\textbf{4. revisar}\\ relire};
\node[draw=poppink, fill=poppinkL, rounded corners, minimum width=2.2cm, minimum height=1cm, align=center, right=of d] (e) {\textbf{5. presentar}\\ présenter};
\draw[-{Stealth}, thick, popdark] (a) -- (b);
\draw[-{Stealth}, thick, popdark] (b) -- (c);
\draw[-{Stealth}, thick, popdark] (c) -- (d);
\draw[-{Stealth}, thick, popdark] (d) -- (e);
\end{tikzpicture}
\legende{Pasos de la producción : leer, planificar, redactar, revisar, presentar.}
\end{popfigure}

\begin{aretenir}[Grille d'évaluation du commentaire et de l'oral]
\begin{center}
\begin{tabularx}{\linewidth}{|X|X|}
\hline
\rowcolor{popgoldL} Critère & Ce qu'attend le correcteur \\
\hline
Respect du plan & introduction, deux idées avec citations, conclusion personnelle \\
\hline
Exactitude grammaticale & emploi correct de \esp{ser}/\esp{estar} et du présent ; accords \\
\hline
Richesse du lexique & mots de la leçon : \esp{patria, bandera, tradición, diversidad, patrimonio, mestizaje} \\
\hline
Avis personnel & une opinion claire avec \esp{en mi opinión} et une ouverture \\
\hline
\end{tabularx}
\end{center}
\end{aretenir}

\begin{exercice}[À vous de produire]
\begin{enumerate}
\item Rédigez votre commentaire de 10 à 15 lignes sur « El Ngondo, fiesta del pueblo sawa » en suivant le plan imposé. Soulignez vos connecteurs et mettez vos citations entre guillemets.
\item Relisez-vous avec la frise : avez-vous \esp{revisado} vos verbes \esp{ser} et \esp{estar} ? Vos présents sont-ils corrects ?
\item Préparez votre présentation orale d'une minute sur une tradition de votre région et chronométrez-vous.
\end{enumerate}
\end{exercice}

\begin{corrige}[Exercice : indications de correction]
\begin{enumerate}
\item Votre commentaire doit contenir : une introduction avec \esp{El texto trata de...} ; une première idée sur la dimension spirituelle du Ngondo avec la citation « rinden homenaje a los espíritus del agua » ; une deuxième idée sur la transmission entre générations avec la citation « los ancianos transmiten a los jóvenes las historias » ; une conclusion avec \esp{En mi opinión...} et une ouverture sur une fête de votre région.
\item Vérifiez en particulier : \esp{el Ngondo es} (identité, donc \esp{ser}), \esp{Douala está a orillas del Wouri} (lieu, donc \esp{estar}), \esp{se celebra} (présent passif réfléchi), \esp{estoy orgulloso} (sentiment, donc \esp{estar}).
\item Une minute correspond à environ huit à dix phrases simples. Si vous dépassez, supprimez les détails ; si vous êtes trop court, ajoutez une activité et un sentiment.
\end{enumerate}
\end{corrige}

\begin{savaistu}[Une tâche typique du Bac]
Savoir présenter une fête de son pays en espagnol est une tâche fréquente à l'épreuve orale et écrite du Baccalauréat, dans la partie « expression personnelle ». Les examinateurs apprécient les candidats qui parlent de leur culture réelle : le Ngondo de Douala, la fête de la jeunesse du 11 février, le festival des arts et de la culture. Préparer à l'avance un petit discours d'une minute sur une tradition camerounaise, c'est gagner des points le jour de l'examen — et c'est aussi être un ambassadeur de la diversité culturelle du Cameroun dans la langue de Cervantes.
\end{savaistu}

\newpage

\coursec{Activité d'intégration}

\coursec{Leçon EA09 : El patriotismo y la diversidad cultural}

\courssub{Objectifs de l'activité}

À la fin de cette activité d'intégration, tu seras capable de :
\begin{itemize}
\item mobiliser le lexique du patriotisme et de la diversité culturelle (\esp{patria, bandera, himno, tradición, fiesta, lengua, patrimonio, diversidad, mestizaje}) dans une situation de communication réelle ;
\item employer correctement \esp{ser} et \esp{estar} pour décrire une fête et situer un événement ;
\item utiliser le présent de l'indicatif pour parler des traditions ;
\item présenter oralement une tradition devant un public et répondre à des questions en espagnol ;
\item rédiger un court article promotionnel structuré (introduction, développement, conclusion) pour le journal du lycée ;
\item développer une attitude de respect et de valorisation de la diversité culturelle du Cameroun et du monde hispanophone.
\end{itemize}

\courssub{Texte support : « Mi patria, mis culturas »}

Voici l'annonce officielle de la manifestation, affichée au tableau d'information du lycée :

\begin{textebox}[Anuncio oficial — Feria de las culturas]
\textbf{¡GRAN FERIA DE LAS CULTURAS!}\par
El Lycée bilingüe de Yaundé organiza su primera Feria de las culturas el viernes 15 de marzo, de 9 a 14 horas, en el patio central.\par
Cada clase presenta un stand sobre una tradición: una fiesta, una danza, una comida típica o una lengua de Camerún, de España, de México o de Guinea Ecuatorial.\par
Cada stand tiene un cartel con el nombre de la fiesta, la fecha, el lugar, las actividades y los símbolos. Los alumnos explican su tradición en español durante dos minutos y responden a las preguntas del público.\par
El mejor artículo, titulado « La diversidad es nuestra riqueza », será publicado en el periódico del liceo.\par
\textbf{¡La diversidad es nuestra riqueza! ¡Participa!}
\end{textebox}

\emph{Traduction :} « Grande foire des cultures ! Le Lycée bilingue de Yaoundé organise sa première foire des cultures le vendredi 15 mars, de 9 h à 14 h, dans la cour centrale. Chaque classe présente un stand sur une tradition : une fête, une danse, un plat typique ou une langue du Cameroun, d'Espagne, du Mexique ou de Guinée équatoriale. Chaque stand possède une affiche avec le nom de la fête, la date, le lieu, les activités et les symboles. Les élèves expliquent leur tradition en espagnol pendant deux minutes et répondent aux questions du public. Le meilleur article, intitulé "La diversité est notre richesse", sera publié dans le journal du lycée. »

\courssub{Compréhension du texte}

\begin{exercice}[Questions de compréhension]
Réponds en espagnol, par des phrases complètes, aux questions suivantes :
\begin{enumerate}
\item ¿Cuándo y dónde tiene lugar la Feria de las culturas?
\item ¿Qué cuatro países deben estar representados en los stands?
\item ¿Qué elementos obligatorios debe tener el cartel de cada stand?
\item ¿Cuánto dura la presentación oral de cada grupo?
\item ¿Cuál es el título del artículo que se publicará en el periódico del liceo?
\end{enumerate}
\end{exercice}

\begin{corrige}[Exercice 1]
\begin{enumerate}
\item \esp{La Feria de las culturas tiene lugar el viernes 15 de marzo, de 9 a 14 horas, en el patio central del Lycée bilingüe de Yaundé.}
\item \esp{Los cuatro países son: Camerún, España, México y Guinea Ecuatorial.}
\item \esp{El cartel debe tener: el nombre de la fiesta, la fecha, el lugar, las actividades y los símbolos.}
\item \esp{La presentación oral dura dos minutos.}
\item \esp{El título del artículo es « La diversidad es nuestra riqueza ».}
\end{enumerate}
\end{corrige}

\courssub{Lexique thématique : patria, símbolos y tradiciones}

\begin{aretenir}[Lexique essentiel — Patria, símbolos y tradiciones]
\begin{center}
\begin{tabularx}{\linewidth}{|X|X|X|}
\hline
\rowcolor{popblueL} \textbf{Français} & \textbf{Español} & \textbf{Remarque / Exemple} \\
\hline
la patrie & \esp{la patria} & \esp{Mi patria es Camerún.} \\
\hline
le drapeau & \esp{la bandera} & \esp{La bandera de España es roja y amarilla.} \\
\hline
l'hymne & \esp{el himno} & \esp{Cantamos el himno nacional.} \\
\hline
la tradition & \esp{la tradición} & \esp{El Ngondo es una tradición sawa.} \\
\hline
la fête & \esp{la fiesta} & \esp{La fiesta es el 11 de febrero.} \\
\hline
la langue & \esp{la lengua} & \esp{Camerún tiene muchas lenguas.} \\
\hline
le patrimoine & \esp{el patrimonio} & \esp{El Día de los Muertos es patrimonio de la humanidad.} \\
\hline
la diversité & \esp{la diversidad} & \esp{La diversidad cultural es una riqueza.} \\
\hline
le métissage & \esp{el mestizaje} & \esp{México es un país de mestizaje.} \\
\hline
la danse & \esp{la danza} & \esp{La danza del Bikutsi es famosa.} \\
\hline
le plat typique & \esp{el plato típico} & \esp{El ndolé es un plato típico camerunés.} \\
\hline
célébrer & \esp{celebrar} & \esp{Celebramos la fiesta de la juventud.} \\
\hline
partager & \esp{compartir} & \esp{Compartimos la comida con la familia.} \\
\hline
\end{tabularx}
\end{center}
\end{aretenir}

\begin{savaistu}[Cameroun, Espagne, Mexique, Guinée équatoriale : quatre patrimoines]
Le Cameroun compte plus de 250 groupes ethniques et langues : le \esp{Ngondo} (peuple Sawa), la \esp{Fête de la Jeunesse} (11 février), les danses \esp{Bikutsi}, \esp{Makossa}, le \esp{Ndolé} comme plat national. L'Espagne célèbre la \esp{Feria de Abril} (Séville), la \esp{Tomatina} (Buñol), les \esp{Fallas} (Valence). Le Mexique inscrit le \esp{Día de los Muertos} au patrimoine UNESCO ; le \esp{mestizaje} y est central. La Guinée équatoriale, seul pays africain hispanophone, partage la langue et des traditions (danse \esp{Balélé}, \esp{Ibanga}) tout en gardant ses cultures \esp{Fang}, \esp{Bubi}, \esp{Ndowe}. \esp{La diversidad es nuestra riqueza} unit tous ces espaces.
\end{savaistu}

\courssub{Grammaire outillée : ser/estar et le présent des traditions}

\begin{propriete}[Ser ou estar ? — Règle de choix]
\begin{itemize}
\item \textbf{\esp{ser}} : identité, origine, nationalité, caractéristique permanente, \textbf{date/heure}, matière, possession. \\
      \esp{La fiesta es en diciembre.} (date) \quad \esp{El Ngondo es una tradición del pueblo sawa.} (identité) \quad \esp{Yo soy camerunés.} (origine)
\item \textbf{\esp{estar}} : localisation (où?), état temporaire, humeur, résultat d'action, position. \\
      \esp{El stand está en el patio.} (localisation) \quad \esp{Los alumnos están contentos.} (état/humeur) \quad \esp{La puerta está abierta.} (résultat)
\end{itemize}
\emph{Astuce :} \textbf{Date = \esp{ser}} ; \textbf{Où ? = \esp{estar}}.
\end{propriete}

\begin{propriete}[Le présent de l'indicatif pour décrire les traditions]
Pour décrire une coutume, un rite ou ce qui se passe habituellement lors d'une fête, on emploie le \textbf{présent de l'indicatif} (valeur gnomique/habituelle).
\begin{center}
\begin{tabular}{|l|l|l|}
\hline
\rowcolor{popblueL} \textbf{Verbe} & \textbf{Conjugaison (3e pers. sg/pl)} & \textbf{Exemple} \\
\hline
\esp{celebrar} (régulier -ar) & \esp{celebra / celebran} & \esp{El pueblo celebra la fiesta.} \\
\hline
\esp{comer} (régulier -er) & \esp{come / comen} & \esp{La gente come platos típicos.} \\
\hline
\esp{vivir} (régulier -ir) & \esp{vive / viven} & \esp{Viven la tradición con alegría.} \\
\hline
\esp{empezar} (e $\to$ ie) & \esp{empieza / empiezan} & \esp{La fiesta empieza al mediodía.} \\
\hline
\esp{vestirse} (e $\to$ i, pronominal) & \esp{se viste / se visten} & \esp{La gente se viste de colores.} \\
\hline
\esp{tener} (irrégulier) & \esp{tiene / tienen} & \esp{Tienen mucha fe.} \\
\hline
\esp{poner} (irrégulier) & \esp{pone / ponen} & \esp{Ponen flores en el altar.} \\
\hline
\esp{hacer} (irrégulier) & \esp{hace / hacen} & \esp{Hacen altares en casa.} \\
\hline
\end{tabular}
\end{center}
\end{propriete}

\begin{attention}[Erreurs fréquentes des francophones]
\begin{itemize}
\item \textbf{Date :} Ne dis pas \esp{la fiesta está en diciembre} $\rightarrow$ \esp{la fiesta \textbf{es} en diciembre} (\esp{ser} pour la date).
\item \textbf{Lieu de l'événement :} \esp{La fiesta es en Madrid} (c'est un événement qui a lieu là-bas = \esp{ser}) \textbf{mais} \esp{El stand está en Madrid} (localisation physique = \esp{estar}).
\item \textbf{Faux ami « fiesta » :} \esp{Fiesta} = fête (cérémonie, célébration). Pour « faire la fête » (sortir, s'amuser) : \esp{salir de fiesta}, \esp{divertirse}, \esp{pasar un buen rato}. Ne pas dire \esp{hacer una fiesta} pour « faire la fête ».
\item \textbf{Accents :} \esp{tradición}, \esp{diversidad}, \esp{patrimonio} (pas d'accent), \esp{México}, \esp{héroe}, \esp{canción}.
\item \textbf{Genre :} \esp{el himno}, \esp{el patrimonio}, \esp{el mestizaje}, \esp{el plato} (masculins) ; \esp{la patria}, \esp{la bandera}, \esp{la tradición}, \esp{la fiesta}, \esp{la lengua}, \esp{la diversidad}, \esp{la danza} (féminins).
\end{itemize}
\end{attention}

\courssub{Méthode du commentaire de texte}

\begin{methode}[Commentaire de texte : présenter une tradition (plan type)]
\begin{enumerate}
\item \textbf{Introduction} (2--3 lignes) :
      \begin{itemize}
      \item Accroche : nom de la fête, pays, type de tradition (religieuse, civile, ancestrale).
      \item Définition rapide : \esp{Es una tradición antigua / Es parte del patrimonio...}
      \item Annonce du plan : \esp{En este artículo vamos a ver las actividades y los símbolos de esta fiesta.}
      \end{itemize}
\item \textbf{Développement} (2 paragraphes distincts, 3--4 lignes chacun) :
      \begin{itemize}
      \item \textbf{Idée 1 : Les activités} (ce que font les gens). Verbes d'action au \textbf{présent} : \esp{preparan, cantan, bailan, comen, comparten, visitan}. Utiliser \esp{estar} pour la localisation momentanée : \esp{están en las casas, están en la calle}.
      \item \textbf{Idée 2 : Les symboles et les valeurs} (ce que la fête représente). \textbf{\esp{ser}} pour l'identité : \esp{es un símbolo de..., es parte del patrimonio, es mestizaje}. \textbf{\esp{estar}} pour l'ambiance/état : \esp{las calles están llenas de color, la gente está emocionada}.
      \end{itemize}
\item \textbf{Conclusion} (2--3 lignes) :
      \begin{itemize}
      \item Bilan personnel : \esp{En mi opinión..., Me parece que...}
      \item Lien avec le Cameroun / ouverture : \esp{En Camerún también..., La diversidad nos une.}
      \item Phrase choc finale (exclamative) : \esp{¡La diversidad es nuestra riqueza!}
      \end{itemize}
\end{enumerate}
\end{methode}

\courssub{Activité d'intégration : La Feria de las culturas}

\begin{experience}[Organisation des stands — Feria de las culturas]
Travaillez par groupes de 4 élèves. Suivez les étapes dans l'ordre. Durée totale : environ 1 h 10.

\begin{enumerate}
\item \textbf{Choix de la tradition (5 min).} En groupe, choisissez une tradition parmi celles proposées : la fête de la jeunesse (11 février) ou le Ngondo au Cameroun, la Feria de Abril en Espagne, le Día de los Muertos au Mexique, ou une danse traditionnelle de Guinée équatoriale (Balélé, Ibanga). Justifiez votre choix en une phrase en espagnol : \esp{Elegimos... porque...}

\item \textbf{Préparation de l'affiche (15 min).} Rédigez une affiche en espagnol avec obligatoirement : le nom de la fête, la date, le lieu, deux ou trois activités et deux symboles (par exemple \esp{la bandera}, \esp{el himno}, un vêtement, un plat). Utilisez des phrases courtes au présent. \emph{Exemple :} \esp{La fiesta es el 11 de febrero. El lugar es Yaoundé. Actividades: desfile, himno, baile. Símbolos: la bandera, el pagne.}

\item \textbf{Préparation de la présentation orale (15 min).} Rédigez puis répétez une présentation de 2 minutes qui contient : au moins \textbf{8 mots du lexique} de la leçon, \textbf{2 emplois de} \esp{ser}, \textbf{2 emplois de} \esp{estar} et \textbf{3 verbes conjugués au présent} (autres que ser/estar). Répartissez la parole entre les quatre membres du groupe.

\item \textbf{Animation du stand (10 min par groupe).} Présentez votre tradition à la classe. Les autres élèves posent au moins deux questions en espagnol, par exemple : \esp{¿Dónde está el stand?}, \esp{¿Cuándo es la fiesta?}, \esp{¿Qué comen las personas?}, \esp{¿Cuál es el símbolo principal?}

\item \textbf{Vote (5 min).} Votez pour le stand le plus convaincant (clarté, correction de la langue, richesse du lexique). Vous ne pouvez pas voter pour votre propre groupe.

\item \textbf{Rédaction de l'article (20 min).} Chaque groupe rédige un court article de 10 à 12 lignes intitulé \esp{« La diversidad es nuestra riqueza »} pour le journal du lycée, en suivant le plan du commentaire de texte (voir Méthode ci-dessus) : introduction, développement en deux idées (activités puis symboles/valeurs), conclusion personnelle avec lien camerounais et exclamation finale.
\end{enumerate}
\end{experience}

\courssub{Ressources à mobiliser}

\begin{aretenir}[Lexique utile — Rappel rapide]
\esp{la patria} (la patrie), \esp{la bandera} (le drapeau), \esp{el himno} (l'hymne), \esp{la tradición} (la tradition), \esp{la fiesta} (la fête), \esp{la lengua} (la langue), \esp{el patrimonio} (le patrimoine), \esp{la diversidad} (la diversité), \esp{el mestizaje} (le métissage), \esp{la danza} (la danse), \esp{el plato típico} (le plat typique), \esp{celebrar} (célébrer), \esp{compartir} (partager), \esp{el desfile} (le défilé), \esp{el altar} (l'autel), \esp{la vela} (la bougie), \esp{la flor} (la fleur), \esp{el cementerio} (le cimetière), \esp{el pan de muerto} (le pain des morts).
\end{aretenir}

\begin{propriete}[Rappel : Ser / Estar / Présent]
\begin{itemize}
\item \esp{ser} : date (\esp{es el 1 de noviembre}), identité (\esp{es una tradición}), caractéristique (\esp{es importante}).
\item \esp{estar} : lieu (\esp{está en el patio}), état (\esp{están contentos}), ambiance (\esp{está lleno de gente}).
\item Présent habituel : \esp{la gente prepara, canta, baila, come, comparte, visita, pone, hace}.
\end{itemize}
\end{propriete}

\courssub{Commentaire modèle et production guidée}

\begin{exoresolu}[Article modèle pour le journal du lycée]
Rédige un article de 10 à 12 lignes intitulé \esp{« La diversidad es nuestra riqueza »} pour présenter la tradition choisie par ton groupe et promouvoir la diversité culturelle.
\tcblower
\textbf{Production modèle (groupe « Día de los Muertos ») :}\par
\esp{\textbf{La diversidad es nuestra riqueza}\par
Hoy nuestro grupo presenta el Día de los Muertos, una fiesta muy importante de México. Es una tradición antigua y es parte del patrimonio cultural de la humanidad.\par
La fiesta es el 1 y el 2 de noviembre. Las familias están en las casas y en los cementerios, y preparan altares con flores, velas y comida. La gente canta, baila y comparte platos típicos como el pan de muerto.\par
Esta fiesta es también un símbolo de mestizaje: mezcla tradiciones indígenas y españolas. Los colores de las flores están por todas partes y las calles están llenas de música.\par
En conclusión, el Día de los Muertos muestra que cada cultura es un tesoro. En Camerún también tenemos muchas lenguas, danzas y fiestas. La diversidad es nuestra riqueza: ¡descubramos las tradiciones de los demás!}\par
\textbf{Commentaire des choix de langue :}
\begin{itemize}
\item \textbf{Lexique mobilisé (10 mots) :} \esp{fiesta, tradición, patrimonio, símbolo, mestizaje, lenguas, danzas, diversidad, altar, pan de muerto}.
\item \textbf{\esp{ser} (3 emplois) :} \esp{es una tradición antigua} (caractéristique), \esp{la fiesta es el 1 y el 2 de noviembre} (date), \esp{es parte del patrimonio} (identité), \esp{es un símbolo de mestizaje} (identité), \esp{cada cultura es un tesoro} (définition).
\item \textbf{\esp{estar} (2 emplois) :} \esp{las familias están en las casas} (localisation), \esp{las calles están llenas de música} (état/ambiance).
\item \textbf{Présents de l'indicatif (coutumes) :} \esp{preparan, canta, baila, comparte, muestra, tenemos, mezclan (impliqué dans 'mezcla')}.
\item \textbf{Structure respectée :} Introduction (présentation), Idée 1 = activités (préparent autels, chantent, dansent, partagent), Idée 2 = symboles/valeurs (mestizaje, couleurs, musique), Conclusion (lien Cameroun + exclamation hortative \esp{¡descubramos!}).
\end{itemize}
\end{exoresolu}

\courssub{Bilan de l'activité}

Cette activité t'a permis de mettre en pratique, dans une situation de communication proche de la réalité, l'ensemble des acquis de la leçon :
\begin{itemize}
\item tu as lu et compris un document authentique (l'annonce de la foire) et répondu à des questions de compréhension ;
\item tu as réutilisé le lexique du patriotisme et de la diversité culturelle dans un contexte précis (affiche, oral, article) ;
\item tu as consolidé l'emploi de \esp{ser} et \esp{estar} ainsi que le présent des traditions (verbes réguliers, alternants, irréguliers, pronominaux) ;
\item tu t'es exprimé oralement devant un public (2 min) et tu as réagi à des questions imprévues (interaction) ;
\item tu as rédigé un article structuré selon le plan du commentaire de texte (introduction, développement en deux idées, conclusion), compétence essentielle pour la classe de Première et le Baccalauréat.
\end{itemize}

\begin{aretenir}[L'essentiel]
Présenter une tradition, c'est savoir dire \textbf{ce qu'elle est} (\esp{ser}), \textbf{où et comment elle se déroule} (\esp{estar}) et \textbf{ce que font les gens} (présent de l'indicatif). Valoriser la diversité, c'est affirmer que chaque culture — camerounaise, espagnole, mexicaine ou équato-guinéenne — est une richesse pour tous : \esp{La diversidad es nuestra riqueza}.
\end{aretenir}

\newpage

\coursec{Méthodes et exercices résolus}

\coursec{Leçon EA09 : El patriotismo y la diversidad cultural}

\courssub{Texte support : « Mi patria, mis culturas »}

\begin{textebox}[Texte original — Témoignage d'une lycéenne camerounaise]
\esp{Soy camerunesa y estoy orgullosa de mi patria. En mi país hay más de doscientas lenguas y muchas tradiciones diferentes. La bandera verde, roja y amarilla es el símbolo de nuestra unidad. Cada año celebramos la fiesta nacional el veinte de mayo. Mi abuela habla ewondo y mi padre habla fulfuldé, pero todos hablamos francés e inglés en la escuela. Esta diversidad es nuestro patrimonio más precioso. Como en México o en Guinea Ecuatorial, el mestizaje de las culturas hace fuerte a nuestra nación.}
\end{textebox}

\textbf{Traduction :} « Je suis camerounaise et je suis fière de ma patrie. Dans mon pays, il y a plus de deux cents langues et beaucoup de traditions différentes. Le drapeau vert, rouge et jaune est le symbole de notre unité. Chaque année, nous célébrons la fête nationale le 20 mai. Ma grand-mère parle l'ewondo et mon père parle le fulfuldé, mais nous parlons tous français et anglais à l'école. Cette diversité est notre patrimoine le plus précieux. Comme au Mexique ou en Guinée équatoriale, le métissage des cultures rend nos nations fortes. »

\begin{definition}[Lexique clé de la leçon : \cle{patria, símbolos y tradiciones}]
\begin{center}
\begin{tabularx}{\linewidth}{|X|X|X|}
\hline
\rowcolor{popblueL} \textbf{Español} & \textbf{Français} & \textbf{Prononciation / Note} \\
\hline
\esp{la patria} & la patrie & \emph{nom f.} — pas d'accent (llana) \\
\esp{la bandera} & le drapeau & \emph{nom f.} — \esp{bandera nacional} \\
\esp{el himno} & l'hymne & \emph{nom m.} — \esp{himno nacional} (pas d'accent) \\
\esp{la tradición} & la tradition & \emph{nom f.} — \textbf{accent sur ó} (aguda en -n) \\
\esp{la fiesta} & la fête & \emph{nom f.} — \esp{fiesta nacional, fiesta tradicional} \\
\esp{la lengua} & la langue & \emph{nom f.} — syn. \esp{el idioma} \\
\esp{el patrimonio} & le patrimoine & \emph{nom m.} — \textbf{accent sur ó} (esdrújula) \\
\esp{la diversidad} & la diversité & \emph{nom f.} — \textbf{accent sur á} (aguda en -d) \\
\esp{el mestizaje} & le métissage & \emph{nom m.} — \textbf{accent sur é} (aguda en -e) \\
\hline
\end{tabularx}
\end{center}
\end{definition}

\courssub{Savoir-faire 1 : Lire et comprendre un texte sur le patriotisme et la diversité culturelle}

\begin{methode}[Comprendre un texte sur la \cle{patria} y la \cle{diversidad}]
\begin{enumerate}
\item \textbf{Lire le titre et la source.} Le titre \esp{Mi patria, mis culturas} annonce un témoignage personnel sur l'identité culturelle. Repérer les mots-clés : \esp{patria, bandera, himno, tradición, lengua, patrimonio, diversidad, mestizaje}.
\item \textbf{Première lecture globale.} Identifier le type de texte (témoignage / \esp{testimonio}), le locuteur (\esp{¿quién habla?} : une jeune Camerounaise) et le sujet (\esp{¿de qué habla?} : fierté nationale, langues, symboles, métissage).
\item \textbf{Deuxième lecture analytique.} Souligner le champ lexical de la patrie et de la diversité ; repérer les verbes \esp{ser} (identité, caractéristiques permanentes) et \esp{estar} (états, sentiments, localisation).
\item \textbf{Deviner les mots inconnus par le contexte.} Ex. : \esp{El mestizaje enriquece nuestra cultura} → \esp{enriquece} (enrichit) aide à comprendre \esp{mestizaje}.
\item \textbf{Répondre aux questions} en recopiant des éléments du texte, sans phrases entières hors sujet. Répondre en espagnol, phrases complètes : \esp{El autor dice que...}, \esp{Según el texto...}.
\item \textbf{Vérifier} : verbe conjugué à la bonne personne, accords (genre/nombre), accents (\esp{fulfuldé, según, precioso}).
\end{enumerate}
\end{methode}

\begin{exoresolu}[Compréhension de texte : « Mi patria, mis culturas »]
\textbf{Texte :} (voir cadre ci-dessus)

\textbf{Questions :}
\begin{enumerate}
\item \esp{¿De qué nacionalidad es la narradora?}
\item \esp{¿Cuántas lenguas hay en su país?}
\item \esp{¿Cuándo se celebra la fiesta nacional?}
\item \esp{¿Qué lenguas hablan su abuela y su padre?}
\item \esp{¿Qué es, según el texto, el patrimonio más precioso?}
\end{enumerate}
\tcblower
\textbf{Raisonnement :} Questions de repérage direct. On cherche l'information exacte, puis on rédige une phrase complète en reprenant la structure de la question (changement de personne si nécessaire : \esp{mi} → \esp{su}).

\textbf{Réponses rédigées :}
\begin{enumerate}
\item \esp{La narradora es camerunesa.} (La narratrice est camerounaise.) — \esp{Soy camerunesa} (1\iere{} phrase). \emph{Note : forme standard RAE \esp{camerunesa}.}
\item \esp{En su país hay más de doscientas lenguas.} (Il y a plus de deux cents langues.) — \esp{hay} (il y a) invariable ; \esp{doscientas} accord féminin pluriel avec \esp{lenguas}.
\item \esp{La fiesta nacional se celebra el veinte de mayo.} (Elle se célèbre le 20 mai.) — Date avec article défini : \esp{el veinte de mayo}. Voix passive réflexe \esp{se celebra}.
\item \esp{Su abuela habla ewondo y su padre habla fulfuldé.} — Changement \esp{mi} → \esp{su} (3\ieme{} personne). \esp{fulfuldé} garde son accent.
\item \esp{Según el texto, la diversidad es el patrimonio más precioso.} — \esp{según} (accent) ; \esp{patrimonio} masculin → \esp{el patrimonio, precioso}.
\end{enumerate}
\textbf{Conclusion :} Réponses complètes, verbes conjugués, accords et accents respectés.
\end{exoresolu}

\courssub{Savoir-faire 2 : Présenter une tradition ou une fête à l'oral et à l'écrit}

\begin{methode}[Présenter une fête ou une tradition]
\begin{enumerate}
\item \textbf{Nommer la fête et la situer} avec \esp{ser} (identité) et \esp{estar} (lieu physique) : \esp{El Ngondo es una fiesta tradicional. Está en Douala.}
\item \textbf{Donner la date} : \esp{Se celebra el veinte de mayo} / \esp{Tiene lugar en diciembre} (sans article devant le mois).
\item \textbf{Décrire les activités au présent} (habitudes, vérité générale) : \esp{La gente canta el himno, baila danzas tradicionales y come platos típicos.}
\item \textbf{Parler des symboles} : \esp{La bandera tiene tres colores: verde, rojo y amarillo.}
\item \textbf{Exprimer son opinion} : \esp{Me gusta esta fiesta porque...}, \esp{Para mí, es una tradición muy bonita.}
\item \textbf{Conclure sur la valeur culturelle} : \esp{Esta fiesta muestra la riqueza de nuestra cultura.}
\end{enumerate}
\textbf{Structures à retenir :} \esp{se celebra}, \esp{tiene lugar}, \esp{la gente + verbe au singulier}, \esp{durante la fiesta}, \esp{es una tradición muy antigua}.
\end{methode}

\begin{exoresolu}[Production écrite : présenter une fête]
\textbf{Énoncé :} \esp{Escribe un párrafo de 8 a 10 líneas para presentar una fiesta tradicional de tu región a un amigo español. Habla de la fecha, las actividades, los símbolos y tu opinión.}
\tcblower
\textbf{Raisonnement :} Plan de la fiche méthode : nom + date (\esp{ser/estar}), activités au présent, symboles, opinion, conclusion. Présent de l'indicatif pour habitudes.

\textbf{Production modèle :}

\esp{En mi región celebramos el Ngondo, la fiesta tradicional del pueblo sawa. Tiene lugar cada año en diciembre, cerca del río Wouri, en Douala. Durante la fiesta, la gente baila danzas tradicionales, canta canciones antiguas y participa en carreras de piraguas. Los jefes tradicionales llevan ropas muy bonitas y bendicen a la comunidad. La fiesta es muy importante porque une a las familias y transmite nuestras tradiciones a los jóvenes. Me gusta mucho el Ngondo porque muestra la riqueza de nuestra cultura. Para mí, es el patrimonio más valioso de mi región.}

\textbf{Justifications grammaticales :}
\begin{itemize}
\item \esp{celebramos} : présent \esp{celebrar}, 1\iere{} pl. (communauté).
\item \esp{Tiene lugar} : locution figée « avoir lieu » ; \esp{en diciembre} (mois sans article).
\item \esp{la gente baila} : \esp{la gente} = singulier en espagnol (contrairement à « les gens »).
\item \esp{es muy importante} : \esp{ser} (caractéristique permanente de la fête).
\item \esp{Me gusta} : construction inversée, \esp{gusta} (3\ieme{} sg) car sujet = \esp{el Ngondo}.
\end{itemize}
\textbf{Conclusion :} Paragraphe structuré, présent, lexique de la leçon (\esp{tradición, cultura, patrimonio}).
\end{exoresolu}

\courssub{Savoir-faire 3 : Promouvoir la diversité culturelle en espagnol}

\begin{methode}[Défendre et promouvoir la diversité culturelle]
\begin{enumerate}
\item \textbf{Affirmer la valeur de la diversité} : \esp{La diversidad cultural es una riqueza.} / \esp{Cada lengua es un tesoro.}
\item \textbf{Donner des exemples concrets} (Cameroun, Espagne, Mexique, Guinée équatoriale) : \esp{En Camerún hay más de doscientas lenguas. En España se hablan español, catalán, gallego y euskera.}
\item \textbf{Utiliser des verbes d'opinion et de devoir} : \esp{Debemos respetar todas las culturas.} / \esp{Es importante proteger las lenguas minoritarias.} / \esp{Hay que valorar el mestizaje.}
\item \textbf{Employer des connecteurs} : \esp{además} (de plus), \esp{también} (aussi), \esp{por eso} (c'est pourquoi), \esp{sin embargo} (cependant).
\item \textbf{Conclure par un appel} : \esp{¡Defendamos nuestra diversidad!} / \esp{¡Viva la diversidad cultural!}
\end{enumerate}
\textbf{Structures à retenir :} \esp{es importante + infinitif}, \esp{debemos + infinitif}, \esp{hay que + infinitif}, \esp{gracias a la diversidad...}.
\end{methode}

\begin{exoresolu}[Thème : promouvoir la diversité]
\textbf{Énoncé :} Traduisez en espagnol :
\begin{enumerate}
\item La diversité culturelle est une richesse pour notre pays.
\item Au Cameroun, il y a plus de deux cents langues.
\item Nous devons respecter toutes les traditions.
\item Il est important de protéger notre patrimoine.
\item C'est pourquoi je suis fier de ma patrie.
\end{enumerate}
\tcblower
\textbf{Raisonnement et solutions :}
\begin{enumerate}
\item \esp{La diversidad cultural es una riqueza para nuestro país.} — \esp{ser} (caractère permanent) ; \esp{riqueza} f. → \esp{una riqueza}.
\item \esp{En Camerún hay más de doscientas lenguas.} — « il y a » = \esp{hay} (invariable) ; \esp{doscientas} accord f. pl. avec \esp{lenguas}.
\item \esp{Debemos respetar todas las tradiciones.} — \esp{deber + infinitif} ; \esp{todas} f. pl. avec \esp{tradiciones}.
\item \esp{Es importante proteger nuestro patrimonio.} — « il est important de » = \esp{es importante + infinitif} (pas de subjonctif, sujet impersonnel).
\item \esp{Por eso, estoy orgulloso de mi patria.} — « c'est pourquoi » = \esp{por eso} ; « être fier » = \esp{estar orgulloso de} (état/sentiment) ≠ \esp{*ser orgulloso} (orgueilleux) ; accord \esp{orgulloso} (m.) / \esp{orgullosa} (f.).
\end{enumerate}
\textbf{Conclusion :} Structures de la méthode appliquées ; pièges évités (\esp{hay}, \esp{estar orgulloso}, accord \esp{doscientas}).
\end{exoresolu}

\courssub{Savoir-faire 4 : Employer ser/estar et le présent pour décrire traditions et symboles}

\begin{propriete}[Choix entre \cle{ser} et \cle{estar} : règle de décision]
\begin{center}
\begin{tabularx}{\linewidth}{|X|X|X|}
\hline
\rowcolor{popblueL} \textbf{Emploi} & \textbf{Verbe} & \textbf{Exemples} \\
\hline
Identité, origine, nationalité, matière, caractéristique permanente, définition & \esp{SER} & \esp{El himno es un símbolo. Somos cameruneses. La fiesta es en mayo. La bandera es de tela.} \\
\hline
Localisation (personnes, objets, lieux concrets), état temporaire, sentiment, action en cours (\esp{estar + gerundio}) & \esp{ESTAR} & \esp{La bandera está en la plaza. Estamos contentos. La gente está bailando. Yaoundé está en el centro.} \\
\hline
\multicolumn{3}{|l|}{\textbf{Cas piège — \esp{orgulloso}} : \esp{ser orgulloso} = être orgueilleux (défaut) ; \esp{estar orgulloso de} = être fier de (sentiment).} \\
\hline
\end{tabularx}
\end{center}
\end{propriete}

\begin{popfigure}
\begin{tikzpicture}[
  node distance=1.2cm and 1.5cm,
  decision/.style={diamond, draw=popblue, thick, fill=popblueL, text width=2.8cm, align=center, inner sep=4pt},
  action/.style={rectangle, rounded corners, draw=popgreen, thick, fill=popgreenL, text width=2.5cm, align=center, inner sep=4pt},
  start/.style={ellipse, draw=poporange, thick, fill=poporangeL, text width=2cm, align=center},
  arrow/.style={-Stealth, thick, popdark}
]
\node[start] (start) {Verbe \esp{ÊTRE} ?};
\node[decision, below of=start] (q1) {Identité, origine, nationalité, date/heure d'un événement, caractéristique permanente ?};
\node[action, below left=1.5cm and 1cm of q1, align=center] (ser){\textbf{SER}\\\esp{soy, eres, es, somos, sois, son}};
\node[decision, below right=1.5cm and 1cm of q1] (q2) {Localisation physique, état passager, sentiment, action en cours ?};
\node[action, below left=1.5cm and 1cm of q2, align=center] (estar){\textbf{ESTAR}\\\esp{estoy, estás, está, estamos, estáis, están}};
\node[action, below right=1.5cm and 1cm of q2, align=center] (trap){\textbf{PIÈGE}\\\esp{ser orgulloso} (orgueilleux)\\\esp{estar orgulloso de} (fier)};
\draw[arrow] (start) -- (q1);
\draw[arrow] (q1) -- node[left, pos=0.5, align=center] {\textbf{SÍ}} (ser);
\draw[arrow] (q1) -- node[right, pos=0.5, align=center] {\textbf{NO}} (q2);
\draw[arrow] (q2) -- node[left, pos=0.5, align=center] {\textbf{SÍ}} (estar);
\draw[arrow] (q2) -- node[right, pos=0.5, align=center] {\textbf{orgulloso?}} (trap);
\end{tikzpicture}
\legende{Arbre de décision \esp{ser} vs \esp{estar} (niveau A2+/B1).}
\end{popfigure}

\begin{methode}[Conjugaison au présent de l'indicatif (révision)]
\begin{center}
\begin{tabular}{|l|l|l|l|}
\hline
\rowcolor{popblueL} \textbf{Personne} & \textbf{SER (irrégulier)} & \textbf{ESTAR (irrégulier)} & \textbf{CELEBRAR (régulier -ar)} \\
\hline
yo & soy & estoy & celebro \\
tú & eres & estás & celebras \\
él/ella/usted & es & está & celebra \\
nosotros/as & somos & estamos & celebramos \\
vosotros/as & sois & estáis & celebráis \\
ellos/ellas/ustedes & son & están & celebran \\
\hline
\end{tabular}
\end{center}
\textbf{Rappel :} Le présent exprime l'habitude (\esp{Cada año celebramos...}), la vérité générale (\esp{La diversidad es una riqueza}) et l'action immédiate (\esp{Ahora estoy cantando}).
\end{methode}

\begin{exoresolu}[Version et transformation : \esp{ser/estar} au présent]
\textbf{Énoncé A (Version) :} Traduisez en français : \esp{La bandera de Guinea Ecuatorial es un símbolo de la independencia. Durante las fiestas, la gente está muy orgullosa y canta el himno nacional.}

\textbf{Énoncé B (Transformation) :} Complétez avec \esp{ser} ou \esp{estar} au présent :
\begin{enumerate}
\item \esp{México \_\_\_ un país de gran diversidad cultural.}
\item \esp{Los jóvenes \_\_\_ contentos durante la fiesta.}
\item \esp{La fiesta nacional \_\_\_ en mayo.}
\item \esp{Nosotros \_\_\_ orgullosos de nuestras tradiciones.}
\end{enumerate}
\tcblower
\textbf{Solution A :} « Le drapeau de la Guinée équatoriale est un symbole de l'indépendance. Pendant les fêtes, les gens sont très fiers et chantent l'hymne national. »
\begin{itemize}
\item \esp{es} → « est » : caractéristique permanente du drapeau (\esp{ser}).
\item \esp{está muy orgullosa} → « sont très fiers » : état temporaire (\esp{estar}) ; \esp{la gente} singulier → \esp{está} ; accord \esp{orgullosa} (f. sg) avec \esp{gente}.
\end{itemize}

\textbf{Solution B :}
\begin{enumerate}
\item \esp{México \textbf{es} un país de gran diversidad cultural.} — Identité, caractéristique permanente → \esp{ser}.
\item \esp{Los jóvenes \textbf{están} contentos durante la fiesta.} — État temporaire → \esp{estar}.
\item \esp{La fiesta nacional \textbf{es} en mayo.} — \textbf{Règle clé :} Pour les événements, la date ET le lieu de l'événement s'expriment avec \esp{ser} (\esp{La fiesta es en mayo / es en el estadio}). \esp{Estar} s'emploie pour la localisation de personnes/objets (\esp{Nosotros estamos en el estadio}).
\item \esp{Nosotros \textbf{estamos} orgullosos de nuestras tradiciones.} — Sentiment → \esp{estar} ; \esp{orgullosos} m. pl.
\end{enumerate}
\textbf{Conclusion :} Le choix \esp{ser/estar} dépend du sens (permanent/temporaire, événement vs localisation physique) ; le présent décrit les traditions.
\end{exoresolu}

\courssub{Savoir-faire 5 : Rédiger un commentaire de texte sur la citoyenneté culturelle}

\begin{methode}[Commenter un texte sur la patrie et la diversité — Structure Bac]
\begin{enumerate}
\item \textbf{Introduction :} Présenter le texte (type, thème, locuteur) : \esp{Este texto es un testimonio sobre el patriotismo y la diversidad cultural en Camerún.}
\item \textbf{Résumé bref :} 2-3 phrases max : \esp{La narradora habla de su país, de sus lenguas y de sus fiestas.}
\item \textbf{Analyse linguistique :} Relever le champ lexical (\esp{patria, bandera, himno, tradición, patrimonio}), les emplois de \esp{ser/estar}, le présent des habitudes, et commenter leur fonction : \esp{El uso del presente muestra que las tradiciones son costumbres vivas.}
\item \textbf{Opinion personnelle argumentée} : \esp{En mi opinión, este texto muestra que la diversidad es una riqueza. Yo también estoy orgulloso de mi cultura porque...}
\item \textbf{Conclusion :} Ouverture : \esp{Como en México o en Guinea Ecuatorial, el mestizaje hace fuerte a las naciones.}
\end{enumerate}
\end{methode}

\begin{exoresolu}[Commentaire guidé]
\textbf{Énoncé :} À partir du texte « Mi patria, mis culturas », rédigez en espagnol un commentaire de 10 à 12 lignes : présentation, résumé, analyse vocabulaire/verbes, opinion personnelle.
\tcblower
\textbf{Raisonnement :} Suivre les 5 étapes ; phrases liées avec connecteurs (\esp{además, por eso, en mi opinión}).

\textbf{Commentaire modèle :}

\esp{Este texto es un testimonio personal sobre el patriotismo y la diversidad cultural. La narradora, una joven camerunesa, habla de su país con orgullo. Primero, presenta los símbolos de la nación: la bandera verde, roja y amarilla y la fiesta nacional del veinte de mayo. Después, habla de las lenguas de su familia: su abuela habla ewondo y su padre habla fulfuldé. El vocabulario del texto pertenece al campo léxico de la patria: patria, bandera, himno, tradición, patrimonio, diversidad y mestizaje. Los verbos están en presente porque describen costumbres y realidades permanentes: "hay más de doscientas lenguas", "celebramos la fiesta nacional". También notamos el contraste entre "soy camerunesa" (ser, identidad) y "estoy orgullosa" (estar, sentimiento). En mi opinión, este texto es muy interesante porque muestra que la diversidad no es un problema sino una riqueza. Yo también pienso que debemos proteger nuestras lenguas y nuestras tradiciones. Por eso, como dice la narradora, la diversidad es nuestro patrimonio más precioso.}

\textbf{Justifications :} Présent de vérité générale ; \esp{ser} (identité), \esp{estar} (sentiment) ; connecteurs variés ; lexique réemployé ; opinion marquée (\esp{en mi opinión, yo también pienso}).

\textbf{Conclusion :} Structure Bac respectée : présentation, résumé, analyse, opinion.
\end{exoresolu}

\begin{savaistu}[Le Cameroun, l'Espagne, le Mexique et la Guinée équatoriale : quatre mosaïques culturelles]
\begin{itemize}
\item \textbf{Cameroun :} « Afrique en miniature ». 250+ langues (ewondo, fulfuldé, bassa, bamileke...), 2 langues officielles (français, anglais), fêtes : 20 mai (fête nationale), Ngondo (Sawa), Ngouon (Bamoun), Medumba. \esp{Patriotismo} = unité dans la diversité.
\item \textbf{Espagne :} 17 communautés autonomes. Langues co-officielles : \esp{catalán, gallego, euskera} (basque) + \esp{castellano}. Fêtes : \esp{La Tomatina, San Fermín, Fallas}. \esp{Mestizaje} historique (romains, wisigoths, arabes, juifs).
\item \textbf{Mexique :} \esp{Mestizaje} central (indigène + espagnol). 68 langues indigènes (nahuatl, maya...). Fêtes : \esp{Día de Muertos} (patrimoine UNESCO), \esp{Guelaguetza}. \esp{Patria} = \esp{México lindo y querido}.
\item \textbf{Guinée équatoriale :} Seul pays africain hispanophone officiel. Langues : espagnol, français, portugais + langues nationales (fang, bube, ndowe). Fête : 12 octobre (indépendance). Pont culturel entre Afrique et monde hispanique.
\end{itemize}
\textbf{Point commun :} La \cle{diversidad} n'est pas une menace, c'est un \cle{patrimonio} à protéger (\esp{Debemos proteger nuestra diversidad}).
\end{savaistu}

\begin{attention}[Les 5 erreurs qui coûtent des points au Bac]
\begin{enumerate}
\item \textbf{Accents oubliés ou mis au hasard.} \textbf{Mots AVEC accent (tilde) :} \esp{tradición, diversidad, patrimonio, mestizaje, canción, también, está (verbe), según, fulfuldé, orgullós? (non, llana fin s -> pas d'accent), camerunesa (llana fin a -> pas d'accent)}. \textbf{Mots SANS accent :} \esp{patria, bandera, himno, fiesta, lengua, cultura, rico, orgulloso, celebra, hay, esta (dém.)}. \textbf{Piège :} \esp{está} (verbe) ≠ \esp{esta} (adjectif démonstratif).
\item \textbf{Confondre \esp{ser} et \esp{estar}.} \esp{*Soy orgulloso de mi país} (calque français) → \textbf{\esp{Estoy orgulloso de mi país}}. \esp{Ser orgulloso} = « être orgueilleux » (défaut). Pour la date/lieu d'un événement : \esp{La fiesta es en mayo / es en el estadio} (pas \esp{estar}).
\item \textbf{Traduire « il y a » par autre chose que \esp{hay}.} \esp{*Son muchas lenguas} / \esp{*Hay muchas lenguas} (correct) → \esp{hay} est invariable, impersonnel, présent de \esp{haber}.
\item \textbf{Accorder « les gens » au pluriel.} \esp{La gente} = singulier en espagnol : \esp{La gente está contenta / baila / canta}. Jamais \esp{*la gente están contentos}.
\item \textbf{Calques syntaxiques du français.} « Avoir lieu » ≠ \esp{*tener un lugar} mais \esp{tener lugar} / \esp{celebrarse} / \esp{tener lugar}. « C'est pourquoi » = \esp{por eso} (pas \esp{*es por qué}). Faux ami : \esp{actualmente} = « actuellement » (et non « en fait » → \esp{en realidad}).
\end{enumerate}
\end{attention}

\begin{exercice}[Entraînement : Compréhension, Lexique, Grammaire, Expression]
\textbf{A. Vocabulaire :} Associez chaque mot à sa définition.
\begin{enumerate}
\item \esp{Patria} \quad \item \esp{Mestizaje} \quad \item \esp{Patrimonio} \quad \item \esp{Himno} \quad \item \esp{Diversidad}
\end{enumerate}
a) Canción oficial que simboliza a un país. \quad b) Mezcla de razas o culturas. \quad c) Conjunto de bienes culturales heredados. \quad d) País natal, tierra de los antepasados. \quad e) Variedad, multiplicidad de elementos diferentes.

\textbf{B. Grammaire — \esp{Ser / Estar} :} Complétez avec la forme correcte au présent.
\begin{enumerate}
\item \esp{Mi abuela \_\_\_\_ (ser) de etnia bassa.}
\item \esp{La bandera \_\_\_\_ (estar) en la plaza central hoy.}
\item \esp{Nosotros \_\_\_\_ (estar) muy contentos por la fiesta.}
\item \esp{El Ngondo \_\_\_\_ (ser) una tradición antigua.}
\item \esp{¿Dónde \_\_\_\_ (estar) la estatua de la libertad?}
\end{enumerate}

\textbf{C. Thème (Traduction FR $\rightarrow$ ES) :}
\begin{enumerate}
\item Notre hymne national est beau.
\item Il y a beaucoup de cultures au Cameroun.
\item Nous devons valoriser nos traditions.
\item Je suis fier de ma diversité culturelle.
\item La fête a lieu en décembre.
\end{enumerate}

\textbf{D. Expression écrite (10 lignes) :} \esp{Escribe un párrafo presentando la fiesta del "Ngondo" o otra fiesta de tu región a un corresponsal mexicano. Incluye: nombre, fecha, lugar, actividades, símbolos y tu opinión personal. Usa \textbf{ser}, \textbf{estar}, \textbf{presente} y el \textbf{léxico de la lección}.}
\end{exercice}

\begin{corrige}[Exercice d'entraînement]
\textbf{A. Vocabulaire :} 1-d, 2-b, 3-c, 4-a, 5-e.

\textbf{B. Grammaire :}
\begin{enumerate}
\item \esp{es} (origine/identité permanente).
\item \esp{está} (localisation physique temporaire de l'objet « bandera » aujourd'hui).
\item \esp{estamos} (sentiment/état temporaire).
\item \esp{es} (caractéristique permanente de la fête / définition).
\item \esp{está} (localisation physique d'un objet/statue).
\end{enumerate}

\textbf{C. Thème :}
\begin{enumerate}
\item \esp{Nuestro himno nacional es bonito.} (\esp{ser} + adjectif caractéristique).
\item \esp{En Camerún hay muchas culturas.} (\esp{hay} invariable).
\item \esp{Debemos valorar nuestras tradiciones.} (\esp{deber + infinitif}, accord \esp{nuestras} f. pl.).
\item \esp{Estoy orgulloso de mi diversidad cultural.} (\esp{estar orgulloso de} = être fier de ; \esp{mi} accord).
\item \esp{La fiesta tiene lugar en diciembre.} / \esp{La fiesta se celebra en diciembre.} (Locution / Passif réflexe ; mois sans article).
\end{enumerate}

\textbf{D. Expression écrite (Proposition de correction) :}
\esp{Querido amigo: Te cuento sobre el Ngondo, la gran fiesta de mi pueblo, los sawas. Se celebra cada año en diciembre, a orillas del río Wouri en Douala. Es una tradición muy antigua y sagrada. Durante la fiesta, la gente baila el "essewe", canta en duala y los jefes hacen ofrendas al agua. La piragua es el símbolo principal: representa la unión y el viaje. Estoy muy orgulloso de esta fiesta porque une a todas las generaciones. Para mí, el Ngondo es el alma de mi cultura. ¡Espero que un día puedas venir a verlo!}
\end{corrige}

\newpage

\coursec{Exercices}

\courssub{Je vérifie mes connaissances}

\begin{exercice}[QCM : le patriotisme et la diversité]
Pour chaque question, entoure la bonne réponse.
\begin{enumerate}
\item \esp{La \cle{patria}} désigne :
\begin{enumerate}
\item la famille \item le pays d'origine \item la capitale
\end{enumerate}
\item \esp{El \cle{himno} nacional} est :
\begin{enumerate}
\item un drapeau \item un chant officiel \item une danse traditionnelle
\end{enumerate}
\item \esp{El \cle{mestizaje}} signifie :
\begin{enumerate}
\item le mélange des cultures \item la guerre \item l'émigration
\end{enumerate}
\item Pour dire \esp{Yaoundé --- au centre du pays}, on utilise :
\begin{enumerate}
\item \esp{ser} \item \esp{estar} \item \esp{haber}
\end{enumerate}
\item \esp{El \cle{patrimonio} cultural} comprend :
\begin{enumerate}
\item les traditions, les langues et les fêtes \item seulement les monuments \item l'argent du pays
\end{enumerate}
\end{enumerate}
\end{exercice}

\begin{exercice}[Vrai ou faux ?]
Indique si chaque phrase est \textbf{vraie} ou \textbf{fausse} et justifie ta réponse par une courte phrase en espagnol.
\begin{enumerate}
\item \esp{Camerún tiene una sola lengua oficial.}
\item \esp{El 20 de mayo es una fiesta nacional en Camerún.}
\item \esp{Guinea Ecuatorial es un país hispanohablante.}
\item \esp{La diversidad cultural es un problema para un país.}
\item \esp{En México, la fiesta de la independencia se celebra en septiembre.}
\end{enumerate}
\end{exercice}

\begin{exercice}[Vocabulaire : trouve le mot]
Complète chaque définition avec le mot du lexique qui convient : \esp{bandera, tradición, lengua, fiesta, patrimonio, diversidad}.
\begin{enumerate}
\item Símbolo de tela con los colores de un país : \esp{la} \ldots
\item Costumbre transmitida de generación en generación : \esp{la} \ldots
\item Conjunto de bienes culturales de un pueblo : \esp{el} \ldots
\item Sistema de comunicación de un pueblo : \esp{la} \ldots
\item Día de celebración : \esp{la} \ldots
\item Variedad de culturas y de lenguas : \esp{la} \ldots
\end{enumerate}
\end{exercice}

\begin{exercice}[Conjugaison : le présent des traditions]
Conjugue les verbes entre parenthèses au présent de l'indicatif.
\begin{enumerate}
\item \esp{Cada año, los alumnos (desfilar) el 20 de mayo.}
\item \esp{Nosotros (celebrar) las fiestas tradicionales con alegría.}
\item \esp{En mi región, la gente (bailar) danzas antiguas.}
\item \esp{Yo (sentir) orgullo de mi cultura.}
\item \esp{Los griots (contar) la historia del pueblo.}
\item \esp{¿Tú (conocer) las tradiciones de tu aldea?}
\end{enumerate}
\end{exercice}

\courssub{Je m'entraîne}

\begin{exercice}[Ser ou estar ?]
Complète les dix phrases avec la forme correcte de \esp{ser} ou \esp{estar} au présent, puis justifie ton choix (identité / description permanente, lieu, état passager).
\begin{enumerate}
\item \esp{Camerún \ldots\ un país bilingüe.}
\item \esp{Yaoundé \ldots\ en el centro del país.}
\item \esp{Nosotros \ldots\ orgullosos de nuestra cultura.}
\item \esp{La bandera camerunesa \ldots\ verde, roja y amarilla.}
\item \esp{Hoy los alumnos \ldots\ muy contentos porque hay desfile.}
\item \esp{El himno nacional \ldots\ un símbolo de la patria.}
\item \esp{Mi abuela \ldots\ en el pueblo para la fiesta tradicional.}
\item \esp{Las lenguas nacionales \ldots\ una riqueza.}
\item \esp{El mercado de Douala \ldots\ cerca del puerto.}
\item \esp{¡Qué bonita \ldots\ la fiesta este año!}
\end{enumerate}
\end{exercice}

\begin{exercice}[Texte à trous : le 20 mai]
Complète le texte avec les mots suivants : \esp{patria, bandera, himno, desfile, orgullo}.
\begin{textebox}[El 20 de mayo en Camerún]
El 20 de mayo es la fiesta nacional de Camerún. Ese día, los cameruneses celebran su \ldots\ con mucha alegría. Por la mañana, los alumnos y los soldados participan en un gran \ldots\ por las calles de Yaoundé. Todos cantan el \ldots\ nacional y saludan la \ldots\ verde, roja y amarilla. En los ojos de las familias se ve el \ldots\ de ser cameruneses.
\end{textebox}
\end{exercice}

\begin{exercice}[Thème : les traditions camerounaises]
Traduis en espagnol les phrases suivantes.
\begin{enumerate}
\item Le Cameroun est un pays bilingue.
\item Chaque année, les élèves défilent le 20 mai.
\item Yaoundé se trouve au centre du pays.
\item Nous sommes fiers de notre culture.
\item La diversité est une richesse.
\end{enumerate}
\end{exercice}

\begin{exercice}[Appariements : pays et traditions]
Relie chaque pays à la tradition qui lui correspond, puis écris une phrase complète en espagnol pour chaque paire (modèle : \esp{En México, la gente celebra el Día de Muertos.}).

\begin{center}
\begin{tabularx}{\linewidth}{|X|X|}
\hline
\rowcolor{popblueL} \textbf{País} & \textbf{Tradición} \\
\hline
\esp{México} & \esp{el Día de Muertos} \\
\hline
\esp{España} & \esp{las fiestas de los pueblos con música y bailes} \\
\hline
\esp{Guinea Ecuatorial} & \esp{las danzas tradicionales fang y bubi} \\
\hline
\esp{Camerún} & \esp{el desfile del 20 de mayo} \\
\hline
\end{tabularx}
\end{center}
\end{exercice}

\courssub{Je me prépare au Bac}

\begin{exercice}[Compréhension de texte type Bac]
Lis le texte suivant, puis réponds aux questions \textbf{en espagnol}, avec des phrases complètes.

\begin{textebox}[El Grito de Independencia en México]
Cada año, la noche del 15 de septiembre, las plazas de México se llenan de gente. En la Ciudad de México, miles de personas se reúnen frente al Palacio Nacional para escuchar el famoso « Grito de Independencia ». El presidente aparece en el balcón, toca la campana y grita : « ¡Viva México! ». La multitud responde con entusiasmo y agita banderas verdes, blancas y rojas.

Esta ceremonia recuerda la noche del 15 de septiembre de 1810, cuando el sacerdote Miguel Hidalgo llamó al pueblo a luchar por la libertad. Desde entonces, el Grito es el símbolo del amor de los mexicanos por su patria.

Después de la ceremonia, las familias comen platos típicos como los tamales y el pozole, y los fuegos artificiales iluminan el cielo. Los niños llevan ropa con los colores de la bandera y cantan canciones tradicionales.

Para los mexicanos, esta fiesta no es solo una celebración : es una manera de transmitir a los jóvenes la historia y los valores de su país. Como dice un proverbio popular : « El que no conoce su historia, no sabe quién es ».
\end{textebox}

\textbf{Questions de compréhension}
\begin{enumerate}
\item (Globale) ¿Qué fiesta describe el texto y cuándo se celebra?
\item (Globale) ¿Cuál es la idea principal del texto?
\item (Détail) ¿Qué hace el presidente durante la ceremonia?
\item (Détail) ¿Qué comen las familias después de la ceremonia?
\item (Interprétation) Según el texto, ¿por qué esta fiesta es importante para los jóvenes?
\end{enumerate}

\textbf{Vrai ou faux} — Justifie chaque réponse par une citation du texte.
\begin{enumerate}
\item \esp{El Grito de Independencia se celebra por la mañana.}
\item \esp{Miguel Hidalgo era sacerdote.}
\item \esp{Los fuegos artificiales iluminan el cielo después de la ceremonia.}
\end{enumerate}

\textbf{Questions de langue}
\begin{enumerate}
\item Dans la phrase \esp{« las plazas de México se llenan de gente »}, justifie l'emploi de \esp{se}.
\item Releve dans le texte deux verbes au présent qui expriment une tradition et explique cet emploi du présent.
\end{enumerate}
\end{exercice}

\begin{exercice}[Thème et version type Bac]
\textbf{A. Thème.} Traduis en espagnol les phrases suivantes.
\begin{enumerate}
\item Notre patrie est riche de ses nombreuses cultures.
\item Le drapeau et l'hymne sont les symboles de la nation.
\item Chaque région du Cameroun a ses danses et ses langues.
\item Les jeunes doivent connaître le patrimoine de leur pays.
\item La Guinée équatoriale est un pays où l'on parle espagnol.
\end{enumerate}

\textbf{B. Version.} Traduis en français le texte suivant.
\begin{textebox}[Un país de muchas culturas]
Camerún es un país de África Central con más de doscientas lenguas. El francés y el inglés son las lenguas oficiales, pero cada región tiene sus propias tradiciones. En las fiestas, la gente baila, canta y come platos típicos. Para los cameruneses, esta diversidad no es un problema : es una gran riqueza y un motivo de orgullo.
\end{textebox}
\end{exercice}

\begin{exercice}[Expression écrite et orale type Bac]
\textbf{A. Commentaire de texte.} Lis le texte suivant, puis rédige en espagnol un commentaire de \textbf{12 à 15 lignes} avec : une introduction (présentation du texte et du thème), deux idées du texte citées et expliquées, et une conclusion personnelle.

\begin{textebox}[El patrimonio de Guinea Ecuatorial]
Guinea Ecuatorial es el único país hispanohablante de África. El español convive con las lenguas nacionales como el fang y el bubi. Esta mezcla de culturas africanas e hispánicas forma un patrimonio único. En las fiestas, las danzas tradicionales se mezclan con la música moderna. Para los jóvenes ecuatoguineanos, hablar español es una puerta abierta al mundo, pero conservar sus lenguas y sus tradiciones es conservar su identidad.
\end{textebox}

\textbf{B. Expression orale.} Présente en \textbf{2 minutes}, en espagnol, une fête traditionnelle de ta région : son nom, sa date, les activités (danses, chants, plats), et son importance pour ta communauté. Prépare-toi ensuite à répondre à deux questions du jury (par exemple : \esp{¿Cuál es tu plato típico favorito?} ou \esp{¿Por qué es importante conservar las tradiciones?}).
\end{exercice}

\newpage

\coursec{Corrigés des exercices}

\begin{corrige}[Exercice 1 — QCM : le patriotisme et la diversidad]
\textbf{Réponses :}
\begin{enumerate}
\item \textbf{b} — \esp{La patria} désigne le pays d'origine, la terre natale. \esp{Mi patria es Camerún.} « Ma patrie est le Cameroun. »
\item \textbf{b} — \esp{El himno nacional} est un chant officiel. Attention au genre : on dit \esp{el himno} (masculin), comme \esp{el himno nacional de Camerún}.
\item \textbf{a} — \esp{El mestizaje} signifie le mélange des cultures (et des peuples). \esp{El mestizaje enriquece la cultura mexicana.} « Le métissage enrichit la culture mexicaine. »
\item \textbf{b} — Pour situer un lieu, on utilise \esp{estar} : \esp{Yaoundé está en el centro del país.} « Yaoundé se trouve au centre du pays. »
\item \textbf{a} — \esp{El patrimonio cultural} comprend les traditions, les langues et les fêtes (et aussi les monuments, mais pas seulement).
\end{enumerate}
\textbf{Point de vigilance :} ne pas confondre \esp{patria} (pays d'origine) et \esp{padre} (père) ; \esp{patria} est féminin : \esp{la patria}. De même, ne pas confondre \esp{diversidad} (diversité) et \esp{diversión} (divertissement).
\end{corrige}

\begin{corrige}[Exercice 2 — Vrai ou faux ?]
\begin{enumerate}
\item \textbf{Faux.} \esp{Camerún tiene dos lenguas oficiales: el francés y el inglés.} « Le Cameroun a deux langues officielles : le français et l'anglais. »
\item \textbf{Vrai.} \esp{El 20 de mayo es la fiesta nacional de Camerún; hay un gran desfile.} « Le 20 mai est la fête nationale du Cameroun ; il y a un grand défilé. »
\item \textbf{Vrai.} \esp{Guinea Ecuatorial es el único país hispanohablante de África.} « La Guinée équatoriale est le seul pays hispanophone d'Afrique. »
\item \textbf{Faux.} \esp{La diversidad cultural es una riqueza, no un problema.} « La diversité culturelle est une richesse, pas un problème. »
\item \textbf{Vrai.} \esp{En México, la independencia se celebra el 15 y el 16 de septiembre.} « Au Mexique, l'indépendance se fête les 15 et 16 septembre. »
\end{enumerate}
\textbf{Point de vigilance :} en espagnol, les mois ne prennent pas de majuscule : \esp{mayo}, \esp{septiembre}. De même, les langues et les nationalités s'écrivent sans majuscule : \esp{el francés}, \esp{el inglés}, \esp{español}.
\end{corrige}

\begin{corrige}[Exercice 3 — Vocabulaire : trouve le mot]
\begin{enumerate}
\item \esp{la \textbf{bandera}} — le drapeau (symbole en tissu aux couleurs d'un pays).
\item \esp{la \textbf{tradición}} — la tradition (coutume transmise de génération en génération).
\item \esp{el \textbf{patrimonio}} — le patrimoine (ensemble des biens culturels d'un peuple).
\item \esp{la \textbf{lengua}} — la langue (système de communication d'un peuple).
\item \esp{la \textbf{fiesta}} — la fête (jour de célébration).
\item \esp{la \textbf{diversidad}} — la diversité (variété de cultures et de langues).
\end{enumerate}
\textbf{Point de vigilance :} attention aux genres : \esp{el patrimonio} est masculin ; \esp{la bandera}, \esp{la lengua}, \esp{la fiesta}, \esp{la tradición}, \esp{la diversidad} sont féminins. Ne pas traduire « langue » par \esp{idioma} dans ce contexte : les deux mots existent, mais \esp{lengua} est le mot du lexique de la leçon. Attention aussi à l'accent de \esp{tradición} (accent sur la dernière syllabe).
\end{corrige}

\begin{corrige}[Exercice 4 — Conjugaison : le présent des traditions]
\begin{enumerate}
\item \esp{Cada año, los alumnos \textbf{desfilan} el 20 de mayo.} « Chaque année, les élèves défilent le 20 mai. » (verbe régulier en \esp{-ar}, 3\textsuperscript{e} personne du pluriel)
\item \esp{Nosotros \textbf{celebramos} las fiestas tradicionales con alegría.} « Nous célébrons les fêtes traditionnelles avec joie. » (régulier, 1\textsuperscript{re} personne du pluriel)
\item \esp{En mi región, la gente \textbf{baila} danzas antiguas.} « Dans ma région, les gens dansent des danses anciennes. » (\esp{la gente} est un nom singulier : 3\textsuperscript{e} personne du singulier)
\item \esp{Yo \textbf{siento} orgullo de mi cultura.} « Je ressens de la fierté pour ma culture. » (\esp{sentir} : verbe à changement de voyelle \esp{e $\to$ ie} : \esp{siento, sientes, siente, sentimos, sentís, sienten})
\item \esp{Los griots \textbf{cuentan} la historia del pueblo.} « Les griots racontent l'histoire du peuple. » (\esp{contar} : changement \esp{o $\to$ ue} : \esp{cuento, cuentas, cuenta, contamos, contáis, cuentan})
\item \esp{¿Tú \textbf{conoces} las tradiciones de tu aldea?} « Connais-tu les traditions de ton village ? » (\esp{conocer} : irrégulier seulement à la 1\textsuperscript{re} personne \esp{conozco} ; ici 2\textsuperscript{e} personne régulière \esp{conoces})
\end{enumerate}
\textbf{Point de vigilance :} le présent s'emploie ici pour exprimer une habitude ou une tradition (valeur d'habitude), souvent signalée par \esp{cada año}, \esp{siempre}, \esp{normalmente}. Ne pas oublier le point d'interrogation ouvrant : \esp{¿}. Rappel : le changement de voyelle ne se produit jamais à \esp{nosotros} ni à \esp{vosotros} : \esp{sentimos}, \esp{contamos}.
\end{corrige}

\begin{corrige}[Exercice 5 — Ser ou estar ?]
\begin{enumerate}
\item \esp{Camerún \textbf{es} un país bilingüe.} — \esp{ser} : identité, caractéristique permanente.
\item \esp{Yaoundé \textbf{está} en el centro del país.} — \esp{estar} : localisation géographique.
\item \esp{Nosotros \textbf{estamos} orgullosos de nuestra cultura.} — \esp{estar} : sentiment, état ressenti (\esp{estar orgulloso de} = être fier de).
\item \esp{La bandera camerunesa \textbf{es} verde, roja y amarilla.} — \esp{ser} : description permanente (couleurs du drapeau).
\item \esp{Hoy los alumnos \textbf{están} muy contentos porque hay desfile.} — \esp{estar} : état passager (aujourd'hui).
\item \esp{El himno nacional \textbf{es} un símbolo de la patria.} — \esp{ser} : identité, définition.
\item \esp{Mi abuela \textbf{está} en el pueblo para la fiesta tradicional.} — \esp{estar} : lieu où se trouve une personne.
\item \esp{Las lenguas nacionales \textbf{son} una riqueza.} — \esp{ser} : définition, jugement général.
\item \esp{El mercado de Douala \textbf{está} cerca del puerto.} — \esp{estar} : localisation.
\item \esp{¡Qué bonita \textbf{está} la fiesta este año!} — \esp{estar} : appréciation sur une situation ponctuelle (cette année).
\end{enumerate}
\textbf{Point de vigilance :} règle résumée : \esp{ser} = identité, origine, caractéristique permanente ; \esp{estar} = lieu, état passager. Attention au piège : \esp{estar orgulloso} (être fier, sentiment du moment) et non \esp{ser orgulloso} (être orgueilleux, trait de caractère). Ne pas oublier l'accent de \esp{está}, \esp{están}.
\end{corrige}

\begin{corrige}[Exercice 6 — Texte à trous : le 20 mai]
\begin{textebox}[El 20 de mayo en Camerún — corrigé]
El 20 de mayo es la fiesta nacional de Camerún. Ese día, los cameruneses celebran su \textbf{patria} con mucha alegría. Por la mañana, los alumnos y los soldados participan en un gran \textbf{desfile} por las calles de Yaoundé. Todos cantan el \textbf{himno} nacional y saludan la \textbf{bandera} verde, roja y amarilla. En los ojos de las familias se ve el \textbf{orgullo} de ser cameruneses.
\end{textebox}
\textbf{Traduction :} « Le 20 mai est la fête nationale du Cameroun. Ce jour-là, les Camerounais célèbrent leur patrie avec beaucoup de joie. Le matin, les élèves et les soldats participent à un grand défilé dans les rues de Yaoundé. Tous chantent l'hymne national et saluent le drapeau vert, rouge et jaune. Dans les yeux des familles, on voit la fierté d'être Camerounais. »
\textbf{Justification :} \esp{su patria} (féminin singulier) ; \esp{un gran desfile} (masculin, avec l'article indéfini \esp{un} ; \esp{gran} est la forme apocopée de \esp{grande} devant un nom masculin singulier) ; \esp{el himno nacional} (masculin) ; \esp{la bandera} (féminin, article défini) ; \esp{el orgullo} (masculin, nom abstrait).
\end{corrige}

\begin{corrige}[Exercice 7 — Thème : les traditions camerounaises]
\begin{enumerate}
\item \esp{Camerún es un país bilingüe.} — \esp{ser} pour l'identité ; \esp{bilingüe} avec tréma sur le \esp{ü} (le \esp{u} se prononce).
\item \esp{Cada año, los alumnos desfilan el 20 de mayo.} — présent d'habitude ; article défini devant la date : \esp{el 20 de mayo}.
\item \esp{Yaoundé está en el centro del país.} — \esp{estar} pour la localisation ; « se trouver » ne se traduit pas mot à mot.
\item \esp{Estamos orgullosos de nuestra cultura.} — \esp{estar orgulloso de} ; accord : \esp{orgullosos} (masculin pluriel).
\item \esp{La diversidad es una riqueza.} — \esp{ser} pour une définition ; \esp{riqueza} est féminin : \esp{una riqueza}.
\end{enumerate}
\textbf{Point de vigilance :} ne pas calquer le français « est riche de » ou « se trouve » ; en espagnol, on dit simplement \esp{está en}. Penser aux accents : \esp{bilingüe} (tréma), \esp{está} (accent sur la dernière syllabe). « Chaque année » se traduit \esp{cada año} : \esp{cada} est invariable et suivi d'un nom au singulier.
\end{corrige}

\begin{corrige}[Exercice 8 — Appariements : pays et traditions]
\textbf{Appariements :} México — el Día de Muertos ; España — las fiestas de los pueblos con música y bailes ; Guinea Ecuatorial — las danzas tradicionales fang y bubi ; Camerún — el desfile del 20 de mayo.

\textbf{Phrases complètes :}
\begin{enumerate}
\item \esp{En México, la gente celebra el Día de Muertos.} « Au Mexique, les gens célèbrent le Jour des Morts. »
\item \esp{En España, la gente celebra las fiestas de los pueblos con música y bailes.} « En Espagne, les gens célèbrent les fêtes des villages avec de la musique et des danses. »
\item \esp{En Guinea Ecuatorial, la gente baila las danzas tradicionales fang y bubi.} « En Guinée équatoriale, les gens dansent les danses traditionnelles fang et bubi. »
\item \esp{En Camerún, los alumnos participan en el desfile del 20 de mayo.} « Au Cameroun, les élèves participent au défilé du 20 mai. »
\end{enumerate}
\textbf{Point de vigilance :} \esp{la gente} est singulier, le verbe se conjugue à la 3\textsuperscript{e} personne du singulier : \esp{la gente celebra}, \esp{la gente baila}. Devant un nom de pays : \esp{en México}, \esp{en España}, sans article. « Participer à » se traduit \esp{participar en} (et non \esp{participar a}).
\end{corrige}

\begin{corrige}[Exercice 9 — Compréhension de texte type Bac]
\textbf{Barème indicatif (sur 20) :} questions de compréhension : 8 pts (2 + 2 + 1 + 1 + 2) ; vrai ou faux : 6 pts (2 par item : 1 pour la réponse, 1 pour la citation) ; questions de langue : 6 pts (3 + 3).

\textbf{Questions de compréhension — réponses modèles :}
\begin{enumerate}
\item \esp{El texto describe la fiesta del Grito de Independencia de México, que se celebra la noche del 15 de septiembre.} « Le texte décrit la fête du Cri de l'Indépendance du Mexique, qui se célèbre la nuit du 15 septembre. »
\item \esp{La idea principal del texto es que esta fiesta es un símbolo del amor de los mexicanos por su patria y una manera de transmitir la historia a los jóvenes.} « L'idée principale est que cette fête est un symbole de l'amour des Mexicains pour leur patrie et une manière de transmettre l'histoire aux jeunes. »
\item \esp{El presidente aparece en el balcón, toca la campana y grita: "¡Viva México!".} « Le président apparaît au balcon, sonne la cloche et crie : "Vive le Mexique !" »
\item \esp{Después de la ceremonia, las familias comen platos típicos como los tamales y el pozole.} « Après la cérémonie, les familles mangent des plats typiques comme les tamales et le pozole. »
\item \esp{Esta fiesta es importante para los jóvenes porque les transmite la historia y los valores de su país.} « Cette fête est importante pour les jeunes parce qu'elle leur transmet l'histoire et les valeurs de leur pays. »
\end{enumerate}

\textbf{Vrai ou faux :}
\begin{enumerate}
\item \textbf{Faux.} Citation : \esp{"la noche del 15 de septiembre, las plazas de México se llenan de gente"}. La fête se célèbre la nuit, pas le matin.
\item \textbf{Vrai.} Citation : \esp{"el sacerdote Miguel Hidalgo llamó al pueblo a luchar por la libertad"}.
\item \textbf{Vrai.} Citation : \esp{"los fuegos artificiales iluminan el cielo"}, après la cérémonie.
\end{enumerate}

\textbf{Questions de langue :}
\begin{enumerate}
\item Dans \esp{"las plazas de México se llenan de gente"}, \esp{se} est la marque du verbe pronominal \esp{llenarse} (se remplir) : le sujet \esp{las plazas} subit et réalise l'action. Construction : \esp{llenarse de} + nom = « se remplir de ».
\item Exemples de verbes au présent : \esp{se llenan}, \esp{responde}, \esp{agita}, \esp{comen}, \esp{iluminan}, \esp{llevan}, \esp{cantan}. Le présent exprime ici une tradition, une action qui se répète chaque année (présent d'habitude), signalée par \esp{cada año}.
\end{enumerate}
\textbf{Points de vigilance :} répondre par des phrases complètes en reprenant les mots de la question ; pour le vrai ou faux, la citation doit être exacte et entre guillemets ; ne pas traduire le texte mais répondre en espagnol. Attention à l'accord dans \esp{les transmite} : \esp{les} reprend \esp{a los jóvenes} (complément indirect pluriel).
\end{corrige}

\begin{corrige}[Exercice 10 — Thème et version type Bac]
\textbf{Barème indicatif (sur 20) :} thème : 10 pts (2 par phrase) ; version : 10 pts (respect du sens : 6, qualité du français : 4).

\textbf{A. Thème — traductions modèles :}
\begin{enumerate}
\item \esp{Nuestra patria es rica en sus numerosas culturas.} — « riche de » se traduit \esp{rica en} ; \esp{ser} pour la caractéristique permanente.
\item \esp{La bandera y el himno son los símbolos de la nación.} — attention aux genres : \esp{la bandera} (féminin), \esp{el himno} (masculin).
\item \esp{Cada región de Camerún tiene sus danzas y sus lenguas.} — \esp{cada} + nom singulier ; verbe \esp{tener} : \esp{tiene} (changement \esp{e $\to$ ie}).
\item \esp{Los jóvenes deben conocer el patrimonio de su país.} — \esp{deber} + infinitif = devoir faire ; \esp{conocer} = connaître (une chose, un lieu).
\item \esp{Guinea Ecuatorial es un país donde se habla español.} — tournure impersonnelle \esp{se habla español} pour « on parle espagnol ».
\end{enumerate}

\textbf{B. Version — traduction modèle :}
« Le Cameroun est un pays d'Afrique centrale qui compte plus de deux cents langues. Le français et l'anglais sont les langues officielles, mais chaque région a ses propres traditions. Pendant les fêtes, les gens dansent, chantent et mangent des plats typiques. Pour les Camerounais, cette diversité n'est pas un problème : c'est une grande richesse et un motif de fierté. »

\textbf{Points de vigilance :} \esp{más de doscientas lenguas} = « plus de deux cents langues » (accord féminin : \esp{doscientas lenguas}) ; \esp{la gente} = « les gens » (singulier en espagnol, pluriel en français) ; \esp{un motivo de orgullo} = « un motif de fierté » (ne pas traduire \esp{orgullo} par « orgueil », qui est péjoratif en français) ; \esp{platos típicos} = « plats typiques » et non « plats typés » ; \esp{propio} = « propre » au sens de « qui appartient en propre », placé après le nom : \esp{sus propias tradiciones}.
\end{corrige}

\begin{corrige}[Exercice 11 — Expression écrite et orale type Bac]
\textbf{Barème indicatif (sur 20) :} respect de la consigne et de la longueur : 4 pts ; organisation (introduction, développement en deux idées, conclusion) : 4 pts ; correction grammaticale (ser/estar, présent, accords) : 6 pts ; richesse du lexique de la leçon : 4 pts ; citations du texte : 2 pts.

\textbf{A. Commentaire de texte — proposition de rédaction modèle :}

\esp{El texto "El patrimonio de Guinea Ecuatorial" presenta la riqueza cultural del único país hispanohablante de África. El tema principal es la mezcla de las culturas africanas e hispánicas.}

\esp{Primero, el autor explica que "el español convive con las lenguas nacionales como el fang y el bubi". Esta idea muestra que Guinea Ecuatorial es un país donde varias lenguas viven juntas. El español une al país con el mundo hispanohablante, pero las lenguas nacionales conservan las raíces africanas.}

\esp{Después, el texto dice que "conservar sus lenguas y sus tradiciones es conservar su identidad". Para los jóvenes, hablar español es una puerta abierta al mundo, pero las danzas y las tradiciones son el corazón de su cultura.}

\esp{En conclusión, este texto me parece muy interesante porque muestra que la diversidad es una riqueza. Como en Camerún, donde conviven muchas lenguas y culturas, cada pueblo debe estar orgulloso de su patrimonio.}

\textbf{Traduction du modèle :} « Le texte "Le patrimoine de la Guinée équatoriale" présente la richesse culturelle du seul pays hispanophone d'Afrique. Le thème principal est le mélange des cultures africaines et hispaniques. D'abord, l'auteur explique que "l'espagnol cohabite avec les langues nationales comme le fang et le bubi". Cette idée montre que la Guinée équatoriale est un pays où plusieurs langues vivent ensemble. L'espagnol relie le pays au monde hispanophone, mais les langues nationales conservent les racines africaines. Ensuite, le texte dit que "conserver ses langues et ses traditions, c'est conserver son identité". Pour les jeunes, parler espagnol est une porte ouverte sur le monde, mais les danses et les traditions sont le cœur de leur culture. En conclusion, ce texte me semble très intéressant parce qu'il montre que la diversité est une richesse. Comme au Cameroun, où cohabitent de nombreuses langues et cultures, chaque peuple doit être fier de son patrimoine. »

\textbf{B. Expression orale — trame de présentation modèle :}

\esp{Buenos días. Hoy voy a presentar una fiesta tradicional de mi región. Se llama \ldots\ y se celebra cada año en el mes de \ldots. Durante la fiesta, la gente baila danzas tradicionales, canta canciones antiguas y come platos típicos como \ldots. Los jóvenes llevan ropa con los colores de la comunidad. Esta fiesta es muy importante porque une a las familias y transmite nuestra historia a los niños. Para mí, conservar las tradiciones es conservar nuestra identidad. Muchas gracias.}

« Bonjour. Aujourd'hui, je vais présenter une fête traditionnelle de ma région. Elle s'appelle \ldots\ et se célèbre chaque année au mois de \ldots. Pendant la fête, les gens dansent des danses traditionnelles, chantent des chansons anciennes et mangent des plats typiques comme \ldots. Les jeunes portent des vêtements aux couleurs de la communauté. Cette fête est très importante parce qu'elle unit les familles et transmet notre histoire aux enfants. Pour moi, conserver les traditions, c'est conserver notre identité. Merci beaucoup. »

\textbf{Réponses modèles aux questions du jury :}
\begin{itemize}
\item \esp{¿Cuál es tu plato típico favorito?} — \esp{Mi plato típico favorito es el ndolé con camarones, porque es delicioso y me recuerda a mi familia.} « Mon plat typique préféré est le ndolé aux crevettes, parce qu'il est délicieux et qu'il me rappelle ma famille. »
\item \esp{¿Por qué es importante conservar las tradiciones?} — \esp{Es importante conservar las tradiciones porque son nuestra identidad y transmiten los valores de nuestros abuelos a los jóvenes.} « Il est important de conserver les traditions parce qu'elles sont notre identité et transmettent les valeurs de nos grands-parents aux jeunes. »
\end{itemize}

\textbf{Points de vigilance :} employer le lexique de la leçon (\esp{patria, patrimonio, tradición, diversidad, mestizaje, orgullo}) ; vérifier chaque \esp{ser/estar} ; citer le texte entre guillemets ; à l'oral, parler lentement, saluer le jury (\esp{Buenos días}) et conclure (\esp{Muchas gracias}) ; ne pas réciter mot à mot mais s'appuyer sur un plan. Attention : \esp{recordar a} = « rappeler, faire penser à » (\esp{me recuerda a mi familia}) ; \esp{une a las familias} : verbe \esp{unir} avec la préposition \esp{a} devant le complément de personne.
\end{corrige}

\newpage

\coursec{Fiche bilan}

\begin{popfigure}
\begin{tikzpicture}[mindmap, grow cyclic,
  every node/.style={concept, font=\small},
  root concept/.append style={concept color=popblue, font=\bfseries\small, minimum size=2.6cm},
  level 1 concept/.append style={sibling angle=51, level distance=3.4cm, minimum size=2.1cm}]
\node[root concept]{El patriotismo y la diversidad cultural}
  child[concept color=poporange]{ node[align=center]{Texto \\ «Mi patria, mis culturas»} }
  child[concept color=popgreen]{ node[align=center]{Léxico \\ patria, bandera, himno, patrimonio} }
  child[concept color=poppurple]{ node[align=center]{Ser / Estar \\ describir y situar} }
  child[concept color=poppink]{ node[align=center]{Presente \\ de las tradiciones} }
  child[concept color=popgold]{ node[align=center]{Método \\ comentario de texto} }
  child[concept color=popblue!60]{ node[align=center]{Culturas \\ Camerún, España, México, G. Ecuatorial} }
  child[concept color=popgreen!60]{ node[align=center]{Producción \\ presentar una fiesta} };
\end{tikzpicture}
\legende{Carte mentale de la leçon : le patriotisme et la diversité culturelle}
\end{popfigure}

\begin{bilan}
\textbf{1. Lexique clé (à connaître par cœur)}

\begin{center}
\begin{tabularx}{\linewidth}{|X|X|X|X|}
\hline
\rowcolor{popblueL} Español & Français & Español & Français \\
\hline
la patria & la patrie & la bandera & le drapeau \\
\hline
el himno & l'hymne & la tradición & la tradition \\
\hline
la fiesta & la fête & la lengua & la langue \\
\hline
el patrimonio & le patrimoine & la diversidad & la diversité \\
\hline
el mestizaje & le métissage & el símbolo & le symbole \\
\hline
la costumbre & la coutume & la identidad & l'identité \\
\hline
\end{tabularx}
\end{center}

\textbf{2. Grammaire : ser ou estar ?}

\begin{center}
\begin{tabularx}{\linewidth}{|X|X|X|}
\hline
\rowcolor{popgreenL} Verbe & Emploi & Ejemplo \\
\hline
\esp{ser} & identité, origine, caractère permanent, nationalité & \esp{Cameroon es un país diverso.} « Le Cameroun est un pays divers. » \\
\hline
\esp{estar} & localisation, état temporaire, sentiment du moment & \esp{La fiesta está en el estadio.} « La fête est au stade. » \\
\hline
\end{tabularx}
\end{center}

\begin{itemize}
\item \esp{ser} + \esp{de} = origine : \esp{Soy de Yaoundé.} « Je suis de Yaoundé. »
\item \esp{estar} + \esp{en} = lieu : \esp{El museo está en Douala.} « Le musée est à Douala. »
\item Présent de vérité générale pour les traditions : \esp{Los cameruneses celebran la fiesta nacional el 20 de mayo.} « Les Camerounais célèbrent la fête nationale le 20 mai. »
\item Attention à l'accent : \esp{celebran} (ils célèbrent) ≠ \esp{celebraban} (ils célébraient).
\end{itemize}

\textbf{3. Méthode express : commenter un texte}
\begin{enumerate}
\item Lire deux fois et souligner les mots-clés (\esp{patria, diversidad, tradición}).
\item Présenter le texte : thème, idée principale, structure.
\item Relever le lexique du patriotisme et les procédés (exemples, énumérations).
\item Donner son avis argumenté : \esp{En mi opinión... / Pienso que... / Estoy de acuerdo porque...}
\item Conclure en ouvrant sur la diversité culturelle.
\end{enumerate}

\textbf{4. Méthode express : présenter une fête}
\begin{enumerate}
\item Nommer la fête et le pays : \esp{En mi país celebramos...}
\item Situer (date, lieu) avec \esp{estar} et le présent.
\item Décrire les activités : \esp{la gente baila, canta, come...}
\item Dire ce qu'elle représente : \esp{Esta fiesta es un símbolo de nuestra identidad.}
\end{enumerate}
\end{bilan}

\begin{aretenir}[Checklist avant le Bac]
\begin{itemize}
\item[$\square$] Je connais le lexique : \esp{patria, bandera, himno, tradición, fiesta, lengua, patrimonio, diversidad, mestizaje}.
\item[$\square$] Je sais choisir entre \esp{ser} (identité, origine) et \esp{estar} (lieu, état).
\item[$\square$] Je conjugue \esp{ser} (soy, eres, es, somos, sois, son) et \esp{estar} (estoy, estás, está, estamos, estáis, están) sans erreur.
\item[$\square$] J'utilise le présent pour décrire une tradition ou une fête.
\item[$\square$] Je sais dire l'origine avec \esp{ser de} et le lieu avec \esp{estar en}.
\item[$\square$] Je peux citer des exemples culturels du Cameroun, de l'Espagne, du Mexique et de la Guinée équatoriale.
\item[$\square$] Je connais les étapes du commentaire de texte (présenter, analyser, donner son avis, conclure).
\item[$\square$] Je sais exprimer mon opinion : \esp{en mi opinión, pienso que, me parece que}.
\item[$\square$] J'évite les faux amis : \esp{la lengua} = la langue, pas « la langue de boeuf » ; \esp{la fiesta} = la fête.
\item[$\square$] Je n'oublie ni les accents (\esp{tradición, está}) ni les signes ¿ et ¡.
\item[$\square$] Je sais rédiger un paragraphe de 6 à 8 phrases pour présenter une fête de mon pays.
\item[$\square$] Je valorise la diversité culturelle avec des phrases positives : \esp{La diversidad es una riqueza.}
\end{itemize}
\end{aretenir}

\findecours

\end{document}
